% Les foyers économiques en expansion — Géographie Tle C — Cours signé OnBuch+
% Programme officiel MINESEC. Compiler avec : tectonic N14-foyers-economiques-en-expansion.tex
\newcommand{\DOCMATIERE}{Géographie}
\newcommand{\DOCNIVEAU}{Tle C}
\newcommand{\DOCMODULE}{Module 2 : La libéralisation des échanges (la mondialisation)}
\newcommand{\DOCLECON}{N14}
\newcommand{\DOCTITRE}{Les foyers économiques en expansion}
% OnBuch+ — préambule « cours signés » ENRICHI (compilé avec tectonic / XeLaTeX)
% Système visuel « Pop » d'OnBuch+ : fond crème (jamais blanc), bordures encre
% épaisses, ombres « dures » décalées, titres Archivo Black en capitales.
% Version MATHÉMATIQUES : graphes (pgfplots/tikz), boîtes pédagogiques
% (définition, propriété, méthode, attention, le savais-tu, exercice résolu,
% exercice, TP, bilan), page de garde et signature OnBuch+ ancrée en bas.
%
% Macros attendues dans main.tex AVANT \input{preamble} :
%   \DOCMATIERE  \DOCNIVEAU  \DOCMODULE  \DOCLECON (numéro)  \DOCTITRE
\documentclass[11pt]{article}

\usepackage{fontspec}
\usepackage[french]{babel}
\usepackage[a4paper,margin=2cm,top=3.4cm,bottom=2.8cm,headheight=1.4cm,footskip=1.3cm]{geometry}
\usepackage{amsmath,amssymb,amsfonts}
\usepackage{mathtools} % \xmapsto, \coloneqq... souvent utilisés par les modèles
\usepackage{cancel} % simplifications barrées (bilans de forces)
% \abs{z} (module, valeur absolue) et \norm{u} (norme), utilisés par les modèles
\providecommand{\abs}{}\let\abs\relax\DeclarePairedDelimiter{\abs}{\lvert}{\rvert}
\providecommand{\norm}{}\let\norm\relax\DeclarePairedDelimiter{\norm}{\lVert}{\rVert}
\usepackage[table]{xcolor} % [table] : fournit aussi \rowcolors (alternance de lignes)
\usepackage{pagecolor}
\usepackage{tikz}
\usetikzlibrary{arrows.meta,positioning,calc,shapes.geometric,shapes.misc,shapes.arrows,decorations.pathmorphing,decorations.pathreplacing,decorations.markings,patterns,shadows,mindmap,trees,fit,backgrounds,matrix,intersections,angles,quotes}
\usepackage{pgfplots}
\usepackage[version=4]{mhchem} % équations nucléaires \ce{^{235}_{92}U}
\usepackage[european,siunitx]{circuitikz} % schémas électriques
\pgfplotsset{compat=1.18}
\usepgfplotslibrary{fillbetween,groupplots} % aires entre courbes (calcul intégral)
\usepackage{enumitem}
\usepackage{fancyhdr}
% [most] charge toutes les bibliothèques tcolorbox (listings, minted, raster,
% documentation...) et gonfle inutilement la mémoire TeX sur un document long
% (~300 boîtes) : on ne charge que ce qui est réellement utilisé.
\usepackage[skins,breakable]{tcolorbox}
\usepackage{graphicx}
\graphicspath{{/home/user/t-t/geo-onbuch/images/}}
\usepackage{colortbl}
\usepackage{booktabs}
\usepackage{tabularx}
\usepackage{multirow}
\usepackage{diagbox} % case d'en-tête diagonale des tableaux à double entrée
% Colonnes extensibles centrée / gauche / droite (C, L, R), souvent utilisées par les modèles
\newcolumntype{C}{>{\centering\arraybackslash}X}
\newcolumntype{L}{>{\raggedright\arraybackslash}X}
\newcolumntype{R}{>{\raggedleft\arraybackslash}X}
\usepackage{array}
\usepackage{multicol}
\usepackage{siunitx}
% parse-numbers=false : accepte aussi \SI{2,5 \times 10^{-3}}{...}
\sisetup{locale=FR,output-decimal-marker={,},group-minimum-digits=4,parse-numbers=false}
\DeclareSIUnit\ohm{\ensuremath{\symup{\Omega}}} % U+2126 absent de la police
\DeclareSIUnit\division{div} % réglages d'oscilloscope (ms/div, V/div)
\DeclareSIUnit\tr{tr} % tours (vitesse de rotation, ex. \SI{1500}{\tr\per\minute})
\DeclareSIUnit\molar{M} % concentration molaire (ex. \SI{3}{\molar}, \SI{1}{\micro\molar})
\DeclareSIUnit\an{an} % année (le modèle confond parfois avec \year, un compteur TeX) — ex. \SI{5}{\kilo\meter\per\an}
\DeclareSIUnit\FCFA{FCFA} % franc CFA (ex. \SI{500}{\FCFA\per\kilo\gram})
\DeclareSIUnit\degreeBrix{\textdegree Brix} % teneur en sucre d'un sirop/jus (réfractométrie)
\DeclareSIUnit\month{mois} % le modèle utilise parfois \month, un compteur TeX, comme unité de durée
\DeclareSIUnit\microliter{\micro\liter} % le modèle écrit parfois \microliter en un seul token
\DeclareSIUnit\million{millions} % dénombrement (ex. \SI{15}{\million\per\mL} spermatozoïdes)
\usepackage{qrcode}
% \degree : commande fréquemment utilisée par les modèles pour un angle ou une
% température en dehors de \SI{}{} (non fournie par les paquets chargés).
\providecommand{\degree}{^{\circ}}
% \qed (fin de démonstration), utilisé par les modèles sans amsthm
\providecommand{\qed}{\hfill$\square$}
% \barre : double barre des valeurs interdites dans un tableau de variations (tkz-tab)
\providecommand{\barre}{\ensuremath{\|}}
% environnement proof (amsthm non chargé) utilisé par les modèles
\ifdefined\proof\else\newenvironment{proof}[1][Démonstration]{\par\noindent\textit{#1.}\ }{\qed\par}\fi
\usepackage{lastpage}
\usepackage{hyperref}
\hypersetup{hidelinks}

% ============================================================
% Palette « Pop » OnBuch+ — mêmes hex que la classe P (plus_ui.dart)
% ============================================================
\definecolor{poporange}{HTML}{F59321}
\definecolor{popdark}{HTML}{E07A0C}
\definecolor{popcream}{HTML}{FFF6E8}
\definecolor{popink}{HTML}{1C1714}
\definecolor{poppurple}{HTML}{7A5AE0}
\definecolor{popgreen}{HTML}{2E9E6A}
\definecolor{poppink}{HTML}{E0457B}
\definecolor{popblue}{HTML}{2F7DE1}
\definecolor{popgold}{HTML}{C98A12}
\definecolor{popmuted}{HTML}{8A7F76}
\definecolor{popline}{HTML}{EADFD2}
\definecolor{popgray}{HTML}{8A7F76}
\colorlet{popgrey}{popgray}
% Alias tolérants (le modèle nomme parfois les couleurs différemment)
\colorlet{popwhite}{white}
\colorlet{popblack}{popink}
\colorlet{popred}{poppink}
\colorlet{popyellow}{popgold}
% Teintes claires (fonds de tableaux/encadrés) — le modèle nomme parfois
% les couleurs avec un suffixe L (ex. popblueL) sans les avoir définies.
\colorlet{poporangeL}{poporange!20}
\colorlet{popdarkL}{popdark!20}
\colorlet{poppurpleL}{poppurple!20}
\colorlet{popgreenL}{popgreen!20}
\colorlet{poppinkL}{poppink!20}
\colorlet{popblueL}{popblue!20}
\colorlet{popgoldL}{popgold!20}
\colorlet{popredL}{popred!20}
\colorlet{popgrayL}{popgray!20}
% Teintes claires pour fonds de tableaux / zones
\colorlet{popblueL}{popblue!10!white}
\colorlet{poporangeL}{poporange!14!white}
\colorlet{popgreenL}{popgreen!12!white}
\colorlet{poppurpleL}{poppurple!12!white}
\colorlet{poppinkL}{poppink!12!white}
\colorlet{popgoldL}{popgold!14!white}

\pagecolor{popcream}

% ============================================================
% Typographie
% ============================================================
\defaultfontfeatures{Ligatures=TeX}
\setmainfont{PlusJakartaSans-Medium.ttf}[Path=../../../fonts/,BoldFont=PlusJakartaSans-Bold.ttf]
% Maths en sans-serif (Fira Math) avec chiffres et lettres droites de Plus
% Jakarta, pour que les formules se fondent dans le texte.
\usepackage[math-style=ISO,bold-style=ISO]{unicode-math}
\setmathfont{FiraMath-Regular.otf}[Path=../../../fonts/]
\setmathfont{PlusJakartaSans-Medium.ttf}[Path=../../../fonts/,range={up/num,up/{Latin,latin},\mathup}]
\usepackage{icomma} % virgule décimale sans espace en mode math
% Caractères mathématiques Unicode absents des polices de texte (Plus Jakarta,
% Archivo Black) : rendus par leur équivalent mathématique, en texte comme en
% formule — sinon ils s'affichent en carré vide (ex. « Arithmétique dans ℤ »).
\usepackage{newunicodechar}
\newunicodechar{ℝ}{\ensuremath{\mathbb{R}}}
\newunicodechar{ℤ}{\ensuremath{\mathbb{Z}}}
\newunicodechar{ℂ}{\ensuremath{\mathbb{C}}}
\newunicodechar{ℕ}{\ensuremath{\mathbb{N}}}
\newunicodechar{ℚ}{\ensuremath{\mathbb{Q}}}
\newunicodechar{𝔻}{\ensuremath{\mathbb{D}}}
\newunicodechar{α}{\ensuremath{\alpha}}
\newunicodechar{β}{\ensuremath{\beta}}
\newunicodechar{γ}{\ensuremath{\gamma}}
\newunicodechar{δ}{\ensuremath{\delta}}
\newunicodechar{ε}{\ensuremath{\varepsilon}}
\newunicodechar{θ}{\ensuremath{\theta}}
\newunicodechar{λ}{\ensuremath{\lambda}}
\newunicodechar{μ}{\ensuremath{\mu}}
\newunicodechar{σ}{\ensuremath{\sigma}}
\newunicodechar{τ}{\ensuremath{\tau}}
\newunicodechar{φ}{\ensuremath{\varphi}}
\newunicodechar{ψ}{\ensuremath{\psi}}
\newunicodechar{ω}{\ensuremath{\omega}}
\newunicodechar{Σ}{\ensuremath{\Sigma}}
% majuscules produites par \MakeUppercase dans les titres (α-aminés -> Α)
\newunicodechar{Α}{\ensuremath{\alpha}}
\newunicodechar{Β}{\ensuremath{\beta}}
\newunicodechar{Γ}{\ensuremath{\Gamma}}
\newunicodechar{Δ}{\ensuremath{\Delta}}
\newunicodechar{Ε}{\ensuremath{\varepsilon}}
\newunicodechar{Θ}{\ensuremath{\Theta}}
\newunicodechar{Λ}{\ensuremath{\Lambda}}
\newunicodechar{Μ}{\ensuremath{\mu}}
\newunicodechar{Τ}{\ensuremath{\tau}}
\newunicodechar{Φ}{\ensuremath{\Phi}}
\newunicodechar{Ψ}{\ensuremath{\Psi}}
\newunicodechar{Ω}{\ensuremath{\Omega}}
\newunicodechar{↦}{\ensuremath{\mapsto}}
\newunicodechar{∈}{\ensuremath{\in}}
\newunicodechar{∉}{\ensuremath{\notin}}
\newunicodechar{∧}{\ensuremath{\wedge}}
\newunicodechar{⇔}{\ensuremath{\Leftrightarrow}}
\newunicodechar{⟺}{\ensuremath{\Longleftrightarrow}}
\newunicodechar{⇒}{\ensuremath{\Rightarrow}}
\newunicodechar{⟹}{\ensuremath{\Longrightarrow}}
\newunicodechar{∘}{\ensuremath{\circ}}
\newunicodechar{∪}{\ensuremath{\cup}}
\newunicodechar{∩}{\ensuremath{\cap}}
\newunicodechar{≡}{\ensuremath{\equiv}}
\newunicodechar{⊂}{\ensuremath{\subset}}
\newunicodechar{⊥}{\ensuremath{\perp}}
\newunicodechar{∃}{\ensuremath{\exists}}
\newunicodechar{∀}{\ensuremath{\forall}}
\newunicodechar{∛}{\ensuremath{\sqrt[3]{\,}}}
\newunicodechar{∜}{\ensuremath{\sqrt[4]{\,}}}
\newunicodechar{⃗}{\ensuremath{{}^{\to}}}
\newunicodechar{ⁿ}{\ifmmode^{n}\else\textsuperscript{\ensuremath{n}}\fi}
\newunicodechar{ˣ}{\ifmmode^{x}\else\textsuperscript{\ensuremath{x}}\fi}
\newunicodechar{ᵇ}{\ifmmode^{b}\else\textsuperscript{\ensuremath{b}}\fi}
\newunicodechar{ʳ}{\ifmmode^{r}\else\textsuperscript{\ensuremath{r}}\fi}
\newunicodechar{ᵘ}{\ifmmode^{u}\else\textsuperscript{\ensuremath{u}}\fi}
\newunicodechar{ʸ}{\ifmmode^{y}\else\textsuperscript{\ensuremath{y}}\fi}
\newunicodechar{ᵃ}{\ifmmode^{a}\else\textsuperscript{\ensuremath{a}}\fi}
\newunicodechar{ᵉ}{\ifmmode^{e}\else\textsuperscript{\ensuremath{e}}\fi}
\newunicodechar{ᵏ}{\ifmmode^{k}\else\textsuperscript{\ensuremath{k}}\fi}
\newunicodechar{ᵖ}{\ifmmode^{p}\else\textsuperscript{\ensuremath{p}}\fi}
\newunicodechar{ᴺ}{\ifmmode^{N}\else\textsuperscript{\ensuremath{N}}\fi}
\newunicodechar{ᶜ}{\ifmmode^{c}\else\textsuperscript{\ensuremath{c}}\fi}
\newunicodechar{ʰ}{\ifmmode^{h}\else\textsuperscript{\ensuremath{h}}\fi}
\newunicodechar{ᵠ}{\ifmmode^{q}\else\textsuperscript{\ensuremath{q}}\fi}
\newunicodechar{ₙ}{\ifmmode_{n}\else\textsubscript{\ensuremath{n}}\fi}
\newunicodechar{ₐ}{\ifmmode_{a}\else\textsubscript{\ensuremath{a}}\fi}
\newunicodechar{ᵢ}{\ifmmode_{i}\else\textsubscript{\ensuremath{i}}\fi}
\newunicodechar{ᵣ}{\ifmmode_{r}\else\textsubscript{\ensuremath{r}}\fi}
\newunicodechar{ₖ}{\ifmmode_{k}\else\textsubscript{\ensuremath{k}}\fi}
\newunicodechar{ᵦ}{\ifmmode_{\beta}\else\textsubscript{\ensuremath{\beta}}\fi}
\newunicodechar{ₓ}{\ifmmode_{x}\else\textsubscript{\ensuremath{x}}\fi}
\newfontfamily\popdisplay{ArchivoBlack-Regular.ttf}[Path=../../../fonts/]
\newfontfamily\poplogo{SpaceGrotesk-Bold.ttf}[Path=../../../fonts/]
\newfontfamily\popsemi{PlusJakartaSans-SemiBold.ttf}[Path=../../../fonts/]
\newfontfamily\popxbold{PlusJakartaSans-ExtraBold.ttf}[Path=../../../fonts/]
\newfontfamily\popxboldls{PlusJakartaSans-ExtraBold.ttf}[Path=../../../fonts/,LetterSpace=3.0]

\setlist[enumerate]{leftmargin=1.7em,itemsep=5pt,topsep=4pt}
\setlist[itemize]{leftmargin=1.7em,itemsep=4pt,topsep=4pt,label={\color{poporange}\textbullet}}
\setlength{\parindent}{0pt}
\setlength{\parskip}{5pt}
\renewcommand{\arraystretch}{1.25}

% pgfplots : style Pop par défaut
\pgfplotsset{
  every axis/.append style={axis line style={popink,line width=1pt},tick style={popink},
    grid style={popline,line width=0.5pt},label style={font=\small},tick label style={font=\footnotesize}},
  popcurve/.style={poporange,line width=1.8pt,no markers,smooth},
}
% Même style disponible dans un \draw TikZ simple (hors pgfplots)
\tikzset{popcurve/.style={poporange,line width=1.8pt}}
\tikzset{popgrid/.style={popline,very thin}}
\colorlet{popcurve}{poporange} % le nom du style est parfois utilisé comme couleur

% ============================================================
% Pill / squiggle
% ============================================================
\newcommand{\poppill}[3]{%
  \tikz[baseline=-0.55ex]\node[draw=popink,line width=1.2pt,rounded corners=2.6pt,%
    fill=#2,inner xsep=6.5pt,inner ysep=3pt]{%
    \popxboldls\fontsize{7}{7}\selectfont\color{#3}\MakeUppercase{#1}};%
}
\newcommand{\popsquiggle}[1]{%
  \tikz[baseline=-0.4ex]{
    \draw[#1,line width=2.6pt,line cap=round]
      plot [smooth] coordinates {(0,0.09) (0.34,0.01) (0.68,0.21) (1.02,0.01) (1.36,0.10)};
  }%
}
% Mot-clé surligné (marqueur orange sous le texte)
\newcommand{\cle}[1]{{\popxbold\color{popdark}#1}}

% ============================================================
% En-tête / pied
% ============================================================
\pagestyle{fancy}
\fancyhf{}
\renewcommand{\headrulewidth}{0pt}
\renewcommand{\footrulewidth}{0pt}
\lhead{\raisebox{-0.15ex}{\poplogo\fontsize{13}{13}\selectfont\color{popink}OnBuch\color{poporange}+}}
\rhead{\poppill{\DOCMATIERE\ \textbullet\ \DOCNIVEAU\ \textbullet\ Leçon \DOCLECON}{white}{popink}}
\lfoot{\popxbold\fontsize{7.6}{7.6}\selectfont\color{popmuted}\MakeUppercase{Cours signé OnBuch+}\ \textbullet\ {\popsemi\DOCTITRE}}
\rfoot{\tikz[baseline=-0.6ex]\node[rounded corners=8.5pt,fill=popink,minimum height=17pt,minimum width=17pt,inner xsep=5pt,inner sep=0]{\popxbold\fontsize{8}{8}\selectfont\color{popcream}\thepage\,/\,\pageref*{LastPage}};}

% ============================================================
% Titres
% ============================================================
\newcommand{\chaptitre}[2]{%
  \begin{tcolorbox}[enhanced,colback=poporange,colframe=popink,boxrule=2.6pt,arc=11pt,
      drop shadow=popink,left=15pt,right=15pt,top=12pt,bottom=11pt,before skip=3pt,after skip=18pt]
    \poppill{#2}{popink}{popcream}\par\vspace{9pt}
    {\raggedright\hyphenpenalty=10000\exhyphenpenalty=10000\popdisplay\fontsize{22}{24}\selectfont\color{popcream}\MakeUppercase{#1}\par}
  \end{tcolorbox}
}
% Section numérotée : pastille orange avec le numéro + titre Archivo
\newcounter{popsec}
\newcommand{\coursec}[1]{%
  \stepcounter{popsec}\setcounter{popsub}{0}%
  \par\vspace{16pt}\noindent
  \begin{minipage}{\linewidth}
  \tikz[baseline=-0.7ex]\node[draw=popink,line width=1.6pt,fill=poporange,rounded corners=4pt,minimum size=22pt,inner sep=0]{\popdisplay\fontsize{12}{12}\selectfont\color{popcream}\thepopsec};%
  \hspace{8pt}{\popdisplay\fontsize{15}{17}\selectfont\color{popink}\MakeUppercase{#1}}\par
  \vspace{4pt}\noindent{\color{poporange}\rule{36pt}{3.6pt}}
  \end{minipage}\par\nopagebreak\vspace{8pt}
}
\newcounter{popsub}
\newcommand{\courssub}[1]{%
  \stepcounter{popsub}%
  \par\vspace{9pt}\noindent{\popxbold\fontsize{11.5}{13}\selectfont\color{popink}{\color{poporange}\thepopsec.\thepopsub}\hspace{6pt}#1}\par\nopagebreak\vspace{4pt}
}

% ============================================================
% Boîtes pédagogiques — même recette : fond blanc, bordure épaisse colorée,
% ombre dure encre, badge plein. Titre optionnel [#1].
% ============================================================
\tcbset{popbase/.style={breakable,enhanced,colback=white,boxrule=2.3pt,arc=9pt,
  shadow={3pt}{-3pt}{0pt}{popink},left=14pt,right=14pt,top=11pt,bottom=11pt,before skip=11pt,after skip=15pt,
  fontupper=\popsemi\fontsize{10.6}{14.5}\selectfont\color{popink},
  fontlower=\popsemi\fontsize{10.6}{14.5}\selectfont\color{popink}}}
% Coche dessinée (le glyphe ✓ n'existe pas dans les polices du document)
\newcommand{\popcheck}[1]{\tikz[baseline=-0.5ex]\draw[#1,line width=1.6pt,line cap=round,line join=round](-0.13,0.0)--(-0.04,-0.09)--(0.13,0.1);}
\providecommand{\coche}{\popcheck{popgreen}}
\providecommand{\resultat}[1]{\par\smallskip\noindent\fcolorbox{popgreen}{popgreenL}{\parbox{\dimexpr\linewidth-2\fboxsep-2\fboxrule\relax}{\textbf{Résultat.} #1}}\par\smallskip}
\newcommand{\popboxhead}[3]{\poppill{#1}{#2}{white}\ifx\relax#3\relax\else\hspace{6pt}{\popxbold\fontsize{10.6}{12}\selectfont\color{popink}#3}\fi\par\vspace{7pt}}

\newtcolorbox{aretenir}[1][]{popbase,colframe=popgreen,before upper={\popboxhead{À retenir}{popgreen}{#1}}}
\newtcolorbox{exemplebox}[1][]{popbase,colframe=poporange,before upper={\popboxhead{Exemple}{poporange}{#1}}}
\newtcolorbox{definition}[1][]{popbase,colframe=popblue,before upper={\popboxhead{Définition}{popblue}{#1}}}
\newtcolorbox{propriete}[1][]{popbase,colframe=poppurple,before upper={\popboxhead{Propriété}{poppurple}{#1}}}
\newtcolorbox{methode}[1][]{popbase,colframe=popgold,before upper={\popboxhead{Méthode}{popgold}{#1}}}
\newtcolorbox{attention}[1][]{popbase,colframe=poppink,before upper={\popboxhead{Attention}{poppink}{#1}}}
\newtcolorbox{experience}[1][]{popbase,colframe=popblue,colback=popblueL,before upper={\popboxhead{Expérience}{popblue}{#1}}}
% « Le savais-tu ? » : carte violette pleine (contexte, culture, Cameroun)
\newtcolorbox{savaistu}[1][]{popbase,colback=poppurple,colframe=popink,
  fontupper=\popsemi\fontsize{10.4}{14.2}\selectfont\color{white},
  before upper={\poppill{Le savais-tu\,?}{white}{poppurple}\ifx\relax#1\relax\else\hspace{6pt}{\popxbold\color{white}#1}\fi\par\vspace{7pt}}}
% Exercice résolu : énoncé en haut, \tcblower puis la solution sur fond crème
\newcounter{popexo}\newcounter{popexr}
\newtcolorbox{exoresolu}[1][]{popbase,colframe=poporange,colbacklower=popcream,
  segmentation style={popink,line width=1.2pt,dashed},
  before upper={\stepcounter{popexr}\popboxhead{Exercice résolu \thepopexr}{poporange}{#1}},
  before lower={\poppill{Solution}{popink}{popcream}\par\vspace{6pt}}}
% Exercice (non résolu en place) — la correction est dans la partie Corrigés
\newtcolorbox{exercice}[1][]{popbase,colframe=popink,
  before upper={\stepcounter{popexo}\popboxhead{Exercice \thepopexo}{popink}{#1}}}
% Corrigé
\newtcolorbox{corrige}[1][]{popbase,colframe=popgreen,colback=popgreenL,
  before upper={\popboxhead{Corrigé}{popgreen}{#1}}}
% Objectifs / prérequis / bilan
\newtcolorbox{objectifs}{popbase,colframe=popink,colback=poporangeL,
  before upper={\popboxhead{Objectifs de la leçon}{popink}{}}}
\newtcolorbox{prerequis}{popbase,colframe=popink,colback=popblueL,
  before upper={\popboxhead{Prérequis}{popblue}{}}}
\newtcolorbox{bilan}{popbase,colframe=popink,colback=white,boxrule=2.8pt,
  before upper={\popboxhead{Fiche bilan}{poporange}{L'essentiel à connaître pour l'examen}}}
% Figure encadrée avec légende
\newtcolorbox{popfigure}[1][]{enhanced,breakable=false,colback=white,colframe=popline,boxrule=1.4pt,arc=8pt,
  left=10pt,right=10pt,top=10pt,bottom=8pt,before skip=10pt,after skip=12pt,halign=center,
  lower separated=false,halign lower=center,
  fontlower=\popsemi\fontsize{9}{11.5}\selectfont\color{popmuted},
  #1}
\newcommand{\legende}[1]{\par\vspace{6pt}{\popsemi\fontsize{9}{11.5}\selectfont\color{popmuted}#1}}

% ============================================================
% Signature OnBuch+
% ============================================================
\newcommand{\poplogoinline}[1]{{\poplogo\fontsize{#1}{#1}\selectfont\color{popink}OnBuch\color{poporange}+}}
\newcommand{\signatureonbuch}{%
  \begin{tcolorbox}[enhanced,colback=popink,colframe=popink,boxrule=2.2pt,arc=10pt,
      drop shadow=poporange,left=14pt,right=14pt,top=10pt,bottom=10pt,before skip=0pt,after skip=0pt]
    \tikz[baseline=-0.7ex]\node[circle,fill=popgreen,draw=popcream,line width=1.3pt,minimum size=22pt,inner sep=0]{\popcheck{white}};%
    \hspace{9pt}\begin{minipage}[c]{0.62\linewidth}
      {\popxbold\fontsize{12}{13}\selectfont\color{popcream}Signé\ \textbullet\ {\poplogo OnBuch\color{poporange}+}}\par\vspace{2pt}
      {\raggedright\popsemi\fontsize{8}{10.5}\selectfont\color{popline}\MakeUppercase{Équipe pédagogique OnBuch+}\\\MakeUppercase{Conforme au programme officiel MINESEC}\par}
    \end{minipage}\hfill
    {\poplogo\fontsize{22}{22}\selectfont\color{popcream}OnBuch\color{poporange}+}
  \end{tcolorbox}%
}

% ============================================================
% Page de garde
% #1 titre · #2 matière · #3 niveau · #4 module · #5 n° leçon · #6 durée ·
% #7 description / objectifs (texte ou itemize)
% ============================================================
\newcommand{\pagegarde}[7]{%
  \thispagestyle{empty}
  \newgeometry{a4paper,margin=1.7cm,top=1.6cm,bottom=1.4cm}%
  \begin{tikzpicture}[remember picture,overlay]
    \fill[poporange!18] ([xshift=-3.2cm,yshift=-3.6cm]current page.north east) circle (4.6cm);
    \fill[poppurple!14] ([xshift=2.4cm,yshift=7.6cm]current page.south west) circle (3.8cm);
  \end{tikzpicture}%
  {\poplogo\fontsize{30}{30}\selectfont\color{popink}OnBuch\color{poporange}+}\hfill
  \raisebox{0.6ex}{\poppill{Cours signé \textbullet\ Premium}{popink}{popcream}}\par
  \vspace{1pt}\popsquiggle{poppurple}\par
  \vspace{8pt}
  \begin{tcolorbox}[enhanced,colback=poporange,colframe=popink,boxrule=2.8pt,arc=13pt,
      drop shadow=popink,left=18pt,right=18pt,top=12pt,bottom=13pt,after skip=10pt]
    \poppill{#2 \textbullet\ #3}{popink}{popcream}\hspace{5pt}\poppill{#4}{white}{popink}\par\vspace{8pt}
    {\popdisplay\fontsize{14}{14}\selectfont\color{popink}LEÇON #5}\par\vspace{4pt}
    {\raggedright\hyphenpenalty=10000\exhyphenpenalty=10000\popdisplay\fontsize{25}{27}\selectfont\color{popcream}\MakeUppercase{#1}\par}
    \vspace{8pt}
    {\popxbold\fontsize{9.5}{9.5}\selectfont\color{popink}\MakeUppercase{Durée conseillée : #6}}
  \end{tcolorbox}
  \begin{tcolorbox}[enhanced,colback=white,colframe=popink,boxrule=2.2pt,arc=10pt,
      drop shadow=popink,left=16pt,right=16pt,top=10pt,bottom=10pt,after skip=8pt]
    \poppill{Dans cette leçon}{poppurple}{white}\par\vspace{5pt}
    \popsemi\fontsize{9.3}{12.6}\selectfont\color{popink}#7
  \end{tcolorbox}
  \noindent\begin{minipage}[t]{0.487\linewidth}
    \begin{tcolorbox}[enhanced,colback=popgreen,colframe=popink,boxrule=2.2pt,arc=10pt,
        drop shadow=popink,left=11pt,right=11pt,top=10pt,bottom=10pt,height=3.1cm,valign=center]
      \tcbox[colback=white,colframe=popink,boxrule=1.3pt,arc=5pt,boxsep=0pt,nobeforeafter,
        left=3pt,right=3pt,top=3pt,bottom=3pt,valign=center]{\qrcode[height=1.9cm]{https://play.google.com/store/apps/details?id=com.luvvix.onbuch&hl=fr}}%
      \hspace{8pt}\begin{minipage}[c]{0.5\linewidth}
        {\popdisplay\fontsize{11}{12.5}\selectfont\color{white}\MakeUppercase{Télécharge OnBuch}}\par\vspace{4pt}
        {\raggedright\popsemi\fontsize{8.4}{11}\selectfont\color{white}Gratuit \textbullet\ Google Play\\Tuteur IA, annales, concours, résultats\par}
      \end{minipage}
    \end{tcolorbox}
  \end{minipage}\hfill
  \begin{minipage}[t]{0.487\linewidth}
    \begin{tcolorbox}[enhanced,colback=poppurple,colframe=popink,boxrule=2.2pt,arc=10pt,
        drop shadow=popink,left=11pt,right=11pt,top=10pt,bottom=10pt,height=3.1cm,valign=center]
      \tikz[baseline=-0.7ex]\node[circle,fill=white,draw=popink,line width=1.3pt,minimum size=1.6cm,inner sep=0]{\popdisplay\fontsize{24}{24}\selectfont\color{poppurple}+};%
      \hspace{8pt}\begin{minipage}[c]{0.6\linewidth}
        {\popdisplay\fontsize{11}{12.5}\selectfont\color{white}\MakeUppercase{Passe à OnBuch+}}\par\vspace{4pt}
        {\raggedright\popsemi\fontsize{8.4}{11}\selectfont\color{white}Tous les cours signés, Tuteur IA illimité, fiches PDF — onglet OnBuch+ de l'app\par}
      \end{minipage}
    \end{tcolorbox}
  \end{minipage}
  \vspace{8pt}
  \popcontenu
  \vfill
  \signatureonbuch
  \restoregeometry
  \newpage
}

% Rangée « Ce cours contient » (page de garde)
\newcommand{\poptile}[3]{% #1 couleur #2 gros texte #3 légende
  \begin{tcolorbox}[enhanced,colback=#1,colframe=popink,boxrule=2pt,arc=9pt,shadow={2.5pt}{-2.5pt}{0pt}{popink},
      left=6pt,right=6pt,top=9pt,bottom=9pt,halign=center,valign=center,height=1.75cm,width=\linewidth,nobeforeafter]
    {\popdisplay\fontsize{12.5}{13}\selectfont\color{white}\MakeUppercase{#2}}\par\vspace{3pt}
    {\centering\popsemi\fontsize{7.8}{9.5}\selectfont\color{white}#3\par}
  \end{tcolorbox}}
\newcommand{\popcontenu}{%
  \par\noindent{\popxboldls\fontsize{7.5}{7.5}\selectfont\color{popmuted}CE COURS SIGNÉ CONTIENT}\par\vspace{7pt}
  \noindent\begin{minipage}[t]{0.235\linewidth}\poptile{popblue}{Cours}{complet, illustré, pas à pas}\end{minipage}\hfill
  \begin{minipage}[t]{0.235\linewidth}\poptile{poporange}{Méthodes}{et exercices résolus}\end{minipage}\hfill
  \begin{minipage}[t]{0.235\linewidth}\poptile{poppink}{Exercices}{type examen + corrigés}\end{minipage}\hfill
  \begin{minipage}[t]{0.235\linewidth}\poptile{popgreen}{Fiche bilan}{l'essentiel à retenir}\end{minipage}\par}

% Sommaire
\newtcolorbox{sommaire}{enhanced,colback=white,colframe=popink,boxrule=2.2pt,arc=10pt,
  drop shadow=popink,left=18pt,right=18pt,top=13pt,bottom=13pt,before skip=6pt,after skip=20pt,
  before upper={\poppill{Sommaire}{poppurple}{white}\par\vspace{9pt}\popsemi\fontsize{10.6}{14.5}\selectfont\color{popink}}}

% Signature de fin de cours — poussée tout en bas de la dernière page
\newcommand{\findecours}{%
  \par\vfill\nopagebreak
  \begin{center}\popsquiggle{poporange}\end{center}\vspace{4pt}
  \signatureonbuch
}
\AtBeginDocument{\def\llbracket{\mathopen{[\mkern-3.5mu[}}\def\rrbracket{\mathclose{]\mkern-3.5mu]}}} % intervalles d'entiers (glyphe ⟦ absent de la police)
% Glyphes absents de Fira Math : redéfinis à partir de glyphes disponibles
\AtBeginDocument{%
  \def\setminus{\mathbin{\backslash}}\let\smallsetminus\setminus
  \def\vdots{\mathord{\vbox{\baselineskip3.2pt\lineskiplimit0pt\kern4pt\hbox{.}\hbox{.}\hbox{.}}}}%
  \def\ddots{\mathinner{\mkern1mu\raise6.5pt\hbox{.}\mkern2mu\raise3.6pt\hbox{.}\mkern2mu\raise0.7pt\hbox{.}\mkern1mu}}%
  \def\triangle{\mathord{\scalebox{0.9}{$\symup{\Delta}$}}}%
  \def\nabla{\mathord{\raisebox{\depth}{\scalebox{1}[-1]{$\symup{\Delta}$}}}}%
  \def\checkmark{\popcheck{popdark}}%
  \def\star{\mathbin{\ast}}\def\bigstar{\mathord{\ast}}%
  \def\diamond{\mathbin{\scalebox{0.6}{$\Diamond$}}}%
  \def\bigcup{\mathop{\mathchoice{\raisebox{-0.3em}{\scalebox{1.7}{$\cup$}}}{\scalebox{1.25}{$\cup$}}{\cup}{\cup}}\displaylimits}%
  \def\bigcap{\mathop{\mathchoice{\raisebox{-0.3em}{\scalebox{1.7}{$\cap$}}}{\scalebox{1.25}{$\cap$}}{\cap}{\cap}}\displaylimits}%
  \def\lesseqqgtr{\mathrel{\lessgtr}}\def\bigoplus{\mathop{\oplus}\displaylimits}%
}
\providecommand{\ssi}{si et seulement si} % abréviation fréquente
\AtBeginDocument{\providecommand{\wideparen}{\overparen}} % arc de cercle

%% FIGFIX-BEGIN v3 — correctifs globaux des figures (docs/onbuch-generation/tools/figfix)
\usepackage{adjustbox}\usepackage{etoolbox}
% toute figure TikZ/pgfplots plus large que la ligne est réduite à la largeur disponible
\BeforeBeginEnvironment{tikzpicture}{\begin{adjustbox}{max width=\linewidth}}
\AfterEndEnvironment{tikzpicture}{\end{adjustbox}}
% légendes pgfplots lisibles par-dessus les courbes
\pgfplotsset{every axis/.append style={legend style={fill=white,fill opacity=0.92,text opacity=1,draw=black!15,inner sep=3pt}}}
% étiquettes d'arêtes (syntaxe edge["..."]) avec fond blanc
\tikzset{every edge quotes/.append style={fill=white,inner sep=1.2pt}}
% bulles de cartes mentales : texte à la ligne dans la bulle, bulle un peu plus grande
\tikzset{every concept/.append style={text width=2.0cm,align=center,minimum size=2.3cm,inner sep=2pt}}
%% FIGFIX-END
\begin{document}

\pagegarde{Les foyers économiques en expansion}{Géographie}{Tle C}{Module 2 : La libéralisation des échanges (la mondialisation)}{N14}{2 h}{Cette leçon identifie et localise les pays émergents (BRICS, NPI, Dragons, pays du Golfe, puissances africaines) et analyse les facteurs, les manifestations et les limites de leur insertion croissante dans la mondialisation. Elle montre comment ces foyers économiques en expansion bouleversent la hiérarchie traditionnelle des puissances mondiales.
\vspace{4pt}
\begin{multicols}{2}\small
\begin{itemize}
\item À la fin de cette leçon, je sais définir les notions de pays émergent, NPI, BRICS et pays à revenus intermédiaires
\item À la fin de cette leçon, je sais localiser sur une carte du monde les principaux pays émergents d'Amérique, d'Asie, du Moyen-Orient et d'Afrique
\item À la fin de cette leçon, je sais décrire les manifestations de l'expansion économique de ces pays (industrie, commerce, IDE, firmes multinationales)
\item À la fin de cette leçon, je sais expliquer les facteurs de l'émergence (matières premières, main-d'œuvre, réformes, ouverture)
\item À la fin de cette leçon, je sais construire et légender une carte de localisation des foyers émergents
\item À la fin de cette leçon, je sais construire un diagramme (bâtons, secteurs) à partir de données statistiques sur les BRICS ou les NPI
\end{itemize}\end{multicols}}

\begin{sommaire}
\begin{multicols}{2}
\begin{enumerate}
\item Définir et localiser les foyers économiques en expansion
\item Les pays émergents d'Amérique : Brésil, Mexique, Argentine
\item Les pays émergents d'Asie et du Moyen-Orient
\item Les pays émergents d'Afrique : Afrique du Sud, Nigéria, Égypte, Maroc
\item Facteurs, manifestations et limites de l'émergence
\item TP guidé : cartographier et représenter les foyers émergents
\item Activité d'intégration
\item Méthodes et exercices résolus
\item Exercices
\item Corrigés des exercices
\item Fiche bilan
\end{enumerate}
\end{multicols}
\end{sommaire}

\begin{objectifs}
\begin{itemize}
\item À la fin de cette leçon, je sais définir les notions de pays émergent, NPI, BRICS et pays à revenus intermédiaires
\item À la fin de cette leçon, je sais localiser sur une carte du monde les principaux pays émergents d'Amérique, d'Asie, du Moyen-Orient et d'Afrique
\item À la fin de cette leçon, je sais décrire les manifestations de l'expansion économique de ces pays (industrie, commerce, IDE, firmes multinationales)
\item À la fin de cette leçon, je sais expliquer les facteurs de l'émergence (matières premières, main-d'œuvre, réformes, ouverture)
\item À la fin de cette leçon, je sais construire et légender une carte de localisation des foyers émergents
\item À la fin de cette leçon, je sais construire un diagramme (bâtons, secteurs) à partir de données statistiques sur les BRICS ou les NPI
\item À la fin de cette leçon, je sais interpréter des documents (cartes, photographies, tableaux statistiques) pour comparer les trajectoires d'émergence
\item À la fin de cette leçon, je sais distinguer pays émergents et pays au décollage économique difficile (PPTE) et nuancer le terme « émergent » à partir de ses limites sociales et spatiales
\end{itemize}
\end{objectifs}

\begin{prerequis}
\begin{itemize}
\item La notion de mondialisation et les acteurs de la libéralisation des échanges (Module 2)
\item La Triade et la hiérarchie des puissances économiques mondiales (leçons précédentes du chapitre 5)
\item La notion d'IDE (investissements directs à l'étranger) et de firme transnationale
\item Les repères de localisation des continents et des grands ensembles géographiques (Première)
\item La lecture de cartes et de tableaux statistiques (échelles, légendes, pourcentages)
\end{itemize}
\end{prerequis}

\newpage

\coursec{Définir et localiser les foyers économiques en expansion}

\textbf{Situation d'accroche.} À Douala, quartier Deïdo, un jeune entrepreneur déballe ses marchandises : des téléphones \og made in China \fg{}, des pièces détachées venues du Brésil, des dattes importées des Émirats Arabes Unis. Le soir, à la télévision, on entend que les BRICS réclament plus de poids dans les institutions mondiales, tandis que le Nigéria voisin est présenté comme le \og géant économique de l'Afrique \fg{}. Pourtant, dans ces mêmes pays, des millions de personnes vivent encore dans la pauvreté. Comment des pays longtemps classés parmi les pauvres sont-ils devenus des \cle{foyers économiques en expansion} ? Quels sont ces pays, où se situent-ils et sur quoi repose leur essor ? Le Cameroun peut-il s'en inspirer ? C'est à ces questions que répond cette première partie de la leçon : définir ce qu'est un pays émergent, puis le localiser sur la carte du monde.

\courssub{Qu'est-ce qu'un pays émergent ?}

\textbf{L'intuition d'abord.} Observe le conteneur que décharge le jeune entrepreneur de Deïdo : il y a trente ans, presque tout ce qui entrait au port de Douala venait d'Europe ou d'Amérique du Nord. Aujourd'hui, une grande partie vient de Chine, d'Inde, du Brésil ou de Dubaï. Cela signifie que de nouveaux pays sont devenus capables de \emph{produire} et d'\emph{exporter} des biens industriels et des services à l'échelle du monde. Ces pays ne sont plus des pays pauvres en difficulté, mais ils ne sont pas encore des pays riches : ils sont \og en train d'émerger \fg{}, comme un nageur qui remonte vers la surface.

\begin{definition}[Pays émergent]
Un \cle{pays émergent} est un pays en développement qui connaît une \cle{croissance économique forte et durable}, une \cle{industrialisation rapide} et une \cle{insertion croissante dans les échanges mondiaux} (exportations, investissements, finance). L'émergence est un \emph{processus inachevé} : le pays reste marqué par de fortes inégalités internes et par la pauvreté d'une partie de sa population.
\end{definition}

\textbf{Comment le géographe vérifie-t-il qu'un pays émerge ?} Il s'appuie sur des documents statistiques. Trois indices simples suffisent pour un premier diagnostic :
\begin{itemize}
\item le \cle{taux de croissance du PIB} (produit intérieur brut) : un pays émergent croît durablement plus vite que les pays développés (souvent plus de 4 à 5 \% par an, contre 1 à 3 \% dans les pays riches) ;
\item la \cle{part de l'industrie et des services} dans le PIB, qui progresse au détriment de l'agriculture ;
\item la \cle{part dans le commerce mondial} : le pays exporte de plus en plus, y compris des produits manufacturés, et attire des investissements étrangers.
\end{itemize}

\begin{exemplebox}[Le cas du Brésil]
Le Brésil est longtemps resté un exportateur de café et de sucre. Depuis les années 1990-2000, il est devenu une puissance industrielle (avions Embraer, automobiles, agro-industrie), un exportateur majeur de soja, de viande et de minerai de fer, et son PIB le place parmi les dix à douze premières économies mondiales (ordre de grandeur : environ \num{2000} milliards de dollars, à manier avec prudence car le classement varie selon les années). Pourtant, les \emph{favelas} de Rio de Janeiro ou de São Paulo rappellent que l'émergence n'efface pas la pauvreté.
\end{exemplebox}

\courssub{Des NPI aux BRICS : les catégories de l'émergence}

\textbf{Les Nouveaux Pays Industriels (NPI).} Le terme apparaît dans les années 1970-1980 pour désigner des pays qui, contrairement aux autres pays du Sud, ont réussi une industrialisation rapide tournée vers l'exportation.

\begin{definition}[NPI]
Les \cle{Nouveaux Pays Industriels} (NPI) sont des pays ayant connu, à partir des années 1960-1970, une \cle{industrialisation rapide} fondée sur les \cle{exportations} de produits manufacturés à bas coût, puis de haute technologie. On distingue :
\begin{itemize}
\item les quatre \cle{Dragons} d'Asie de l'Est : \textbf{Corée du Sud, Taïwan, Singapour, Hong Kong} (Hong Kong étant, rappelons-le, un territoire et non un État : rattaché à la Chine depuis 1997, il conserve un statut particulier) ;
\item les \og bébés tigres \fg{} (ou nouveaux Dragons) : \textbf{Malaisie, Thaïlande}, parfois Indonésie et Philippines.
\end{itemize}
\end{definition}

\begin{savaistu}[Pourquoi des « Dragons » ?]
Dans la culture asiatique, le dragon symbolise la puissance et la réussite. Les quatre Dragons ont appliqué une recette voisine : État volontariste, éducation de masse, main-d'œuvre disciplinée et bon marché au départ, zones franches attirant les firmes multinationales, puis montée en gamme technologique. La Corée du Sud, pays ravagé par la guerre en 1953, est aujourd'hui un champion mondial de l'électronique (Samsung) et de l'automobile (Hyundai) : elle est même désormais classée parmi les pays développés, preuve qu'un NPI peut achever son émergence.
\end{savaistu}

\textbf{Les BRICS : l'émergence devient puissance politique.}

\begin{definition}[BRICS]
Les \cle{BRICS} sont un groupe informel réunissant les grandes puissances émergentes : \textbf{B}résil, \textbf{R}ussie, \textbf{I}nde, \textbf{C}hine et, depuis 2010, l'Afrique du \textbf{S}ud. Ils revendiquent un poids accru dans la \cle{gouvernance mondiale} (ONU, FMI, Banque mondiale) et coopèrent entre eux (sommets annuels depuis 2009, Nouvelle Banque de développement créée en 2014-2015). Cas particulier : la Russie, héritière d'une économie déjà industrialisée (ex-URSS), est membre du groupe pour son poids démographique, énergétique et diplomatique plutôt que pour un décollage industriel récent.
\end{definition}

\begin{savaistu}[Un acronyme de banquier devenu réalité géopolitique]
Le terme \og BRIC \fg{} a été inventé en 2001 par Jim O'Neill, économiste de la banque d'affaires américaine \textbf{Goldman Sachs}, pour désigner les pays dont la croissance devait transformer l'économie mondiale d'ici 2050. Les pays concernés ont transformé ce simple sigle boursier en véritable club diplomatique : premier sommet en 2009, entrée de l'Afrique du Sud en 2010 (d'où le \og S \fg{}). Lors du sommet de Johannesburg en 2023, le groupe a annoncé un élargissement à de nouveaux membres (dont l'Égypte, l'Éthiopie, l'Iran, l'Arabie Saoudite et les Émirats Arabes Unis, invités à partir de 2024) : c'est un complément utile pour le Bac, à citer avec prudence car la composition évolue.
\end{savaistu}

\textbf{Les pays à revenus intermédiaires.} La Banque mondiale classe les pays selon le \cle{revenu national brut (RNB) par habitant} : pays à faible revenu, pays à \cle{revenu intermédiaire} (tranche inférieure et tranche supérieure), pays à revenu élevé. La plupart des pays émergents (Brésil, Mexique, Chine, Afrique du Sud, Malaisie) appartiennent à la catégorie des pays à revenu intermédiaire : ils ont dépassé le stade de la pauvreté extrême généralisée sans atteindre le niveau de vie des pays riches. À l'opposé, les pays les plus pauvres, souvent très endettés, sont appelés \cle{PPTE} (Pays Pauvres Très Endettés) : ils seront étudiés dans la leçon suivante. Le Cameroun, pour situer, se trouve dans la tranche inférieure des revenus intermédiaires.

\courssub{Localiser les foyers économiques en expansion dans le monde}

\textbf{Observation guidée.} Reportons maintenant ces pays sur un planisphère simplifié. L'exercice est formateur : au Bac, savoir \emph{placer} un pays est aussi important que savoir le définir.

\begin{popfigure}
\includegraphics[width=\linewidth,height=8cm,keepaspectratio]{N14/emergents.png}
\legende{Les marchés émergents et les pays développés dans le monde (carte de 2006). Rose : pays développés ; bleu : marchés émergents (Chine, Inde, Brésil, Mexique, Russie, Afrique du Sud, Asie du Sud-Est, Nigéria\ldots) ; gris : autres pays. On lit une répartition mondiale de l'émergence : Amérique latine, Afrique, Europe orientale et surtout Asie. \\ Source : Wikimedia Commons, AlexCovarrubias, domaine public.}
\end{popfigure}

\textbf{Lecture de la carte.} Trois grands ensembles apparaissent :
\begin{itemize}
\item en \textbf{Amérique} : le Brésil, le Mexique et l'Argentine ;
\item en \textbf{Asie et au Moyen-Orient} : la Chine et l'Inde (géants démographiques), les Dragons et la Malaisie, ainsi que les monarchies pétrolières du Golfe (Arabie Saoudite, Émirats Arabes Unis), enrichies par l'hydrocarbure et reconverties dans la finance, le transport aérien et le tourisme ;
\item en \textbf{Afrique} : l'Afrique du Sud (seul membre africain des BRICS), le Nigéria (première puissance démographique du continent et l'un de ses premiers PIB, le rang exact variant selon les années et les révisions statistiques), l'Égypte et le Maroc (industries, tourisme, agriculture exportatrice).
\end{itemize}

\begin{methode}[Légender une carte de localisation]
Pour construire la légende d'une carte des pays émergents :
\begin{enumerate}
\item Choisir \textbf{un code couleur par catégorie} (ici : BRICS, Dragons, autres émergents, pays du Golfe) et s'y tenir sur toute la carte.
\item \textbf{Ordonner la légende} de la catégorie la plus importante à la moins importante (les BRICS d'abord, car ce sont les chefs de file).
\item Utiliser des \textbf{symboles simples} (points pour les États de petite taille comme Singapour, zones colorées pour les grands pays).
\item Ne pas oublier le \textbf{titre}, l'\textbf{orientation} (flèche du nord) et, si possible, une \textbf{échelle indicative}.
\end{enumerate}
\end{methode}

\courssub{Des caractéristiques communes qui confirment l'émergence}

Au-delà de leur diversité, les pays émergents partagent des traits communs que l'on peut vérifier statistiquement :
\begin{itemize}
\item une \textbf{forte démographie} : la Chine et l'Inde dépassent chacune 1,4 milliard d'habitants (l'Inde ayant dépassé la Chine vers 2023), le Brésil environ 215 millions, le Nigéria plus de 220 millions (ordres de grandeur, années 2020) ; cette population est à la fois une main-d'œuvre abondante et un immense marché intérieur ;
\item une \textbf{croissance du PIB supérieure} à celle des pays développés, comme le montre le graphique ci-dessous ;
\item une \textbf{urbanisation rapide} : São Paulo, Mexico, Shanghai, Mumbai (anciennement Bombay), Lagos ou Le Caire comptent parmi les plus grandes métropoles du monde ;
\item un \textbf{poids croissant dans le commerce mondial} : la Chine est devenue le premier exportateur mondial de marchandises ; les BRICS représentent ensemble, en ordre de grandeur, environ un quart du PIB mondial (davantage si l'on raisonne en parité de pouvoir d'achat) et plus de 40 \% de la population de la planète.
\end{itemize}

\begin{popfigure}
\begin{center}
\begin{tikzpicture}
\begin{axis}[
  ybar, bar width=0.55cm,
  width=0.95\linewidth, height=6.2cm,
  ymin=0, ymax=9,
  ylabel={Croissance moyenne annuelle du PIB (en \%)},
  symbolic x coords={Chine,Inde,Brésil,Afr. du Sud,USA,Zone euro},
  xtick=data, x tick label style={font=\scriptsize, rotate=15, anchor=east},
  ytick={0,2,4,6,8},
  nodes near coords, nodes near coords style={font=\scriptsize, popdark},
  axis lines=left, enlarge x limits=0.12,
  every axis plot/.append style={fill=popblue, draw=popdark}
]
\addplot coordinates {(Chine,7.5) (Inde,6.2) (Brésil,2.3) (Afr. du Sud,1.8) (USA,2.0) (Zone euro,1.2)};
\end{axis}
\end{tikzpicture}
\end{center}
\legende{Taux de croissance moyen annuel du PIB comparé (ordres de grandeur indicatifs sur les deux premières décennies du XXI\textsuperscript{e} siècle, à manier avec prudence : les moyennes varient selon la période retenue ; source d'inspiration : Banque mondiale). Les géants émergents d'Asie croissent nettement plus vite que les économies développées ; le Brésil et l'Afrique du Sud montrent que l'émergence peut aussi connaître des ralentissements.}
\end{popfigure}

\textbf{Lecture du graphique.} Deux enseignements : d'une part, la Chine (environ 7,5 \% par an en moyenne) et l'Inde (environ 6 \%) croissent trois à cinq fois plus vite que les États-Unis ou la zone euro ; d'autre part, le Brésil et l'Afrique du Sud affichent des moyennes proches de celles des pays riches, ce qui rappelle que l'émergence n'est ni automatique ni irréversible : elle peut marquer le pas.

\begin{attention}[Ne pas confondre « émergent » et « développé »]
Erreur fréquente au Bac : présenter la Chine ou le Brésil comme des pays \og développés \fg{}. L'émergence est un \emph{processus inachevé}. Un pays émergent conserve : un revenu par habitant inférieur à celui des pays riches, de \textbf{fortes inégalités internes} (littoral chinois développé contre intérieur rural ; townships sud-africains ; favelas brésiliennes), une part importante d'emplois informels et une dépendance possible aux matières premières. À l'inverse, ne confonds pas non plus « émergent » et « PPTE » : les Pays Pauvres Très Endettés, eux, n'ont pas amorcé de décollage durable.
\end{attention}

\begin{exoresolu}[Identifier et classer des pays]
Énoncé : classe les pays suivants en trois catégories (BRICS ; NPI-Dragons ; pays du Golfe) et justifie en une phrase chaque classement : Singapour, Arabie Saoudite, Inde, Corée du Sud, Afrique du Sud, Émirats Arabes Unis.
\tcblower
\textbf{Solution détaillée.}
\begin{itemize}
\item \textbf{BRICS} : l'\textbf{Inde} (géant démographique, forte croissance, membre fondateur du groupe) et l'\textbf{Afrique du Sud} (admise en 2010, première puissance industrielle d'Afrique subsaharienne).
\item \textbf{NPI-Dragons} : \textbf{Singapour} et la \textbf{Corée du Sud}, industrialisés dès les années 1960-1970 grâce aux exportations manufacturées ; ils sont aujourd'hui si avancés qu'on les range souvent parmi les économies développées.
\item \textbf{Pays du Golfe} : l'\textbf{Arabie Saoudite} et les \textbf{Émirats Arabes Unis}, dont l'expansion repose d'abord sur la rente pétrolière et gazière, réinvestie dans la finance, l'aviation (Emirates, Etihad) et les infrastructures (Dubaï).
\end{itemize}
Piège évité : classer l'Arabie Saoudite parmi les BRICS parce qu'elle est riche — la richesse pétrolière ne suffit pas à définir un BRICS, qui est un club diplomatique précis (même si l'Arabie Saoudite a été invitée à le rejoindre lors de l'élargissement annoncé en 2023).
\end{exoresolu}

\begin{aretenir}[L'essentiel de la section]
Un \cle{pays émergent} est un pays du Sud dont la croissance forte et durable, l'industrialisation et l'insertion dans les échanges mondiaux le rapprochent des pays développés, sans que le processus soit achevé. Les pionniers furent les \cle{NPI} asiatiques (les quatre Dragons : Corée du Sud, Taïwan, Singapour, Hong Kong). Aujourd'hui, les \cle{BRICS} (Brésil, Russie, Inde, Chine, Afrique du Sud) incarnent l'émergence devenue puissance mondiale. Les foyers en expansion se répartissent sur trois continents : Amérique (Brésil, Mexique, Argentine), Asie et Moyen-Orient (Chine, Inde, Malaisie, Arabie Saoudite, Émirats), Afrique (Afrique du Sud, Nigéria, Égypte, Maroc). Leurs points communs : démographie massive, croissance supérieure à celle des pays riches, urbanisation rapide, poids croissant dans le commerce mondial — mais aussi inégalités internes persistantes et risque de ralentissement.
\end{aretenir}

\coursec{Les pays émergents d'Amérique : Brésil, Mexique, Argentine}

Dans la section précédente, nous avons défini les \cle{foyers économiques en expansion} et localisé les principaux \cle{pays émergents} sur le planisphère. Nous avons retenu qu'un pays émergent est un pays en développement dont la croissance rapide, l'industrialisation et l'insertion dans les échanges mondiaux lui confèrent un poids nouveau dans l'économie mondiale. Il est temps maintenant d'entrer dans le détail, continent par continent. Commençons par l'Amérique latine, où trois pays dominent : le \cle{Brésil}, le \cle{Mexique} et l'\cle{Argentine}.

Pour t'y situer d'emblée, imagine un conteneur chargé au port de Santos, près de São Paulo : il contient du soja brésilien destiné à la Chine. Un autre, parti de la frontière américano-mexicaine, transporte des pièces automobiles assemblées dans une usine de Ciudad Juárez vers les États-Unis. Un troisième, embarqué à Buenos Aires, est rempli de viande bovine argentine. Ces trois flux racontent trois visages de l'émergence latino-américaine.

\begin{popfigure}
\includegraphics[width=\linewidth,height=11cm,keepaspectratio]{N14/amlat.jpg}
\legende{Carte politique de l'Amérique latine (en anglais), où se lisent les trois foyers économiques étudiés et leurs principaux pôles urbains. 1 : le Mexique (Mexico), au sud des États-Unis dont il est voisin, atelier manufacturier adossé à ce marché ; 2 : le Brésil (São Paulo, Rio de Janeiro, Brasília), géant agricole, minier et industriel ; 3 : l'Argentine (Buenos Aires), puissance agroalimentaire. La carte ne représente pas les flux d'exportation. \\ Source : Wikimedia Commons, Central Intelligence Agency (États-Unis), domaine public.}
\end{popfigure}

\courssub{Le Brésil, géant sud-américain et membre des BRICS}

\courssub{Une puissance aux dimensions continentales}

Le Brésil est le cinquième pays du monde par sa superficie (environ \num{8,5} millions de km², soit près de 18 fois le Cameroun) et par sa population (environ \num{215} millions d'habitants, ordre de grandeur). C'est la \cle{première puissance économique d'Amérique latine} : son PIB se situe autour de \num{2000} milliards de dollars (ordre de grandeur, variable selon les années et le taux de change), ce qui le place régulièrement parmi les dix premières économies mondiales. Membre fondateur des \cle{BRICS} (Brésil, Russie, Inde, Chine, Afrique du Sud), il prétend à un rôle de puissance mondiale, symbolisé par l'accueil de la Coupe du monde de football en 2014 et des Jeux olympiques de Rio en 2016.

\courssub{Des atouts agricoles, miniers, énergétiques et industriels}

L'émergence brésilienne repose sur quatre piliers complémentaires.

\begin{itemize}
\item \textbf{Une agriculture commerciale tournée vers l'exportation.} Le Brésil est le premier exportateur mondial de \cle{soja}, de \cle{café} et de \cle{sucre} (issu de la canne à sucre), ainsi que de viande bovine et de jus d'orange concentré. De vastes exploitations mécanisées (l'\emph{agribusiness}) occupent le Centre-Ouest et le Sud, tandis que la canne à sucre alimente aussi la production d'\cle{éthanol}, biocarburant dont le Brésil est pionnier.
\item \textbf{Un sous-sol riche.} Le Brésil est l'un des tout premiers producteurs et exportateurs mondiaux de \cle{minerai de fer} (mines de Carajás en Amazonie, quadrilatère ferrifère du Minas Gerais), exploité notamment par le groupe Vale.
\item \textbf{Un pétrole offshore prometteur.} La compagnie nationale \cle{Petrobras} exploite d'importants gisements en mer, au large de Rio de Janeiro (gisements dits \emph{pré-salifères}, enfouis sous une épaisse couche de sel marin), qui ont fait du Brésil l'un des grands producteurs de pétrole d'Amérique latine.
\item \textbf{Une industrie diversifiée.} Automobile, acier, agroalimentaire, chimie, et surtout aéronautique avec \cle{Embraer}. Le pôle industriel dominant est \cle{São Paulo} : la seule région métropolitaine de São Paulo concentre environ un tiers du PIB brésilien (ordre de grandeur couramment admis).
\end{itemize}

\begin{exemplebox}[Le Brésil en chiffres]
Premier exportateur mondial de soja, de café et de sucre ; parmi les trois premiers exportateurs mondiaux de minerai de fer ; la métropole de São Paulo pèse à elle seule environ un tiers du PIB national. Ces chiffres illustrent une insertion réussie dans la mondialisation : le Brésil est devenu un fournisseur incontournable de la Chine, premier client de son soja et de son fer.
\end{exemplebox}

\begin{popfigure}
\begin{center}
\begin{tikzpicture}
% Diagramme circulaire : structure indicative des exportations brésiliennes
% Produits agricoles 40 %, minerais et métaux 25 %, manufacturés 25 %, pétrole et autres 10 %
\fill[popgold!75] (0,0) -- (90:2.4) arc (90:-54:2.4) -- cycle;
\fill[popgreen!65] (0,0) -- (-54:2.4) arc (-54:-144:2.4) -- cycle;
\fill[popblue!60] (0,0) -- (-144:2.4) arc (-144:-234:2.4) -- cycle;
\fill[poporange!65] (0,0) -- (-234:2.4) arc (-234:-270:2.4) -- cycle;
\draw[popink] (0,0) circle (2.4);
\node[font=\small\bfseries,white] at (18:1.5) {40 \%};
\node[font=\small\bfseries,white] at (-99:1.5) {25 \%};
\node[font=\small\bfseries,white] at (-189:1.5) {25 \%};
\node[font=\small\bfseries,white] at (-252:1.6) {10 \%};
% Légende
\fill[popgold!75] (3.2,1.6) rectangle (3.7,1.95); \node[anchor=west,font=\small] at (3.85,1.78) {Produits agricoles (soja, café, sucre, viandes)};
\fill[popgreen!65] (3.2,0.7) rectangle (3.7,1.05); \node[anchor=west,font=\small] at (3.85,0.88) {Minerais et métaux (fer, or)};
\fill[popblue!60] (3.2,-0.2) rectangle (3.7,0.15); \node[anchor=west,font=\small] at (3.85,-0.02) {Produits manufacturés (avions, automobiles)};
\fill[poporange!65] (3.2,-1.1) rectangle (3.7,-0.75); \node[anchor=west,font=\small] at (3.85,-0.92) {Pétrole et autres produits};
\end{tikzpicture}
\end{center}
\legende{Structure indicative des exportations brésiliennes par grandes catégories de produits (ordres de grandeur, moyennes des années récentes ; les parts varient selon les cours mondiaux). La prédominance des produits agricoles et miniers (environ les deux tiers des recettes) révèle une forte dépendance aux matières premières.}
\end{popfigure}

\begin{savaistu}[Embraer, le petit géant de l'aéronautique]
Fondée en 1969 par l'État brésilien puis privatisée, l'entreprise \cle{Embraer} (São José dos Campos) est devenue le troisième constructeur aéronautique mondial, derrière l'américain Boeing et l'européen Airbus. Spécialisée dans les avions régionaux et d'affaires, elle exporte ses appareils sur tous les continents : la preuve qu'un pays émergent peut réussir dans une industrie de pointe, et pas seulement dans les matières premières.
\end{savaistu}

\courssub{Le Mexique, atelier manufacturier adossé aux États-Unis}

Le Mexique (environ \num{128} millions d'habitants) est la deuxième économie d'Amérique latine. Son émergence suit un modèle différent de celui du Brésil : il s'est industrialisé \emph{à l'ombre} du géant américain.

\begin{definition}[Maquiladora]
Une \cle{maquiladora} est une usine de sous-traitance installée au Mexique, le plus souvent près de la frontière avec les États-Unis (Ciudad Juárez, Tijuana, Monterrey), qui assemble des produits (textiles, électronique, pièces automobiles) à partir de composants importés, en bénéficiant d'avantages fiscaux et douaniers et d'une main-d'œuvre bon marché, avant de réexporter l'essentiel de la production vers le marché américain.
\end{definition}

L'intuition est simple : une entreprise américaine paie un ouvrier mexicain beaucoup moins cher qu'un ouvrier américain ; en installant l'usine juste de l'autre côté de la frontière, elle réduit ses coûts tout en restant à quelques heures de camion de ses clients. Ce mécanisme s'est amplifié avec l'\cle{ALENA} (Accord de libre-échange nord-américain, entré en vigueur en 1994), remplacé en 2020 par l'\cle{AEUMC} (Accord Canada--États-Unis--Mexique, dit USMCA en anglais), qui a fait de l'Amérique du Nord une vaste zone de libre-échange. Résultat : le Mexique est devenu un grand exportateur de produits manufacturés (automobiles, appareils électroniques, électroménager), réalisant la majeure partie de son commerce avec les États-Unis.

À cette industrie s'ajoutent deux autres atouts : le \cle{pétrole} (le Mexique est un producteur historique, avec la compagnie nationale Pemex et les gisements du golfe du Mexique) et le \cle{tourisme} international (Cancún, Riviera Maya, Mexico), qui rapporte d'importantes devises. Le Mexique figure ainsi parmi les \cle{NPI} (nouveaux pays industrialisés) retenus par le programme.

\courssub{L'Argentine, puissance agroalimentaire à l'économie instable}

L'Argentine (environ \num{46} millions d'habitants) offre un troisième visage de l'émergence : celui d'une \cle{puissance agroalimentaire} freinée par une instabilité chronique.

Son atout majeur est la \cle{pampa}, immense plaine fertile qui s'étend autour de Buenos Aires : c'est l'une des meilleures terres agricoles du monde. L'Argentine figure parmi les premiers exportateurs mondiaux de \cle{soja} (grains, huile, tourteaux), de \cle{maïs}, de \cle{blé} et de \cle{viande bovine} — le célèbre bœuf argentin. L'agroalimentaire et l'industrie (automobile autour de Córdoba) complètent ce profil. Au début du XX\textsuperscript{e} siècle, l'Argentine comptait parmi les dix pays les plus riches du monde par habitant : son recul relatif depuis lors est un cas d'école.

Car l'économie argentine est marquée par une \cle{instabilité récurrente} : crises monétaires à répétition, inflation très élevée (parfois à deux ou trois chiffres), dévaluations successives du peso, défauts de paiement sur la dette publique (notamment en 2001). Ces crises détruisent l'épargne, découragent l'investissement et plongent périodiquement une partie de la population dans la pauvreté. L'Argentine illustre ainsi qu'un pays peut posséder d'immenses atouts sans parvenir à stabiliser sa trajectoire de développement.

\courssub{Des limites communes : inégalités, dépendance, endettement}

Malgré leurs différences, les trois foyers partagent des fragilités structurelles.

\begin{itemize}
\item \textbf{Des inégalités sociales criantes.} L'Amérique latine est l'une des régions les plus inégalitaires du monde. À Rio de Janeiro ou à São Paulo, les \cle{favelas} (quartiers pauvres et auto-construits) côtoient des quartiers d'affaires ultramodernes ; au Mexique, le Sud rural et indien contraste avec le Nord industriel ; en Argentine, les \emph{villas miserias} ceinturent Buenos Aires.
\item \textbf{Une dépendance aux cours des matières premières.} Soja, fer, pétrole, cuivre : quand les cours mondiaux s'effondrent, les recettes d'exportation et les budgets publics s'effondrent avec eux. La lecture du diagramme circulaire des exportations brésiliennes le montre clairement.
\item \textbf{L'endettement et la dépendance financière.} Les trois pays ont connu de graves crises de la dette (Mexique en 1982 et 1994, Brésil dans les années 1980, Argentine en 2001), qui ont imposé des politiques d'austérité coûteuses socialement.
\end{itemize}

\begin{attention}[Erreur fréquente : un Brésil « entièrement développé »]
Ne présente jamais le Brésil (ni le Mexique, ni l'Argentine) comme un pays développé achevé. C'est un pays \emph{émergent}, c'est-à-dire un territoire \cle{dual} : les gratte-ciel de São Paulo et l'industrie aéronautique coexistent avec les favelas, le Nordeste rural pauvre et des millions de travailleurs informels. Au Baccalauréat, un devoir qui oublie cette dualité perd l'essentiel de la notion d'émergence.
\end{attention}

\begin{methode}[Analyser une photographie : la favela sur fond de gratte-ciel]
Face à une photographie montrant une favela au premier plan et des tours de bureaux à l'arrière-plan (document classique au Bac), procède en trois étapes :
\begin{enumerate}
\item \textbf{Décrire objectivement} ce que tu vois, sans interpréter : au premier plan, des habitations précaires superposées sur une colline (matériaux de récupération, densité élevée) ; à l'arrière-plan, des immeubles modernes de verre et de béton.
\item \textbf{Dégager le contraste spatial} : deux mondes urbains se touchent physiquement mais sont séparés socialement et économiquement ; repère les indices (qualité des constructions, équipements, végétation).
\item \textbf{Conclure sur les inégalités} : la photographie illustre la dualité sociale et spatiale des métropoles des pays émergents, où la richesse mondialisée et la pauvreté cohabitent dans le même espace urbain.
\end{enumerate}
\end{methode}

\courssub{L'insertion dans la mondialisation}

Ces trois économies sont pleinement insérées dans les échanges mondiaux, mais selon des modalités variées.

\begin{itemize}
\item \textbf{L'intégration régionale} : le Brésil et l'Argentine sont membres du \cle{Mercosur} (Marché commun du Sud, créé en 1991, avec le Paraguay et l'Uruguay), qui organise le libre-échange sud-américain ; le Mexique, lui, est arrimé à l'Amérique du Nord par l'AEUMC.
\item \textbf{Le basculement vers la Chine} : la Chine est devenue le premier partenaire commercial du Brésil et un client majeur de l'Argentine (soja), illustrant la montée des échanges \cle{Sud-Sud}.
\item \textbf{La visibilité mondiale} : l'accueil de la Coupe du monde de football (2014) et des Jeux olympiques (2016) par le Brésil, comme celui du Mondial 2026 coorganisé par le Mexique, traduit la volonté de ces pays de compter sur la scène mondiale.
\end{itemize}

\begin{exoresolu}[Comparer les trois modèles d'émergence latino-américaine]
Énoncé : à partir du cours, complète le raisonnement suivant. « Le Brésil, le Mexique et l'Argentine sont trois foyers économiques en expansion, mais leur insertion dans la mondialisation repose sur des bases différentes. » Justifie cette affirmation en une dizaine de lignes.
\tcblower
\textbf{Solution détaillée.} Le Brésil fonde son émergence sur la diversité : agriculture commerciale de masse (premier exportateur mondial de soja, café, sucre), mines (fer), pétrole offshore (Petrobras) et industrie de pointe (Embraer), avec un vaste marché intérieur et un rôle diplomatique (BRICS, Mercosur). Le Mexique, lui, s'est inséré comme atelier manufacturier adossé aux États-Unis : les maquiladoras et l'ALENA/AEUMC ont fait des produits manufacturés l'essentiel de ses exportations, complétées par le pétrole et le tourisme. L'Argentine, enfin, mise sur sa puissance agroalimentaire (pampa, soja, bœuf, blé), mais son émergence est fragilisée par des crises monétaires récurrentes. Les trois pays partagent cependant des limites communes : inégalités sociales (favelas, villas miserias), dépendance aux cours des matières premières et endettement. Conclusion : l'émergence latino-américaine est réelle mais inachevée, car elle reste inégalement partagée et vulnérable aux fluctuations de l'économie mondiale.
\end{exoresolu}

\begin{aretenir}[L'essentiel de la section]
\begin{itemize}
\item Le \cle{Brésil}, première puissance d'Amérique latine et membre des BRICS, cumule agriculture commerciale, mines, pétrole offshore et industrie (Embraer, Petrobras) ; São Paulo concentre environ un tiers de son PIB.
\item Le \cle{Mexique} est un atelier manufacturier (maquiladoras) intégré au marché nord-américain (ALENA puis AEUMC), avec pétrole et tourisme en appoint.
\item L'\cle{Argentine} est une puissance agroalimentaire (pampa) minée par l'instabilité monétaire.
\item Limites communes : inégalités criantes, dépendance aux matières premières, endettement.
\end{itemize}
\end{aretenir}

\begin{exercice}[Pour t'entraîner]
\begin{enumerate}
\item Définis en deux phrases le terme \emph{maquiladora} et explique pourquoi ces usines se concentrent près de la frontière américano-mexicaine.
\item À partir du diagramme circulaire des exportations brésiliennes, calcule la part cumulée des produits agricoles et miniers, puis explique en cinq lignes pourquoi cette structure est une faiblesse pour le Brésil.
\item Rédige un commentaire de dix lignes de la photographie suivante (décrite) : « À Rio de Janeiro, une favela aux maisons de briques et de tôles s'accroche à une colline ; en contrebas, les gratte-ciel du quartier d'affaires. »
\end{enumerate}
\end{exercice}

\begin{corrige}[Exercice]
\begin{enumerate}
\item Une maquiladora est une usine de sous-traitance installée au Mexique, bénéficiant d'avantages fiscaux et douaniers, qui assemble des composants importés pour réexporter les produits finis. Elles se concentrent près de la frontière américaine pour minimiser les coûts et délais de transport vers le marché des États-Unis, principal débouché, tout en profitant d'une main-d'œuvre mexicaine moins chère.
\item Produits agricoles (environ 40 \%) + minerais et métaux (environ 25 \%) = environ 65 \% des exportations. Cette structure est une faiblesse car les cours mondiaux des matières premières sont volatils : une baisse des prix du soja ou du fer réduit mécaniquement les recettes en devises et les budgets publics, rendant la croissance brésilienne dépendante de facteurs extérieurs qu'elle ne maîtrise pas.
\item \emph{Description} : au premier plan, un habitat précaire et dense (briques, tôles, ruelles) occupe une colline ; à l'arrière-plan, des tours modernes de bureaux. \emph{Contraste} : deux espaces urbains juxtaposés mais socialement séparés ; la proximité géographique souligne l'écart de niveau de vie. \emph{Conclusion} : cette photographie illustre la dualité des métropoles des pays émergents, où la richesse issue de la mondialisation coexiste avec une pauvreté massive, signe que l'émergence brésilienne reste socialement inachevée.
\end{enumerate}
\end{corrige}

\coursec{Les pays émergents d'Asie et du Moyen-Orient}

Après avoir traversé l'Atlantique pour étudier le Brésil, le Mexique et l'Argentine, poursuivons notre tour du monde des foyers économiques en expansion vers l'est. Imagine un instant le trajet d'un porte-conteneurs parti de Shanghai : il longe les côtes de la Malaisie, passe devant Singapour, traverse l'océan Indien au large de l'Inde, franchit le détroit d'Ormuz non loin des pétroliers du Golfe, puis le canal de Suez avant de gagner l'Europe. Cette route maritime, la plus fréquentée du monde, borde précisément les foyers émergents que nous allons étudier : l'\cle{Inde}, les quatre \cle{Dragons}, la \cle{Malaisie} et les \cle{pays du Golfe}. Ce n'est pas un hasard : la position sur les grandes routes commerciales est l'un des facteurs majeurs de l'émergence.

\courssub{L'Inde : le géant démographique devenu puissance des services}

\textbf{Une démographie et une économie de géant.} Avec environ 1,4 milliard d'habitants, l'Inde est devenue en 2023 le pays le plus peuplé du monde, devant la Chine. Son PIB total la place parmi les cinq premières économies mondiales (ordre de grandeur : 3 500 milliards de dollars, soit environ 80 fois le PIB du Cameroun). Mais ce PIB, divisé par une population immense, donne un revenu par habitant encore modeste (environ 2 500 dollars par an, ordre de grandeur) : l'Inde illustre parfaitement le paradoxe des pays émergents, à la fois puissants collectivement et pauvres individuellement.

\textbf{La puissance des services : Bangalore et l'externalisation.} Depuis les réformes d'ouverture de 1991, l'Inde s'est spécialisée dans les \cle{services}, et notamment dans l'informatique. La ville de Bangalore (Bengaluru), dans le sud du pays, concentre des milliers d'entreprises de logiciels et de services numériques (Infosys, Wipro, Tata Consultancy Services). Le mécanisme est celui de l'\cle{externalisation} (ou \emph{offshoring}) : les entreprises des pays développés confient à des sociétés indiennes des tâches — développement de logiciels, centres d'appels, comptabilité — parce que la main-d'œuvre indienne y est \textbf{qualifiée, anglophone et bon marché}. Un ingénieur informaticien indien coûte trois à cinq fois moins cher qu'un ingénieur européen à compétence comparable.

\begin{exemplebox}[Bangalore, la « Silicon Valley indienne »]
L'Inde est le premier exportateur mondial de services informatiques externalisés : ses exportations de services informatiques dépassent 150 milliards de dollars par an (ordre de grandeur). Bangalore emploie plus d'un million de personnes dans le secteur numérique. Quand tu appelles le service client d'une grande entreprise européenne, il est possible que ton interlocuteur soit à Bangalore ou à Hyderabad.
\end{exemplebox}

\textbf{Des industries de pointe : pharmacie et espace.} L'Inde ne se contente pas des services. Elle est surnommée la « pharmacie du monde » : premier producteur mondial de \textbf{médicaments génériques} à bas coût, elle fournit une grande partie des vaccins et des médicaments consommés en Afrique, y compris au Cameroun. Son agence spatiale, l'ISRO, a réussi en août 2023 à faire atterrir une sonde (Chandrayaan-3) près du pôle sud de la Lune, exploit que peu de pays ont réalisé.

\textbf{Mais une pauvreté massive.} L'émergence indienne reste profondément inégale : plusieurs centaines de millions d'Indiens vivent encore avec moins de quelques dollars par jour, l'agriculture emploie plus de 40 \% de la population active avec une faible productivité, et les bidonvilles des mégapoles (Mumbai, Delhi, Kolkata) côtoient les quartiers d'affaires ultramodernes. L'Inde est donc un pays émergent \emph{inachevé}, où la modernité la plus avancée coexiste avec une pauvreté de masse.

\courssub{Les quatre Dragons : des émergents devenus économies avancées}

\textbf{Un modèle d'industrialisation exportatrice.} La Corée du Sud, Taïwan, Singapour et Hong Kong sont surnommés les quatre \cle{Dragons} (on parlait aussi de « nouveaux pays industrialisés » ou \cle{NPI} dans les années 1970-1980). Partis d'un niveau de développement comparable à celui de nombreux pays africains dans les années 1960, ils ont suivi une stratégie claire : l'\cle{industrialisation exportatrice}. Plutôt que de produire pour leur seul marché intérieur, ils ont fabriqué massivement des biens (textile, puis électronique, puis produits de haute technologie) destinés à être vendus sur le marché mondial, en s'appuyant sur une main-d'œuvre disciplinée et bon marché, une éducation de masse et des États volontaristes qui ont orienté le crédit et protégé les industries naissantes.

\textbf{Devenus des économies avancées.} Le résultat est spectaculaire : en deux générations, ces territoires ont rejoint le club des économies développées. La Corée du Sud est aujourd'hui une puissance technologique mondiale (Samsung pour l'électronique, Hyundai pour l'automobile) ; Taïwan domine la production mondiale de semi-conducteurs grâce à TSMC, qui fabrique les puces les plus avancées de la planète ; Singapour est l'un des premiers ports et centres financiers du monde ; Hong Kong est une place financière majeure. Leur revenu par habitant dépasse 30 000 dollars par an, et leur IDH est très élevé.

\begin{attention}[Erreur fréquente : classer les Dragons parmi les pays pauvres]
Au Baccalauréat, une erreur classique consiste à ranger encore la Corée du Sud, Taïwan, Singapour et Hong Kong parmi les « pays pauvres » ou les « pays en voie de développement ». C'est faux : ces économies ont un \textbf{IDH très élevé} et un \textbf{revenu par habitant de niveau développé}, parfois supérieur à celui de pays européens. Elles illustrent au contraire qu'un pays émergent peut \emph{achever} son émergence et rejoindre le groupe des pays développés. Les Dragons sont un modèle de réussite, pas un exemple de sous-développement.
\end{attention}

\courssub{La Malaisie : des matières premières à l'industrie électronique}

La Malaisie offre un autre visage de l'émergence asiatique : celui d'une \textbf{montée en gamme progressive}. Ancienne colonie britannique indépendante depuis 1957, elle a d'abord fondé sa richesse sur l'exportation de matières premières tropicales : le \cle{caoutchouc} (elle fut longtemps premier producteur mondial), l'\cle{huile de palme} (dont elle reste le deuxième producteur mondial derrière l'Indonésie), l'étain et le pétrole. À partir des années 1970-1980, l'État malaisien a attiré les investissements directs étrangers (IDE) dans l'\textbf{industrie électronique} : assemblage de composants, d'appareils électriques et de semi-conducteurs, notamment dans la région de Penang. Les tours Petronas de Kuala Lumpur, qui furent les plus hauts édifices du monde à leur achèvement en 1998, symbolisent cette réussite : financées par la compagnie pétrolière nationale, elles affichent l'ambition d'un pays passé du statut de fournisseur de matières premières à celui de \cle{pays à revenu intermédiaire} industrialisé. La Malaisie est aujourd'hui l'une des économies les plus prospères d'Asie du Sud-Est, avec un revenu par habitant voisin de 11 000 dollars (ordre de grandeur).

\courssub{Les pays du Golfe : la rente pétrolière investie dans la diversification}

\textbf{La rente pétrolière et gazière.} L'Arabie Saoudite et les Émirats Arabes Unis (EAU) doivent leur émergence à une ressource unique : les hydrocarbures. L'Arabie Saoudite est l'un des trois premiers producteurs mondiaux de pétrole ; les EAU (fédération de sept émirats, dont Abou Dhabi et Dubaï) possèdent d'importantes réserves de pétrole et de gaz. Cette \cle{rente pétrolière} a permis d'accumuler d'immenses capitaux et de moderniser rapidement des pays qui, il y a soixante ans, étaient des déserts peuplés de bédouins et de pêcheurs de perles.

\textbf{La diversification : le modèle de Dubaï.} Conscients que le pétrole s'épuisera un jour, les pays du Golfe investissent massivement leur rente dans d'autres secteurs. Dubaï est le cas le plus spectaculaire : l'émirat s'est transformé en \textbf{centre financier} international, en destination \textbf{touristique} de luxe (hôtels emblématiques, îles artificielles, tour Burj Khalifa, la plus haute du monde avec 828 m) et en \cle{hub aérien} mondial grâce à la compagnie Emirates et à l'un des aéroports les plus fréquentés de la planète. L'Arabie Saoudite poursuit une stratégie analogue avec son plan « Vision 2030 ».

\begin{savaistu}[Dubaï ne vit plus du pétrole]
Contrairement à une idée reçue, Dubaï ne vit plus principalement du pétrole : les hydrocarbures représentent \textbf{moins de 5 \% de son PIB}. L'essentiel de sa richesse provient désormais du commerce (réexportation), de la finance, de l'immobilier et du tourisme. C'est Abou Dhabi, l'émirat voisin, qui détient l'essentiel des réserves pétrolières des EAU. Dubaï est ainsi devenu le symbole mondial de la reconversion réussie d'une rente pétrolière.
\end{savaistu}

\begin{definition}[Fonds souverain]
Un \cle{fonds souverain} est un fonds d'investissement \textbf{public}, alimenté par les excédents d'un État — souvent la rente pétrolière ou gazière — et qui place ces capitaux \textbf{à l'étranger} (actions, immobilier, infrastructures, entreprises) pour faire fructifier la richesse nationale et préparer l'après-pétrole. Exemples : le fonds d'Abou Dhabi (ADIA), l'un des plus importants du monde, le fonds saoudien PIF, ou encore le fonds norvégien.
\end{definition}

\courssub{Une lecture cartographique : les foyers émergents sur la route Chine--Europe}

\begin{popfigure}
\includegraphics[width=\linewidth,height=8cm,keepaspectratio]{N14/shipping_routes.jpg}
\legende{La route maritime Chine--Europe sur la carte mondiale du trafic maritime commercial (les continents sont en noir ; plus le bleu est foncé, plus le trafic est dense). On suit l'axe le plus fréquenté du monde : de l'Asie orientale au détroit de Malacca, à l'océan Indien, puis à la mer Rouge et à Suez, à la Méditerranée et à l'Europe du Nord-Ouest. Le long de cet axe se trouvent les foyers émergents d'Asie et du Moyen-Orient : Dragons, Inde, Malaisie, Golfe ; Singapour, au débouché du détroit de Malacca, est un point de passage obligé. La carte ne distingue pas ces pays. \\ Source : Wikimedia Commons, T. Hengl, CC BY-SA 3.0.}
\end{popfigure}

\courssub{La skyline de Dubaï : photographie commentée}

Observons maintenant un document emblématique : la vue de Dubaï, que l'on peut analyser comme une photographie de paysage urbain au Baccalauréat.

\begin{popfigure}
\includegraphics[width=\linewidth,height=8cm,keepaspectratio]{N14/dubai_mer.jpg}
\legende{La skyline de Dubaï vue de la mer (photographie). Lecture : (1) la Burj Khalifa, au centre, plus haute tour du monde, vitrine technologique ; (2) un quartier d'affaires dense de gratte-ciel, signe de la fonction financière ; (3) au premier plan, la façade maritime du golfe Persique, rappel du rôle de hub commercial ; (4) des grues encore visibles : une ville construite en quelques décennies grâce aux capitaux de la rente. Cette skyline matérialise la reconversion des rentes pétrolières du Golfe dans les services mondialisés. \\ Source : Wikimedia Commons, Jpbowen, CC BY-SA 4.0.}
\end{popfigure}

\courssub{Facteurs communs et limites de l'émergence asiatique}

\textbf{Les facteurs communs.} Malgré leurs différences, l'Inde, les Dragons, la Malaisie et les pays du Golfe partagent plusieurs ressorts d'émergence :
\begin{itemize}
\item une \textbf{main-d'œuvre abondante et bon marché} (Inde, Dragons au départ, Malaisie), qui attire les entreprises étrangères ;
\item l'\textbf{ouverture aux investissements directs étrangers} (IDE) et l'intégration aux chaînes de valeur mondiales ;
\item des \textbf{stratégies étatiques volontaristes} : réformes indiennes de 1991, plans quinquennaux coréens, zones franches malaisiennes et de Dubaï, « Vision 2030 » saoudienne ;
\item une \textbf{position privilégiée sur les routes maritimes mondiales} : détroit de Malacca (Singapour, Malaisie), océan Indien (Inde), détroit d'Ormuz (Golfe).
\end{itemize}

\textbf{Les limites.} L'émergence asiatique et moyen-orientale connaît aussi de réelles fragilités : la \textbf{dépendance pétrolière} des pays du Golfe (que la diversification ne fait qu'atténuer) ou la \textbf{dépendance technologique} de certains pays vis-à-vis des firmes étrangères ; des \textbf{inégalités} criantes (pauvreté massive en Inde, ouvriers migrants aux droits très limités sur les chantiers du Golfe) ; des \textbf{conditions de travail} parfois très dures dans les usines d'assemblage ; enfin des \textbf{tensions géopolitiques} (rivalité Inde--Chine, question de Taïwan, instabilité du Moyen-Orient) qui menacent la stabilité de ces foyers.

\begin{methode}[Interpréter un tableau statistique comparatif]
Pour analyser un tableau de données chiffrées sur les pays émergents :
\begin{enumerate}
\item \textbf{Repérer l'unité et l'année} des données (milliards de dollars ? dollars par habitant ? quelle date ?) : comparer des grandeurs d'unités ou d'années différentes est une erreur grave.
\item \textbf{Classer les valeurs} (du plus fort au plus faible) pour dégager une hiérarchie.
\item \textbf{Calculer des rapports} pour rendre les écarts parlants : par exemple, le PIB de l'Inde (environ 3 500 milliards de dollars) est environ 80 fois celui du Cameroun (environ 45 milliards de dollars).
\item \textbf{Croiser les indicateurs} : un PIB total élevé peut cacher un revenu par habitant faible (cas de l'Inde) ; toujours confronter richesse globale et richesse par habitant.
\end{enumerate}
\end{methode}

\begin{exoresolu}[Comparer deux foyers émergents]
Énoncé : le PIB de l'Inde est d'environ 3 500 milliards de dollars pour 1,4 milliard d'habitants ; celui de l'Arabie Saoudite d'environ 1 000 milliards de dollars pour 36 millions d'habitants. Calcule le PIB par habitant de chaque pays et commente.
\tcblower
Solution détaillée :
\begin{align*}
\text{Inde : } & \frac{3500 \text{ milliards}}{1400 \text{ millions}} \approx 2\,500 \text{ dollars/habitant} \\
\text{Arabie Saoudite : } & \frac{1000 \text{ milliards}}{36 \text{ millions}} \approx 27\,800 \text{ dollars/habitant}
\end{align*}
Commentaire : le PIB total de l'Inde est 3,5 fois supérieur à celui de l'Arabie Saoudite, ce qui en fait une puissance économique globale plus importante. Mais son PIB par habitant est environ 11 fois inférieur. L'Inde est donc un géant économique dont la population reste en moyenne pauvre, tandis que l'Arabie Saoudite est une puissance moyenne dont la rente pétrolière, répartie sur une population réduite, assure un niveau de vie élevé. Ce calcul montre pourquoi il faut toujours croiser PIB total et PIB par habitant pour qualifier un pays émergent.
\end{exoresolu}

\begin{aretenir}[L'essentiel de la section]
L'Asie et le Moyen-Orient concentrent des foyers émergents aux profils variés : l'\textbf{Inde}, géant démographique devenu puissance mondiale des services (Bangalore) et de la pharmacie, mais confrontée à une pauvreté massive ; les quatre \textbf{Dragons} (Corée du Sud, Taïwan, Singapour, Hong Kong), qui ont achevé leur émergence grâce à l'industrialisation exportatrice ; la \textbf{Malaisie}, passée des matières premières à l'électronique ; les \textbf{pays du Golfe} (Arabie Saoudite, EAU), qui convertissent leur rente pétrolière en finance, tourisme et hubs mondiaux via des fonds souverains. Facteurs communs : main-d'œuvre, IDE, États volontaristes et position maritime. Limites : dépendances, inégalités et tensions géopolitiques.
\end{aretenir}

\begin{exercice}[Pour t'entraîner]
À partir de la carte schématique de la leçon : 1) cite les quatre Dragons et localise-les ; 2) explique pourquoi Singapour occupe une position stratégique ; 3) rédige un commentaire de dix lignes du schéma de la skyline de Dubaï en montrant ce qu'il révèle de la stratégie des pays du Golfe.
\end{exercice}

\coursec{Les pays émergents d'Afrique : Afrique du Sud, Nigéria, Égypte, Maroc}

Après avoir analysé les dynamiques américaines, asiatiques et moyen-orientales, tournons notre regard vers le continent africain. Longtemps perçu uniquement comme une réserve de matières premières, l'Afrique voit émerger, depuis le début des années 2000, des pôles de croissance qui restructurent les échanges intra-africains et attirent les investissements mondiaux. Quatre pays concentrent l'essentiel de cette dynamique : l'\cle{Afrique du Sud}, le \cle{Nigéria}, l'\cle{Égypte} et le \cle{Maroc}. Ils ne forment pas un bloc homogène, mais chacun, à sa manière, incarne une voie d'insertion dans la mondialisation. Le lancement de la \cle{ZLECAf} (Zone de libre-échange continentale africaine, entrée en vigueur en 2021) vise précisément à relier ces foyers entre eux.

\courssub{L'Afrique du Sud : la locomotive industrielle du continent}

L'Afrique du Sud est la première puissance industrielle d'Afrique subsaharienne et le seul membre africain des \cle{BRICS} (Brésil, Russie, Inde, Chine, Afrique du Sud), qu'elle a rejoints en 2011. Son économie, la plus diversifiée du continent, repose sur un socle minier exceptionnel et un secteur tertiaire sophistiqué.

\begin{exemplebox}[L'Afrique du Sud, puissance minière et financière]
Le sous-sol sud-africain recèle les plus grandes réserves mondiales de \textbf{platine} (environ 90\,\% des réserves mondiales connues), ainsi que d'importantes réserves d'or, de chrome et de manganèse. Le bassin du Witwatersrand a fait de Johannesburg la capitale économique du pays. La Bourse de Johannesburg (\cle{JSE}) est la plus grande place financière d'Afrique par sa capitalisation boursière (de l'ordre de 1 000 milliards USD, ordre de grandeur). Des groupes comme \textbf{Anglo American}, \textbf{Sasol} (chimie, énergie) ou \textbf{Naspers} (médias, technologie) y ont leur siège et rayonnent sur tout le continent.
\end{exemplebox}

Pourtant, cette puissance structurelle masque des failles profondes, héritées de l'\cle{apartheid} (1948-1994). Le taux de chômage officiel dépasse 30\,\% (plus de 40\,\% au sens large), touchant massivement la jeunesse noire. Les inégalités (coefficient de Gini voisin de 0,63, parmi les plus élevés au monde) se lisent dans l'espace : townships paupérisés (Soweto, Khayelitsha) face aux quartiers résidentiels sécurisés (Sandton). La crise énergétique (\emph{load shedding} : délestages quotidiens) freine l'industrie depuis la fin des années 2000, révélant la vétusté des centrales à charbon d'\textbf{Eskom}.

\begin{attention}[Piège : confondre puissance industrielle et développement inclusif]
Ne pas écrire : « L'Afrique du Sud est un pays développé ». C'est une \emph{économie émergente} à revenu intermédiaire de la tranche supérieure (PIB/habitant de l'ordre de 6 000 USD en nominal), mais son IDH (environ 0,71) reste moyen. Au Bac, nuancez toujours : « Première puissance industrielle africaine, mais freinée par des inégalités structurelles et une crise énergétique chronique ».
\end{attention}

\courssub{Le Nigéria : le géant démographique et pétrolier}

Avec plus de \textbf{220 millions d'habitants} (estimation 2024), le Nigéria est le pays le plus peuplé d'Afrique et l'une des six premières populations mondiales. Il figure, avec l'Afrique du Sud et l'Égypte, parmi les trois premières économies du continent (PIB nominal de l'ordre de 360 à 390 milliards USD en 2023, selon les sources et les effets de change), porté par la rente pétrolière.

\begin{exemplebox}[Le delta du Niger : cœur énergétique et zone de tensions]
Le pétrole représente environ 90\,\% des recettes d'exportation et près de la moitié des recettes budgétaires. Les terminaux d'exportation (Bonny, Forcados) et les raffineries (Port Harcourt, Warri, et la nouvelle méga-raffinerie \textbf{Dangote} à Lekki, d'une capacité de 650 000 barils/jour, entrée en service progressif à partir de 2023) concentrent l'activité. Mais cette rente alimente la corruption (le Nigéria se classe parmi les pays les moins bien notés de l'indice de perception de la corruption de Transparency International) et a nourri des conflits armés (MEND, vols de pétrole ou \emph{bunkering} illégal) qui polluent gravement les mangroves du delta.
\end{exemplebox}

\begin{attention}[Erreur fréquente : « Le Nigéria est riche parce qu'il a du pétrole »]
C'est une confusion entre \textbf{richesse nationale} (PIB) et \textbf{bien-être de la population}. Environ \textbf{40\,\% des Nigérians vivent sous le seuil de pauvreté monétaire} (moins de 2,15 USD/jour selon le seuil international révisé de la Banque mondiale). La rente est captée par une élite politico-économique ; la redistribution échoue (faible pression fiscale, absence d'État-providence). Au Bac, écrivez : « Une rente pétrolière mal redistribuée, source de croissance sans développement inclusif ».
\end{attention}

\begin{savaistu}[Nollywood : la « Dream Factory » de Lagos]
Le Nigéria produit environ \textbf{2 500 films par an}, ce qui en fait l'une des toutes premières industries cinématographiques mondiales en volume, aux côtés de Bollywood (Inde) et devant Hollywood en nombre de films. \cle{Nollywood} génère des centaines de millions de dollars, emploie plus d'un million de personnes (acteurs, techniciens, distributeurs informels) et exporte sa culture (vidéos, musique \emph{Afrobeats} avec Burna Boy, Wizkid, Davido) dans toute la diaspora africaine et au-delà. C'est un atout majeur de \cle{soft power} nigérian.
\end{savaistu}

Au-delà du pétrole, Lagos (plus de 20 millions d'habitants dans son aire urbaine) devient un pôle \cle{fintech} africain (\textbf{Flutterwave}, \textbf{Paystack}, rachetée par Stripe). L'agriculture (cacao, caoutchouc, igname) et l'industrie naissante (ciment Dangote, engrais) tentent de diversifier l'économie, mais l'insécurité (Boko Haram au Nord-Est, banditisme au Nord-Ouest, tensions séparatistes au Sud-Est) freine l'investissement.

\courssub{L'Égypte : carrefour stratégique entre Afrique, Europe et Asie}

L'Égypte (environ 110 millions d'habitants) occupe une position géostratégique unique : le \cle{canal de Suez} (193 km) relie la Méditerranée à la mer Rouge, raccourcissant considérablement la route maritime Europe-Asie. Environ 12\,\% du commerce maritime mondial y transite en temps normal.

\begin{exemplebox}[Le canal de Suez : une rente de situation modernisée]
L'élargissement achevé en 2015 (nouveau chenal parallèle d'environ 35 km) permet le croisement des navires sur une partie du parcours et réduit le temps de transit. Le trafic a dépassé \textbf{25 000 navires par an}, générant un record d'environ \textbf{9,4 milliards USD} de recettes pour l'exercice 2022-2023. La \cle{Zone économique du canal de Suez (SCZone)} attire des zones franches (industries textile, automobile, agroalimentaire) pour capter de la valeur ajoutée au-delà du simple péage. À noter : les attaques en mer Rouge fin 2023-2024 ont fortement perturbé ce trafic, illustrant la fragilité d'une rente de situation.
\end{exemplebox}

L'agriculture, concentrée sur environ 4\,\% du territoire (vallée et delta du Nil), est hyper-productive (coton, canne à sucre, fruits, riz) grâce à l'irrigation, mais l'Égypte importe plus de la moitié de son blé : c'est une dépendance alimentaire stratégique, révélée par la guerre en Ukraine. L'industrie (textile, ciment, chimie, assemblage automobile), le gaz naturel (champ offshore de Zohr) et le tourisme (pyramides, mer Rouge, Louxor ; environ 15 millions de visiteurs en 2023) sont les autres piliers. Le \cle{barrage de la Renaissance (GERD)} construit par l'Éthiopie sur le Nil Bleu inquiète Le Caire pour sa sécurité hydrique : c'est une source majeure de tensions diplomatiques.

\courssub{Le Maroc : l'atelier automobile et aéronautique, hub logistique et énergétique}

Le Maroc (environ 37 millions d'habitants) a fait le pari réussi de l'industrialisation par l'intégration aux chaînes de valeur mondiales (CVM), notamment européennes. Il est devenu le \textbf{premier producteur de voitures particulières d'Afrique} (capacité installée de l'ordre de 700 000 véhicules par an, objectif 1 million).

\begin{exemplebox}[Tanger Med : la porte de l'Afrique sur l'Atlantique]
Inauguré en 2007, le port \textbf{Tanger Med} est le \textbf{premier port à conteneurs d'Afrique et de Méditerranée}, avec plus de \textbf{10 millions d'EVP} (Équivalent Vingt Pieds) traités en 2024, après 8,6 millions en 2023. Il relie près de 180 ports dans le monde. La zone franche attenante héberge \textbf{Renault} (usine de Melloussa, Dacia Sandero), \textbf{Stellantis} (Kénitra, Peugeot 208), et un écosystème de centaines d'équipementiers (câblage, sièges, pneumatiques). Le taux d'intégration locale dépasse 65\,\%.
\end{exemplebox}

Le Maroc domine aussi le marché mondial des \textbf{phosphates} (environ 70\,\% des réserves mondiales connues, via l'\cle{OCP} -- Office Chérifien des Phosphates), en s'orientant vers la transformation locale (engrais phosphatés, acide phosphorique) plutôt que l'export brut. En énergie, le complexe solaire \textbf{Noor Ouarzazate} (environ 580 MW, l'un des plus grands au monde) symbolise la stratégie de transition énergétique (objectif de plus de 50\,\% d'énergies renouvelables dans le mix électrique à l'horizon 2030). Le secteur aéronautique (Boeing, Safran, Hexcel à Casablanca/Nouaceur) emploie environ 20 000 personnes.

\courssub{Facteurs et limites de l'émergence africaine}

Ces quatre foyers partagent des moteurs communs, mais butent sur des obstacles structurels similaires.

\begin{definition}[Classe moyenne africaine]
Catégorie de population dont les dépenses journalières se situent, selon la définition de la Banque africaine de développement (BAD), entre \textbf{2 et 20 USD par jour en parité de pouvoir d'achat} (la tranche 4-20 USD étant la plus stable). Elle représente environ \textbf{350 millions de personnes} en Afrique (ordre de grandeur, soit près d'un tiers de la population). Sa consommation (télécoms, logement, éducation, santé, automobile) tire la demande intérieure et attire la grande distribution (Carrefour, Shoprite, Label'Vie) et les services financiers (mobile money : M-Pesa, Orange Money, MTN Mobile Money).
\end{definition}

\begin{aretenir}[Les 4 moteurs de l'émergence africaine]
\begin{enumerate}
    \item \textbf{Ressources naturelles :} hydrocarbures (Nigéria, Algérie, Angola), minerais (Afrique du Sud, RDC, Zambie, Guinée), phosphates (Maroc), terres arables irriguées (vallée du Nil).
    \item \textbf{Dividende démographique :} population en âge de travailler en croissance rapide (environ 1,4 milliard d'Africains en 2023, 2,5 milliards projetés en 2050). Urbanisation galopante (Lagos, Le Caire, Kinshasa, Luanda, Johannesburg).
    \item \textbf{Montée de la classe moyenne :} moteur de la consommation de masse, de l'immobilier et des services (banque, assurance, télécoms).
    \item \textbf{IDE chinois et nouvelles routes de la soie :} la Chine est le premier partenaire commercial de l'Afrique (échanges sino-africains de l'ordre de 280 milliards USD en 2023). Infrastructures (chemins de fer Addis-Abeba--Djibouti et Mombasa--Nairobi, ports de Kribi au Cameroun, Lekki au Nigéria), zones industrielles (Éthiopie, Égypte).
\end{enumerate}
\end{aretenir}

\begin{propriete}[Limites structurelles de l'émergence africaine]
Malgré la croissance, l'émergence reste \textbf{fragile et inachevée} :
\begin{itemize}
    \item \textbf{Dépendance aux matières premières :} les produits primaires dominent les exportations de la plupart des pays africains, d'où une vulnérabilité aux chocs de prix (pétrole, cuivre, cacao).
    \item \textbf{Faible industrialisation :} la part de l'industrie manufacturière dans le PIB africain reste voisine de 10\,\% (contre environ 25\,\% en Asie du Sud-Est). Importation massive de biens manufacturés (médicaments, machines, véhicules).
    \item \textbf{Dette publique et risque de change :} endettement en devises (eurobonds, prêts chinois) ; la dépréciation des monnaies (naira, livre égyptienne, rand) alourdit le fardeau de la dette.
    \item \textbf{Gouvernance et instabilité :} coups d'État (Sahel), conflits armés (Soudan, RDC, Sahel, Nord du Mozambique), corruption et insécurité juridique freinent les IDE productifs.
    \item \textbf{Croissance démographique non maîtrisée :} fécondité élevée (plus de 4 enfants par femme en moyenne en Afrique subsaharienne) ; forte pression sur l'emploi (10 à 12 millions de jeunes entrent chaque année sur le marché du travail), l'éducation, la santé et l'alimentation.
\end{itemize}
\end{propriete}

\begin{methode}[Construire un croquis de synthèse : « L'Afrique : foyers émergents et corridors d'intégration »]
\begin{enumerate}
    \item \textbf{Contour et repères :} tracez le contour simplifié de l'Afrique (polygone grossier). Placez les repères incontournables : canal de Suez (NE), détroit de Gibraltar (N), Cap de Bonne-Espérance (S), golfe de Guinée (O), Corne de l'Afrique (E).
    \item \textbf{Localisation des 4 foyers :} placez un symbole distinct (carré, losange, triangle, cercle) pour \textbf{Johannesburg/Pretoria} (Afrique du Sud), \textbf{Lagos/Abuja} (Nigéria), \textbf{Le Caire/Alexandrie} (Égypte), \textbf{Casablanca/Tanger/Rabat} (Maroc). Nommez les villes.
    \item \textbf{Flux et corridors :} dessinez des flèches épaisses pour les corridors majeurs :
    \begin{itemize}
        \item Corridor Nord-Sud : Le Caire $\rightarrow$ Khartoum $\rightarrow$ Addis-Abeba $\rightarrow$ Nairobi $\rightarrow$ Lusaka $\rightarrow$ Johannesburg (axe route/rail, en partie projeté).
        \item Corridor transsaharien : Lagos $\rightarrow$ Niamey $\rightarrow$ Alger (projet route/rail).
        \item Corridor atlantique : Tanger $\rightarrow$ Casablanca $\rightarrow$ Dakar $\rightarrow$ Abidjan $\rightarrow$ Lagos (axe côtière, projet de gazoduc Nigéria-Maroc).
    \end{itemize}
    \item \textbf{Ressources et ports :} ajoutez des symboles pour les mines (or et platine en Afrique du Sud, pétrole au Nigéria, phosphates au Maroc) et les ports hubs (Tanger Med, Durban, Lagos/Lekki, Port-Saïd/Suez, Mombasa, Djibouti).
    \item \textbf{Légende organisée :} trois colonnes : \emph{signes ponctuels} (villes, mines, ports), \emph{signes linéaires} (axes routiers, ferroviaires, pipelines), \emph{signes de surface} (ZLECAf, CEDEAO, SADC, COMESA). Titre précis, échelle indicative, flèche du Nord.
\end{enumerate}
\end{methode}

\begin{popfigure}
\includegraphics[width=\linewidth,height=11cm,keepaspectratio]{N14/afrique_pol.jpg}
\legende{Carte politique de l'Afrique (en anglais) : États, capitales (étoiles) et principales villes, du Maroc à l'Afrique du Sud, avec le golfe de Guinée à l'ouest et la mer Rouge à l'est. Elle permet de localiser les quatre foyers émergents d'Afrique (Nigéria, Égypte, Afrique du Sud, Maroc) et leurs voisins ; les corridors d'intégration de la ZLECAf sont à tracer par l'élève entre ces pôles (Le Caire, Lagos et Abuja, Johannesburg et Pretoria, Casablanca et Rabat), car la carte n'en représente aucun. \\ Source : Wikimedia Commons, Central Intelligence Agency (États-Unis), domaine public.}
\end{popfigure}

\begin{popfigure}
\begin{tikzpicture}
\begin{axis}[
    ybar,
    bar width=14pt,
    width=\linewidth,
    height=9cm,
    enlargelimits=0.2,
    ylabel={PIB (milliards USD, 2023)},
    symbolic x coords={Afrique du Sud, Nigéria, Égypte, Maroc},
    xtick=data,
    nodes near coords,
    nodes near coords align={vertical},
    every node near coord/.append style={font=\tiny, popdark},
    ymin=0,
    ymax=1900,
    ymajorgrids=true,
    grid style={popline},
    tick label style={font=\small},
    label style={font=\small},
    axis line style={popdark},
    tick style={popdark},
    legend style={font=\small, at={(0.5,-0.18)}, anchor=north, legend columns=2},
]
\addplot[fill=popblue!70, draw=popblue] coordinates {
    (Afrique du Sud, 377)
    (Nigéria, 363)
    (Égypte, 396)
    (Maroc, 141)
};
\addplot[fill=poporange!70, draw=poporange] coordinates {
    (Afrique du Sud, 994)
    (Nigéria, 1390)
    (Égypte, 1700)
    (Maroc, 380)
};
\legend{PIB nominal, PIB en PPA}
\end{axis}
\end{tikzpicture}
\legende{Comparaison du PIB nominal et du PIB en parité de pouvoir d'achat (PPA) des quatre foyers émergents africains (ordres de grandeur). En nominal, les trois premières économies (Nigéria, Égypte, Afrique du Sud) sont proches ; en PPA, l'Égypte et le Nigéria dominent grâce à leur démographie. L'Afrique du Sud reste la plus industrialisée et la plus riche par habitant. Source : FMI, \emph{World Economic Outlook}, 2024 (estimations).}
\end{popfigure}

\begin{exoresolu}[Analyse comparative des trajectoires d'émergence]
\textbf{Énoncé :} À partir des documents ci-dessus (croquis, graphique, encadrés) et de vos connaissances, comparez les modèles d'émergence de l'Afrique du Sud et du Maroc. Montrez en quoi ils illustrent deux voies distinctes d'insertion dans la mondialisation.
\tcblower
\textbf{Solution détaillée :}
\begin{enumerate}
    \item \textbf{Identification des modèles :}
    \begin{itemize}
        \item \textbf{Afrique du Sud :} modèle « rente minière + services financiers avancés ». Héritage minier (or, platine) + place financière (JSE) + industrie lourde (sidérurgie, chimie avec Sasol, automobile BMW/Ford/VW pour l'export). Insertion par l'export de minerais à haute valeur et de services financiers vers le monde. Membre des BRICS : voix politique Sud-Sud.
        \item \textbf{Maroc :} modèle « plateforme d'assemblage + transformation des phosphates + tourisme/énergies vertes ». Insertion par l'intégration aux chaînes de valeur européennes (automobile Renault/Stellantis, aéronautique) grâce à la proximité géographique, aux accords de libre-échange (UE, États-Unis, Turquie) et à des infrastructures portuaires de classe mondiale (Tanger Med). Transformation locale du phosphate (OCP) : montée en gamme.
    \end{itemize}
    \item \textbf{Similitudes (facteurs communs) :} stabilité politique relative, infrastructures portuaires et logistiques performantes, cadres réglementaires attractifs pour les IDE (zones franches), classe moyenne urbaine dynamique.
    \item \textbf{Différences structurelles :}
    \begin{itemize}
        \item \textbf{Dépendances :} l'Afrique du Sud dépend du cycle des matières premières (cours du platine et de l'or) et de l'électricité (charbon). Le Maroc dépend de l'agriculture (pluviométrie), de la demande automobile européenne et des importations énergétiques, mais se diversifie (solaire Noor).
        \item \textbf{Inégalités :} Afrique du Sud : inégalités extrêmes, héritées de l'apartheid (Gini $\approx$ 0,63). Maroc : inégalités territoriales (urbain/rural) et sociales, mais Gini plus faible ($\approx$ 0,39).
        \item \textbf{Intégration régionale :} l'Afrique du Sud est le moteur de la SADC et le pont vers l'Afrique australe. Le Maroc est un hub vers l'Afrique de l'Ouest (projet de gazoduc Nigéria-Maroc, banques Attijariwafa et BMCE présentes au Sahel) et a réintégré l'Union africaine en 2017.
    \end{itemize}
    \item \textbf{Conclusion synthétique :} l'Afrique du Sud incarne une \textbf{émergence par la rente minière et la finance}, freinée par une transformation structurelle inachevée. Le Maroc incarne une \textbf{émergence par l'industrialisation sous-traitée (\emph{nearshoring}) et la logistique}, plus vulnérable à la conjoncture européenne mais créatrice d'emplois manufacturiers massifs.
\end{enumerate}
\end{exoresolu}

\begin{exercice}[Entraînement Bac : commentaire de document statistique]
\textbf{Document :} extrait adapté du rapport \emph{African Economic Outlook 2024} (BAD) et du FMI (ordres de grandeur).
\begin{center}
\begin{tabularx}{\linewidth}{|l|X|X|X|X|}\hline
\rowcolor{popblueL}
\textbf{Indicateur 2023} & \textbf{Afrique du Sud} & \textbf{Nigéria} & \textbf{Égypte} & \textbf{Maroc} \\ \hline
PIB nominal (Md USD) & 377 & 363 & 396 & 141 \\ \hline
PIB/hab. PPA (USD) & 16 600 & 6 100 & 17 000 & 10 400 \\ \hline
Croissance PIB réel (\%) & 0,6 & 2,9 & 3,8 & 3,2 \\ \hline
Inflation (\%) & 6,0 & 24,5 & 33,9 & 6,1 \\ \hline
Dette publique / PIB (\%) & 75 & 38 & 96 & 69 \\ \hline
IDE entrants (Md USD) & 5,2 & 3,0 & 10,0 & 3,5 \\ \hline
\end{tabularx}
\end{center}
\textbf{Consignes :}
\begin{enumerate}
    \item Présentez le document (nature, source, date, thème).
    \item Identifiez les deux pays dont la croissance est la plus forte en 2023 et proposez deux explications pour chacun.
    \item Montrez, en vous appuyant sur le tableau, que la performance macroéconomique ne garantit pas la stabilité sociale. Citez deux indicateurs pertinents.
    \item En quoi la structure de la dette (indicateurs « Dette publique/PIB » et « Inflation ») différencie-t-elle les vulnérabilités de l'Égypte et du Nigéria ?
\end{enumerate}
\end{exercice}

\begin{corrige}[Exercice n°14-4]
\textbf{1. Présentation :} il s'agit d'un tableau statistique extrait de l'\emph{African Economic Outlook 2024} (Banque africaine de développement) et du FMI, comparant six indicateurs macroéconomiques clés pour les quatre foyers émergents africains en 2023.

\textbf{2. Croissance la plus forte :} \textbf{Égypte (3,8\,\%)} et \textbf{Maroc (3,2\,\%)}.
\begin{itemize}
    \item \emph{Égypte :} reprise du tourisme (recettes record), production de gaz (champ de Zohr), investissements dans la SCZone, soutien financier du Golfe et du FMI.
    \item \emph{Maroc :} performances de l'automobile (exportations record), résilience du tourisme, investissements publics (notamment reconstruction après le séisme d'Al Haouz de septembre 2023), malgré une année agricole affectée par la sécheresse.
\end{itemize}

\textbf{3. Performance $\neq$ stabilité sociale :}
\begin{itemize}
    \item \textbf{Nigéria :} croissance de 2,9\,\% mais \textbf{inflation de 24,5\,\%} (érosion du pouvoir d'achat, suppression des subventions carburant, dévaluation du naira) et PIB/habitant en PPA le plus faible des quatre (6 100 USD) : la richesse nationale ne ruisselle pas.
    \item \textbf{Égypte :} croissance de 3,8\,\% mais \textbf{inflation de 33,9\,\%} (crise de change, dévaluations successives) et \textbf{dette de 96\,\% du PIB} (risque de défaut, service de la dette très lourd).
    \item \textbf{Afrique du Sud :} croissance quasi nulle (0,6\,\%) avec un chômage supérieur à 30\,\% et des délestages : tensions sociales majeures.
\end{itemize}

\textbf{4. Vulnérabilités dette/inflation :}
\begin{itemize}
    \item \textbf{Égypte :} \textbf{dette très élevée (96\,\% du PIB)}, dont une grande part en devises (eurobonds), et \textbf{inflation très forte (34\,\%)}. Risque : spirale dette-inflation-dévaluation, dépendance au FMI et aux pays du Golfe.
    \item \textbf{Nigéria :} \textbf{dette modérée (38\,\% du PIB)} mais \textbf{inflation forte (24,5\,\%)}. Risque : le service de la dette absorbe l'essentiel des recettes fédérales (assiette fiscale très faible, de l'ordre de 7 à 10\,\% du PIB), l'insécurité freine les recettes pétrolières, pression sur le naira. Vulnérabilité plus fiscale et structurelle que de solvabilité brute.
\end{itemize}
\end{corrige}

\coursec{Les foyers économiques en expansion}

\courssub{Définir et localiser les foyers économiques en expansion}

\begin{definition}[Pays émergent]
Un \cle{pays émergent} est un pays en développement qui connaît une \textbf{croissance économique soutenue et rapide} (généralement $> 4\%$ par an sur une longue période), une \textbf{industrialisation accélérée}, une \textbf{insertion croissante dans les échanges mondiaux} (exportations manufacturées et services) et qui voit naître ses propres \cle{firmes transnationales émergentes} (FTN du Sud). Il se distingue des NPI (Nouveaux Pays Industrialisés) par sa taille démographique, son poids géopolitique récent et un développement encore inachevé (IDH moyen, fortes inégalités).
\end{definition}

\begin{propriete}[Critères de repérage d'un foyer émergent]
Pour identifier un foyer émergent sur une carte ou dans un tableau, on croise quatre indicateurs :
\begin{enumerate}
\item \textbf{Taux de croissance du PIB réel} nettement supérieur à la moyenne mondiale (souvent $4\text{--}8\%$).
\item \textbf{Part de l'industrie et des services dans le PIB} en hausse, part de l'agriculture en baisse.
\item \textbf{Taux d'ouverture} (Exportations + Importations / PIB) élevé et croissant.
\item \textbf{Présence de FTN du Sud} (siège social dans le pays) et d'infrastructures de rang mondial (ports, aéroports, autoroutes, câbles sous-marins).
\end{enumerate}
\end{propriete}

\begin{popfigure}
\includegraphics[width=\linewidth,height=8cm,keepaspectratio]{N14/nic.png}
\legende{Les pays nouvellement industrialisés (NPI) en rouge : Chine, Inde, Brésil, Mexique, Afrique du Sud, Turquie, Thaïlande, Malaisie, Indonésie et Philippines (liste de 2013). La carte localise les principaux foyers d'expansion des années 2000-2020 ; elle ne représente ni le poids économique de chaque pays ni les flux Sud-Sud, que le tableau et le texte de la leçon permettent d'ajouter. \\ Source : Wikimedia Commons, RidwanFadilArif, CC BY-SA 3.0.}
\end{popfigure}

\begin{experience}[TP Guidé : Construire un croquis des foyers émergents]
\textbf{Objectif :} Réaliser un croquis planisphérique des foyers émergents à l'échelle mondiale.
\textbf{Matériel :} Fond de carte monde A4 (projection Robinson ou Mercator simplifiée), crayons de couleur, règle, calculatrice.
\textbf{Protocole :}
\begin{enumerate}
\item \textbf{Choix des pays :} Sélectionner 10 pays émergents représentatifs (Brésil, Mexique, Argentine, Afrique du Sud, Nigéria, Égypte, Maroc, Inde, Chine, Indonésie, Turquie, Arabie Saoudite, Émirats Arabes Unis, Malaisie, Vietnam).
\item \textbf{Données :} Recueillir le PIB nominal (FMI/Banque mondiale, année 2022 ou 2023) et le taux de croissance moyen récent.
\item \textbf{Sémiologie :} 
    \begin{itemize}
    \item \textbf{Surface du cercle} (ou carré) proportionnelle au PIB (échelle : $\SI{1}{cm^2}$ pour $\SI{1000}{milliards}\ \$US$).
    \item \textbf{Couleur du contour} par continent (Bleu Amériques, Orange Asie, Vert Afrique, Violet Europe/Moyen-Orient).
    \item \textbf{Hachure intérieure} si taux de croissance $> 5\%$ (dynamique forte).
    \end{itemize}
\item \textbf{Flux :} Tracer 3 flèches majeures Sud-Sud (ex: Chine $\to$ Afrique matières 1ères ; Brésil $\to$ Chine soja/fer ; Inde $\to$ Afrique pharmacie/services).
\item \textbf{Légende :} Titre, Auteur, Date, Échelle des cercles, Signification des couleurs/hachures, Flux.
\end{enumerate}
\textbf{Attendu :} Un croquis lisible, légendé, montrant la concentration des émergents en Asie et la percée africaine/latino-américaine.
\end{experience}

\courssub{Les pays émergents d'Amérique : Brésil, Mexique, Argentine}

\begin{itemize}
\item \textbf{Brésil :} 1ère puissance d'Amérique latine (PIB $\approx \SI{1900}{milliards}\ \$US$, 2022). \cle{Moteurs} : Agrobusiness (1er exportateur mondial de soja, café, sucre, viande), Mines (fer Vale, 2ème producteur mondial), Industrie aéronautique (\textbf{Embraer}, 3ème constructeur mondial), Pétrole (\textbf{Petrobras}, pré-sal). \cle{Atouts} : Taille continentale, ressources, marché intérieur de 215 M hab. \cle{Défis} : Inégalités extrêmes (Gini $\approx 0,53$), déforestation Amazonie, instabilité politique récurrente, infrastructure logistique insuffisante.
\item \textbf{Mexique :} 2ème puissance régionale (PIB $\approx \SI{1400}{milliards}\ \$US$). \cle{Moteurs} : \textbf{Maquiladoras} (usines d'assemblage frontalières), Industrie automobile (1er exportateur vers USA, 7ème producteur mondial), Électronique, Tourisme. \cle{Atout majeur} : Accès privilégié au marché US (Accord \textbf{ACEUM} / ex-ALÉNA), main-d'œuvre qualifiée et moins chère qu'aux USA. \cle{Défis} : Dépendance excessive aux USA ($> 80\%$ exportations), insécurité (cartels), pauvreté rurale, informel.
\item \textbf{Argentine :} PIB $\approx \SI{630}{milliards}\ \$US$. \cle{Moteurs} : Agriculture (soja, maïs, blé, viande), Énergie (Vaca Muerta : 2ème réserves mondiales gaz de schiste, 4ème pétrole), Industrie agroalimentaire, Tech (licornes : MercadoLibre, Globant). \cle{Freins} : Instabilité macroéconomique chronique (inflation $> 100\%$, dette, défauts souverains répétés), contrôles de change, fuite des capitaux.
\end{itemize}

\begin{savaistu}[L'Argentine et le « mal hollandais »]
L'Argentine illustre le piège de la \cle{maladie hollandaise} : la rente agricole et gazière fait monter la monnaie, tue l'industrie manufacturière non compétitive, et rend l'économie vulnérable aux cours mondiaux. La diversification y est plus difficile qu'au Brésil ou au Mexique.
\end{savaistu}

\courssub{Les pays émergents d'Asie et du Moyen-Orient}

\begin{itemize}
\item \textbf{Les « Quatre Dragons » (Corée du Sud, Taïwan, Hong Kong, Singapour) :} \cle{NPI de 1ère génération} (années 1960-1990). Aujourd'hui \textbf{pays développés} (revenu élevé, IDH très haut). Modèle : industrialisation par substitution aux importations puis exportation, éducation massive, État stratège. \textbf{Samsung, TSMC, Hyundai} sont nées là.
\item \textbf{Malaisie, Thaïlande, Indonésie, Vietnam :} \cle{NPI de 2ème génération} / Émergents ASEAN. \textbf{Malaisie} : Électronique (puces), Pétrole/Gaz (Petronas), Halal hub. \textbf{Vietnam} : « Usine du monde » délocalisée de Chine (textile, électronique Samsung/Intel), croissance $6\text{--}7\%$.
\item \textbf{Inde :} 1ère population mondiale (1,43 Md, 2023), 5ème PIB nominal ($\approx \SI{3400}{milliards}\ \$US$). \cle{Moteurs} : Services IT (\textbf{TCS, Infosys, Wipro} : 1er exportateur mondial services informatiques), Pharmacie générique (« pharmacie du monde »), Industrie (automobile \textbf{Tata}, \textbf{Mahindra}), Démarrage « Make in India ». \cle{Atouts} : Dividende démographique (âge médian 28 ans), diaspora qualifiée, anglais. \cle{Défis} : Infrastructures déficientes, bureaucratie, inégalités de genre/castes, emploi informel ($> 80\%$).
\item \textbf{Chine :} Cas à part (2ème PIB mondial $\approx \SI{18000}{milliards}\ \$US$). « Atelier du monde » $\to$ montée en gamme (Haute tech, IA, VE, Nucléaire). \textbf{FTN} : Huawei, BYD, CATL, ICBC. Moteur principal de la demande mondiale de matières premières.
\item \textbf{Pays du Golfe (Arabie Saoudite, Émirats Arabes Unis, Qatar) :} \cle{Rente pétrolière/gazière} réinvestie dans diversification (\textbf{Vision 2030} Arabie : NEOM, tourisme, finance ; \textbf{Dubaï} : Hub logistique/aérien/finance mondial, Port Jebel Ali, Emirates Airlines ; \textbf{Qatar} : GNL, Finance). Main-d'œuvre migrante massive (système \emph{Kafala} en réforme).
\end{itemize}

\begin{aretenir}[NPI vs Pays Émergents : la distinction Bac]
\begin{itemize}
\item \textbf{NPI (Nouveaux Pays Industrialisés) :} Industrialisation achevée, revenu élevé, IDH très élevé. Ex : 4 Dragons (Corée du Sud, Taïwan, Singapour, Hong Kong). \textbf{Ne plus les classer comme « émergents » au Bac.}
\item \textbf{Pays Émergents (au sens strict actuel) :} Industrialisation en cours, revenu intermédiaire (transsupérieur ou infrasupérieur), IDH moyen/élevé, forte croissance, poids géopolitique nouveau. Ex : \textbf{BRICS} (Brésil, Russie, Inde, Chine, Afrique du Sud + nouveaux membres 2024 : Égypte, Émirats, Éthiopie, Iran, Arabie Saoudite), Mexique, Indonésie, Turquie, Vietnam, Nigeria.
\end{itemize}
\end{aretenir}

\courssub{Les pays émergents d'Afrique : Afrique du Sud, Nigéria, Égypte, Maroc}

\begin{itemize}
\item \textbf{Afrique du Sud :} 1ère puissance industrielle du continent (PIB $\approx \SI{400}{milliards}\ \$US$). \cle{Moteurs} : Mines (Or, Platine, Charbon, Diamants - \textbf{Anglo American}, \textbf{Sibanye}), Finance (\textbf{Standard Bank}, \textbf{Sanlam}), Industrie (Auto BMW, Mercedes, VW ; Acier \textbf{ArcelorMittal}), Retail (\textbf{Shoprite}), Télécoms (\textbf{MTN}, \textbf{Vodacom} - leaders au Cameroun). \cle{Statut} : Seule africaine dans \textbf{G20} et \textbf{BRICS}. \cle{Crise} : Chômage $\approx 32\%$, coupures d'électricité (\emph{load shedding} Eskom), criminalité, inégalités post-apartheid (Gini $\approx 0,63$).
\item \textbf{Nigéria :} 1ère population Afrique (220 M), 1er PIB Afrique ($\approx \SI{470}{milliards}\ \$US$, 2022). \cle{Moteurs} : Pétrole/Gaz ($> 90\%$ exportations, $> 50\%$ recettes État), Agriculture (Cacao, Caoutchouc, Ignames), FinTech (\textbf{Flutterwave}, \textbf{Opay} - licornes), Cinéma (\textbf{Nollywood} 2ème mondial), Musique (Afrobeats). \cle{Défis} : Dépendance pétrole, insécurité (Boko Haram, banditisme), corruption, inflation, infrastructures déficientes, électricité ($< 5000$ MW pour 220 M hab.).
\item \textbf{Égypte :} 110 M hab., PIB $\approx \SI{470}{milliards}\ \$US$. \cle{Moteurs} : Canal de Suez (revenus records $\SI{9,4}{milliards}\ \$US$ 2023), Gaz naturel (champ Zohr), Tourisme, Agriculture (deltas Nil), Industrie textile/agro, Zone économique de Suez (SCZone). \cle{Rôle} : Pivot géopolitique Méditerranée/Mer Rouge, membre BRICS 2024.
\item \textbf{Maroc :} PIB $\approx \SI{130}{milliards}\ \$US$. \cle{Moteurs} : Phosphates (\textbf{OCP} : 70\% réserves mondiales, 1er exportateur engrais), Automobile (\textbf{Renault}, \textbf{Stellantis}, \textbf{Citroën} : 1er exportateur Afrique vers UE), Aéronautique (\textbf{Bombardier}, \textbf{Safran}, \textbf{Boeing} - écosystème Tanger/Casablanca), Énergies renouvelables (Noor Ouarzazate), Tourisme, Finance islamique. \cle{Atout} : Stabilité politique, proximité UE, accords de libre-échange (UE, USA, Turquie, Afrique via ZLECAf).
\end{itemize}

\begin{exemplebox}[Le Cameroun face aux émergents africains]
Le Cameroun importe du ciment \textbf{Dangote} (Nigéria), de l'engrais \textbf{OCP} (Maroc), utilise le réseau \textbf{MTN} (Afrique du Sud), voit des banques \textbf{Attijariwafa} / \textbf{BCP} (Maroc) et \textbf{Standard Bank} (Afrique du Sud) s'installer. Le port de Kribi est construit par \textbf{CMEC} (Chine) et géré par \textbf{Dubai Port World} (Émirats). L'émergence africaine n'est pas abstraite : elle structure l'économie camerounaise quotidienne.
\end{exemplebox}

\courssub{Facteurs, manifestations et limites de l'émergence}

Nous venons de parcourir les foyers émergents d'Amérique, d'Asie, du Moyen-Orient et d'Afrique. Une question s'impose : comment des pays longtemps classés parmi les « pays en développement » sont-ils devenus des acteurs incontournables ? Et pourquoi cette ascension reste-t-elle fragile ?

\courssub{Les facteurs de l'émergence : une combinaison de causes internes et externes}

\textbf{Les facteurs internes : les atouts propres à chaque pays}\par

L'émergence repose d'abord sur des atouts que les pays mobilisent eux-mêmes.

\begin{itemize}
\item \textbf{Les ressources naturelles.} Beaucoup de pays émergents disposent de matières premières recherchées : pétrole et gaz (Nigéria, Arabie Saoudite, Émirats Arabes Unis, Mexique), minerais (fer au Brésil \textbf{Vale}, or/platine en Afrique du Sud), produits agricoles (soja brésilien/argentin, cacao ivoirien/ghanéen/nigérian). Ces ressources, exportées, rapportent des devises qui financent l'investissement.
\item \textbf{Une main-d'œuvre nombreuse et compétitive.} L'Inde (1,43 Md), l'Indonésie (275 M), le Brésil, le Mexique, le Nigéria offrent une main-d'œuvre abondante, jeune (dividende démographique) et à coût unitaire faible. C'est l'atout qui a permis à l'Asie de devenir « l'atelier du monde ».
\item \textbf{Les réformes structurelles d'ouverture.} Années 1980-1990 : libéralisation des échanges, \cle{privatisations}, accueil des IDE, assainissement budgétaire (Plans d'ajustement structurel FMI/BM puis réformes autonomes). Ex : \textbf{Inde 1991} (fin \emph{Licence Raj}), \textbf{Brésil 1990} (Plan Collor/Real), \textbf{Vietnam 1986} (\emph{Doi Moi}), \textbf{Maroc 1983/2000}.
\item \textbf{Stabilité politique et institutionnelle relative.} Fin des dictatures militaires (Brésil 1985, Argentine 1983, Chili 1990), fin de l'apartheid (Afrique du Sud 1994), monarchies stables (Maroc, Golfe). La prévisibilité juridique attire le capital.
\item \textbf{Investissement dans le capital humain et les infrastructures.} Dragons : scolarisation massive précoce $\to$ ingénieurs. Chine/Inde : millions de diplômés STEM/an. Infrastructures : Barrage des Trois Gorges (Chine), Autoroutes Brésil/Mexique, Port Tanger Med (Maroc), Ligne à grande vitesse Maroc, Métro Lagos/Abidjan, Câbles sous-marins (2Africa, METISS).
\end{itemize}

\textbf{Les facteurs externes : l'impulsion venue de l'extérieur}\par

Aucun pays n'émerge seul. L'environnement international joue un rôle décisif.

\begin{itemize}
\item \textbf{Les investissements directs étrangers (IDE) des firmes transnationales (FTN) du Nord.} Toyota, Volkswagen, Samsung, Nestlé, TotalEnergies, Apple (sous-traitants Foxconn/Pegatron) installent usines/R\&D pour bas salaires + marchés locaux. Apport : capitaux, technologies, management, accès réseaux mondiaux.
\item \textbf{Les délocalisations et fragmentation des chaînes de valeur mondiales (CVM).} Déplacement production standardisée (textile $\to$ Bangladesh/Vietnam/Éthiopie ; électronique $\to$ Malaisie/Vietnam/Mexique ; auto $\to$ Mexique/Brésil/Slovaquie/Maroc). « \cle{Factory Asia} », « \cle{Factory North America} » (Mexique), « \cle{Factory Europe} » (Maroc/Europe de l'Est/Turquie).
\item \textbf{Le « Super-cycle » des matières premières (années 2000-2014).} Hausse structurelle prix (pétrole, fer, cuivre, soja) tirée par urbanisation chinoise. Enrichissement exportateurs (Brésil, Chili, Pérou, Afrique australe, Golfe, Russie). \cle{Effet de manne} : réserves de change, fonds souverains (Norvège, Arabie Saoudite \textbf{PIF}, Émirats \textbf{ADIA}, Chine \textbf{CIC}).
\item \textbf{La demande chinoise comme locomotive.} Chine = 1er partenaire commercial de > 120 pays. Appétit insatiable : soja/fer Brésil, cuivre Chili/Zambie/RCongo, pétrole Angola/Arabie/Iran/Russie, GNL Qatar/Australie/USA. Création d'interdépendance forte (risque de « piège à dette » via BRI / Nouvelles Routes de la Soie).
\end{itemize}

\begin{attention}[Erreur fréquente : la cause unique]
Au Baccalauréat, beaucoup de copies attribuent l'émergence à \emph{une seule} cause (« le Brésil émerge grâce au pétrole », « l'Inde grâce à sa population »). C'est une erreur d'analyse sanctionnée. L'émergence résulte \textbf{toujours d'une combinaison} de facteurs internes (atouts nationaux, réformes, éducation) et de facteurs externes (IDE, CVM, conjoncture matières premières, demande chinoise). Dans un commentaire ou une dissertation, organise ton raisonnement en distinguant explicitement ces deux registres de causes.
\end{attention}

\begin{popfigure}
\begin{center}
\begin{tikzpicture}[font=\small, >=Stealth, node distance=0.5cm]
% Facteurs Internes
\node[draw=popblue, fill=popblueL, rounded corners, text width=4.5cm, align=center] (internes) at (0,2.5) {\textbf{Facteurs Internes (Offre nationale)}\\[3pt] Ressources naturelles (Énergie, Mines, Agro)\\ Dividende démographique (Main-d'œuvre jeune)\\ Réformes libérales, Stabilisation macro\\ Investissement Éducation (STEM) \& Infrastructures\\ État stratège (Planification industrielle)};
% Facteurs Externes
\node[draw=poporange, fill=poporangeL, rounded corners, text width=4.5cm, align=center] (externes) at (0,-2.5) {\textbf{Facteurs Externes (Demande/Conjoncture)}\\[3pt] IDE des FTN du Nord (Capitaux, Tech, Réseaux)\\ Délocalisations / Fragmentation CVM\\ Super-cycle matières premières (2000-2014)\\ Demande chinoise massive (Moteur mondial)\\ Intégration régionale (ASEAN, Mercosur, ZLECAf)};
% Croissance
\node[draw=poppurple, fill=poppurpleL, rounded corners, text width=3.5cm, align=center] (croissance) at (7,0) {\textbf{Croissance économique forte}\\[3pt] PIB $\uparrow$ 4-8\%/an\\ Industrialisation / Tertiarisation\\ Accumulation capital physique/humain\\ Excédents commerciaux / Réserves};
% Insertion
\node[draw=popgreen, fill=popgreenL, rounded corners, text width=3.8cm, align=center] (insertion) at (12.5,0) {\textbf{Insertion puissance dans la mondialisation}\\[3pt] Exportations diversifiées (Haute tech)\\
Naissance FTN du Sud (Petrobras, Tata, Huawei, MTN, OCP)\\
Métropoles mondiales (São Paulo, Dubaï, Shanghai, Mumbai)\\
Poids géopolitique : G20, BRICS, NBD, AIIB};
\draw[->, very thick, popblue] (internes.east) -- node[above, sloped, pos=0.5, font=\tiny] {Mobilisation atouts} (croissance.north west);
\draw[->, very thick, poporange] (externes.east) -- node[below, sloped, pos=0.5, font=\tiny] {Environnement favorable} (croissance.south west);
\draw[->, very thick, poppurple] (croissance) -- node[above, font=\tiny] {Richesse \& Capacités} (insertion);
\end{tikzpicture}
\end{center}
\legende{Schéma explicatif : la convergence des facteurs internes (offre) et externes (demande/conjoncture) alimente la croissance, qui permet l'insertion par le haut dans la mondialisation (FTN du Sud, gouvernance mondiale).}
\end{popfigure}

\courssub{Les manifestations de l'émergence}

Trois types de signes trahissent concrètement l'émergence.

\textbf{Des manifestations économiques : croissance, exportations et firmes transnationales du Sud}\par

Les pays émergents affichent des taux de croissance du \cle{PIB} nettement supérieurs à ceux des pays riches (souvent $4\text{--}8\%$ vs $1\text{--}2\%$ en zone OCDE), et leurs exportations progressent plus vite que le commerce mondial. Structure exportations : montée en gamme (Brésil : avions Embraer, Mexique : auto, Inde : pharma/IT, Maroc : auto/aéro, Vietnam : électronique).

Fait nouveau et capital : ces pays ne se contentent plus d'accueillir les FTN du Nord, ils font naître leurs propres géants.

\begin{definition}[Firme transnationale émergente (FTN du Sud)]
Une \cle{firme transnationale émergente} est une entreprise originaire d'un pays émergent qui possède des filiales et des activités de production/vente/R\&D dans plusieurs pays. Elles investissent massivement au Nord (rachats) et au Sud (infrastructures, ressources).
\begin{itemize}
\item \textbf{Amérique :} \textbf{Petrobras} (Pétrole, Brésil), \textbf{Embraer} (Aéronautique), \textbf{Ambev} (Boissons), \textbf{Natura} (Cosmétiques).
\item \textbf{Asie :} \textbf{Tata} (Acier, Auto, IT, Inde), \textbf{Mahindra} (Auto/Tracteurs), \textbf{Infosys/TCS/Wipro} (Services IT), \textbf{Huawei/BYD/CATL/ICBC} (Chine), \textbf{Samsung/Hyundai} (Corée - ex-NPI), \textbf{Petronas} (Malaisie).
\item \textbf{Moyen-Orient :} \textbf{Emirates/Etihad/Qatar Airways} (Transport aérien), \textbf{DP World} (Ports/Logistique - gère Kribi Cameroun), \textbf{Aramco} (Pétrole), \textbf{PIF/ADIA/Mubadala} (Fonds souverains investisseurs mondiaux).
\item \textbf{Afrique :} \textbf{MTN Group} (Télécoms, Af. Sud - leader Cameroun), \textbf{Dangote Group} (Ciment/Sucre/Raffinage, Nigéria), \textbf{OCP Group} (Engrais/Phosphates, Maroc), \textbf{Shoprite} (Retail, Af. Sud), \textbf{Standard Bank} (Banque, Af. Sud), \textbf{Attijariwafa Bank / BCP} (Banques, Maroc - leaders UEMOA/CEMAC).
\end{itemize}
\end{definition}

\begin{exemplebox}[Le poids des BRICS+ en chiffres (2023-2024, ordres de grandeur FMI/Banque Mondiale)]
Les \cle{BRICS} (Brésil, Russie, Inde, Chine, Afrique du Sud + Égypte, Émirats, Éthiopie, Iran, Arabie Saoudite depuis 2024) regroupent :
\begin{itemize}
\item $\approx$ \textbf{45 \% de la population mondiale} ($\approx 3,6$ Md).
\item $\approx$ \textbf{35-37 \% du PIB mondial} (en Parité de Pouvoir d'Achat - PPA) ; $\approx 28-30\%$ en nominal \$US.
\item $\approx$ \textbf{25-30 \% des exportations mondiales} de marchandises.
\item $\approx$ \textbf{40-45 \% de la production pétrolière} (avec nouveaux membres).
\end{itemize}
Ces chiffres confirment le déplacement du centre de gravité économique vers le Sud et l'Est.
\end{exemplebox}

\textbf{Des manifestations spatiales : la métropolisation et les corridors}\par

L'émergence se lit dans les paysages : de gigantesques \cle{métropoles mondiales} se dressent, vitrines de la réussite, nœuds de commandement et de flux.
\begin{itemize}
\item \textbf{Amérique :} \textbf{São Paulo} (22 M hab. aire urbaine) : cœur financier/industriel Brésil ; \textbf{Mexico} (22 M) : commande économie nationale ; \textbf{Buenos Aires} (15 M).
\item \textbf{Asie :} \textbf{Shanghai} (29 M), \textbf{Pékin} (22 M), \textbf{Shenzhen} (18 M, ex-village pêcheurs $\to$ Silicon Valley chinoise), \textbf{Mumbai} (21 M), \textbf{Délhi} (33 M), \textbf{Jakarta} (35 M), \textbf{Manille} (28 M).
\item \textbf{Moyen-Orient :} \textbf{Dubaï} (3,5 M mais hub mondial : Port Jebel Ali 1er port artificiel, Aéroport DXB 1er trafic international, Burj Khalifa, Finance, Tourisme) ; \textbf{Riyad} (projet NEOM, Vision 2030).
\item \textbf{Afrique :} \textbf{Johannesburg/Pretoria} (15 M, cœur financier/industriel), \textbf{Lagos} (21 M, capitale économique Nigéria, FinTech, Port, Nollywood), \textbf{Le Caire} (22 M), \textbf{Casablanca/Rabat} (12 M, hub finance/auto/aéro/logistique Maroc), \textbf{Nairobi} (5 M, hub Est Afrique/Tech « Silicon Savannah »).
\end{itemize}
Ces métropoles sont reliées par des \cle{corridors de développement} (ex: Corridor Maputo - Gauteng, Corridor Lagos-Abidjan, Corridor Tanger-Casablanca-Kénitra, Corridor Pékin-Shanghai-Guangzhou).

\textbf{Des manifestations politiques : un nouveau poids dans la gouvernance mondiale}\par

Les pays émergents réclament et obtiennent une place dans les instances qui dirigent l'économie mondiale :
\begin{itemize}
\item Le \cle{G20} (créé 1999, sommet chefs d'État 2008) : réunit 19 pays + UE + UA (Union Africaine 2023). Inclut Chine, Inde, Brésil, Mexique, Argentine, Arabie Saoudite, Afrique du Sud, Indonésie, Turquie, Corée du Sud. A supplanté le G7/G8.
\item La \textbf{réforme des quotes-parts FMI/Banque Mondiale} (2010, appliquée 2016) : revalorisation droits de vote Chine (6\% $\to$ 12\%), Inde, Brésil, Russie. L'Europe cède des sièges au CA.
\item Les \textbf{BRICS} créent leurs propres institutions : \cle{Nouvelle Banque de Développement (NBD)} (créée 2015, siège \textbf{Shanghai}, capital $\SI{100}{milliards}\ \$US$, présidence tournante - 1ère Dilma Rousseff Brésil), \cle{Accord de réserve de contingences (ARC)} (fonds de réserve monétaire $\SI{100}{milliards}\ \$US$), système de paiement \textbf{BRICS Pay} (alternative SWIFT), think-tanks, agence de notation.
\end{itemize}

\begin{savaistu}[La Nouvelle Banque de Développement (NBD) : une alternative au Sud]
La NBD finance des projets d'infrastructures (énergie propre, transport, eau, assainissement, numérique) dans les pays membres et autres pays en développement, \textbf{sans conditionnalité politique} (gouvernance, droits de l'homme) contrairement à la BM/FMI. Elle emprunte sur les marchés en devises locales (yuan, roupie, rand, real) pour réduire le risque de change. Preuve que les émergents veulent \textbf{écrire les règles} de la mondialisation, pas seulement y participer.
\end{savaistu}

\begin{popfigure}
\begin{center}
\begin{tikzpicture}[font=\small, grow=right, level distance=3.8cm,
level 1/.style={sibling distance=3.0cm},
level 2/.style={sibling distance=1.3cm, level distance=4.2cm},
edge from parent/.style={draw=popmuted, thick, ->}]
\node[draw=popdark, fill=popcream, rounded corners, align=center, text width=2.8cm] {\textbf{L'Émergence :\\Synthèse}}
child { node[draw=popgreen, fill=popgreenL, rounded corners, align=center, text width=2.5cm] {\textbf{Manifestations}}
  child { node[draw=popgreen, rounded corners, align=center, text width=3.5cm] {\textbf{Économiques}\\ PIB $\uparrow$, Export $\uparrow$, FTN Sud (Petrobras, Tata, MTN, OCP, Huawei)} }
  child { node[draw=popgreen, rounded corners, align=center, text width=3.5cm] {\textbf{Spatiales}\\ Métropoles mondiales (São Paulo, Dubaï, Shanghai, Lagos, Casablanca)\\ Corridors logistiques} }
  child { node[draw=popgreen, rounded corners, align=center, text width=3.5cm] {\textbf{Politiques}\\ G20, BRICS+, NBD, AIIB\\ Réforme FMI/BM\\ Diplomatie économique (BRI)} }
}
child { node[draw=popink, fill=poppinkL, rounded corners, align=center, text width=2.5cm] {\textbf{Limites \& Fragilités}}
  child { node[draw=popink, rounded corners, align=center, text width=3.5cm] {\textbf{Sociales \& Spatiales}\\ Inégalités extrêmes (Gini $> 0,5$)\\ Pauvreté persistante (Inde, Nigéria)\\ Fracture Métropoles / Périphéries\\ Chômage jeunes (Af. Sud 60\%, MENA 25\%)} }
  child { node[draw=popink, rounded corners, align=center, text width=3.5cm] {\textbf{Économiques}\\ Dépendance matières 1ères (Pétrole, Soja, Minerais)\\ Dépendance IDE / Technologie cœur (Semi-conducteurs)\\ Endettement \$US / Fuite capitaux (Sudden Stop)\\ Piège revenu intermédiaire} }
  child { node[draw=popink, rounded corners, align=center, text width=3.5cm] {\textbf{Environnementales}\\ Déforestation (Amazonie, Bornéo, Bassin Congo)\\ Pollution air/eau/sols (Mégapoles, Delta Niger)\\ Émissions GES en forte hausse (Chine 1er, Inde 3ème)\\ Stress hydrique (Moyen-Orient, Afrique du Nord, Inde)} }
  child { node[draw=popink, rounded corners, align=center, text width=3.5cm] {\textbf{Géopolitiques / Systémiques}\\ Vulnérabilité crises mondiales (2008, 2020, Guerre Ukraine)\\ Tensions commerciales (USA/Chine)\\ Instabilité politique interne (Brésil, Afrique du Sud, Moyen-Orient)\\ Rivalités régionales (Inde/Chine, Iran/Arabie, Éthiopie/Égypte)} }
};
\end{tikzpicture}
\end{center}
\legende{Arbre synthétique complet : Manifestations (Éco, Spatiales, Politiques) et Limites (Sociales, Éco, Enviro, Géopolitiques) de l'émergence. Structure type pour une dissertation ou un commentaire de document.}
\end{popfigure}

\courssub{Les limites de l'émergence : une ascension inachevée}

L'émergence ne doit pas être confondue avec le développement achevé (pays développés). Quatre limites majeures la fragilisent, formant souvent un \cle{piège à revenu intermédiaire}.

\begin{itemize}
\item \textbf{Les inégalités sociales et spatiales structurelles.} La richesse créée est très mal répartie (coefficients de Gini souvent $> 0,50$). \textbf{Brésil} : Favelas (Rio, São Paulo) vs quartiers d'affaires ; Nord-Est pauvre vs Sud dynamique. \textbf{Afrique du Sud} : Chômage $\approx 32\%$ (jeunes $> 60\%$), townships vs Sandton. \textbf{Inde} : Des centaines de millions dans la pauvreté multidimensionnelle vs milliardaires (Ambani, Adani). \textbf{Nigéria} : 40\% sous seuil pauvreté ($<\SI{2,15}{jour}$) vs élites pétrolières. \textbf{Chine} : Hukou (régime résidence) crée 290 M migrants ruraux sans droits urbains pleins.
\item \textbf{La dépendance économique structurelle (spécialisation régressive).} Beaucoup d'émergents restent coincés en bas des chaînes de valeur : exportation matières premières brutes (pétrole Nigéria/Angola/Golfe, minerais Zambie/RCongo/Chili, soja Brésil/Argentine/Paraguay) et importation biens d'équipement/haute tech. Bénéfices IDE rapatriés vers sièges Nord. Faible valeur ajoutée locale. \cle{Piège du revenu intermédiaire} : incapacité à passer de l'imitation/assemblage à l'innovation de rupture (sauf Corée/Taïwan/Chine partielle).
\item \textbf{Les dégradations environnementales majeures (coût externe de la croissance).} \textbf{Déforestation} : Amazonie (soja/élevage), Bornéo/Sumatra (palmier à huile), Bassin du Congo (bois/mines). \textbf{Pollution} : Air (New Delhi, Pékin, Lahore, Le Caire, Mexico - PM2.5 records), Eau (Citarum Indonésie, Yamuna Inde, Delta Niger marées noires), Sols (métaux lourds). \textbf{Émissions GES} : Chine (30\% mondial), Inde (7\%), Russie, Iran, Indonésie, Arabie Saoudite, Brésil (déforestation). Conflit \textbf{Croissance vs Climat} (Accord Paris : responsabilités communes mais différenciées).
\item \textbf{La vulnérabilité aux chocs extérieurs (interdépendance asymétrique).} Crise financière 2008 (chute IDE/export), Chute cours pétrole 2014-2016 (récession Brésil, Russie, Nigéria, Golfe), Pandémie 2020 (rupture CVM, dette), Guerre Ukraine 2022 (inflation alimentaire/énergétique, sanctions Russie). \cle{Fuite des capitaux} (\"Sudden Stop\") : investisseurs de portefeuille (« hot money ») partent au 1er signe de risque $\to$ dépréciation monnaie, crise balance paiements (Argentine, Turquie, Sri Lanka, Pakistan, Ghana, Zambie défauts 2020-2023).
\end{itemize}

\begin{attention}[Piège de rédaction majeur au Bac]
Ne présente \textbf{jamais} les pays émergents comme des « pays riches » ou « pays développés ». Le vocabulaire rigoureux attendu est :
\begin{itemize}
\item « Pays à \textbf{revenu intermédiaire} » (transsupérieur : Brésil, Chine, Russie, Mexique, Malaisie, Afrique du Sud ; infrasupérieur : Inde, Indonésie, Vietnam, Maroc, Égypte, Philippines).
\item « Pays en \textbf{forte croissance} » / « \textbf{Économies émergentes} ».
\item « Développement \textbf{inachevé} » / « \textbf{Développement inégal} » (IDH moyen/élevé, pas très élevé).
\end{itemize}
Écrire « Le Brésil est un pays développé » ou « La Chine est un pays riche » est une \textbf{erreur factuelle grave} sanctionnée par une note faible. Nuance : La Chine est une \textbf{superpuissance émergente} (2ème PIB, 1er exportateur, armée, nucléaire, spatial) mais revenu/habitant encore intermédiaire ($\approx \SI{12\,700}{\$US}$ nominal 2023).
\end{attention}

\courssub{Le revers de l'émergence : les pays au décollage économique difficile}

Tous les pays du Sud ne sont pas emportés par cette dynamique. À côté des émergents subsistent des pays \cle{au décollage économique difficile} : les \cle{PPTE} (Pays Pauvres Très Endettés), pour la plupart africains (Tchad, Centrafrique, Burundi, Soudan du Sud, Niger, Mali, Burkina Faso, Mozambique, Malawi, Madagascar, Haïti, Yémen, Afghanistan), et les pays à faible revenu (LIC).

\begin{propriete}[Caractéristiques des pays au décollage difficile (PPTE / LIC)]
\begin{itemize}
\item \textbf{Économie extravertie non diversifiée :} Dépendance 1 à 3 produits primaires (coton, or, uranium, café, cacao, pétrole brut) aux cours volatils.
\item \textbf{Fardeau de la dette :} Service de la dette $> 20-30\%$ recettes budgétaires $\to$ éviction dépenses santé/éducation/infrastructure. Initiative PPTE (1996) + PPTE renforcée (1999) + Initiative d'allègement dette G20 (DSSI 2020-2021) + Cadre commun (2020) : allègements partiels, mais reprise endettement rapide (Chine créancier bilatéral majeur, Eurobonds).
\item \textbf{Fragilité institutionnelle et sécuritaire :} Conflits armés, terrorisme (Sahel, Lac Tchad, Corne Afrique), coups d'État, gouvernance faible, corruption.
\item \textbf{Déficit infrastructurel massif :} Électricité ($< 20\%$ accès rural Sahel), Routes, Ports, Numérique, Santé, École.
\item \textbf{Marginalisation dans la mondialisation :} Part commerce mondial $< 1\%$ (Afrique subsaharienne hors Afrique du Sud/Nigéria). Captent peu d'IDE (hors secteur extractif). Pèsent peu dans instances mondiales (voix FMI/BM faibles).
\end{itemize}
\end{propriete}

\begin{exemplebox}[Le paradoxe nigérian : Émergent et PPTE à la fois]
Le Nigéria illustre la fracture interne au Sud : c'est la 1ère économie et 1ère population d'Afrique, membre du G20 (invité permanent), hub FinTech, Nollywood, Dangote (raffinerie $\SI{650\,000}{b/j}$). \textbf{MAIS} : 40\% pauvres, 20 M enfants non scolarisés, insécurité généralisée, infrastructure électrique effondrée, dette service $> 90\%$ recettes (2023). C'est un \textbf{pays émergent à deux vitesses} : un noyau moderne (Lagos, Abuja, Port Harcourt) et une immense périphérie exclue.
\end{exemplebox}

Cette fracture entre \textbf{Sud émergent} (intégré, dynamique, acteur) et \textbf{Sud exclu / en difficulté} (marginalisé, fragile, pourvoyeur de matières premières) fera l'objet de la leçon suivante : \emph{Les pays au décollage économique difficile} (Leçon 15).

\begin{methode}[Construire un diagramme en bâtons comparatif (PIB, Exportations, IDE, Population)]
Pour représenter, par exemple, le PIB nominal de plusieurs pays émergents :
\begin{enumerate}
\item \textbf{Choisir une échelle adaptée à la valeur maximale} : repère la plus grande valeur du tableau et fixe l'échelle pour que le bâton le plus haut tienne dans le cadre (ex. 1 cm pour $\SI{500}{milliards}\ \$US$).
\item \textbf{Ordonner les bâtons} par valeurs décroissantes (ou croissantes) : la lecture comparative devient immédiate.
\item \textbf{Tracer des bâtons de largeur égale}, régulièrement espacés, hachurés ou coloriés uniformément par catégorie (couleur continent).
\item \textbf{Indiquer obligatoirement} : Titre précis (variable, année, unité), Source (FMI, BM, ONU), Échelle graphique, Graduation axe des valeurs, Noms des pays sur axe des catégories.
\end{enumerate}
\end{methode}

\begin{exoresolu}[Calculer et comparer : la part des BRICS dans le PIB mondial]
Énoncé. En 2022 (données FMI, PIB nominal \$US courants, ordres de grandeur), le PIB mondial est d'environ $\SI{101\,000}{milliards}\ \$US$ et celui des 5 BRICS historiques d'environ $\SI{26\,500}{milliards}\ \$US$.
\begin{enumerate}
\item Calcule la part des 5 BRICS dans le PIB mondial nominal en 2022.
\item En 2000, le PIB mondial était d'environ $\SI{33\,500}{milliards}\ \$US$ et celui des 5 BRICS d'environ $\SI{3\,500}{milliards}\ \$US$. Calcule l'évolution de cette part entre 2000 et 2022 et commente en deux phrases.
\end{enumerate}
\tcblower
Solution détaillée.
\begin{enumerate}
\item Part des 5 BRICS en 2022 :
\[ \frac{26\,500}{101\,000}\times 100 \approx 26,2\,\% \]
Les 5 BRICS pèsent donc environ \textbf{un quart du PIB mondial nominal} (contre $\approx 35\%$ en PPA).
\item Part en 2000 :
\[ \frac{3\,500}{33\,500}\times 100 \approx 10,4\,\% \]
La part est donc passée d'environ \textbf{10 \% à 26 \% en vingt-deux ans}, soit une multiplication par 2,5.
\textbf{Commentaire :} Ce bond spectaculaire traduit l'insertion rapide des pays émergents dans l'économie mondiale, portée principalement par la croissance chinoise (passée de 3,6\% à 18\% du PIB mondial) et indienne. Il confirme le déplacement du centre de gravité économique vers l'Asie, même si les trois quarts de la richesse nominale restent produits par les pays avancés et la Chine seule.
\end{enumerate}
\end{exoresolu}

\begin{exercice}[À toi de jouer - Entraînement Bac]
\begin{enumerate}
\item \textbf{Cartographie (Croquis) :} Sur un fond de carte monde fourni (ou dessiné), localise et nomme : les 5 BRICS historiques, le Mexique, l'Indonésie, la Turquie, l'Arabie Saoudite, le Nigeria, le Maroc, le Vietnam. Représente par un cercle proportionnel (échelle : $\SI{1}{cm}$ de rayon pour $\SI{2000}{milliards}\ \$US$ PIB nominal 2022) le poids économique de chacun. Trace deux flèches de flux Sud-Sud majeurs (ex: Brésil $\to$ Chine soja/fer ; Inde $\to$ Afrique pharmacie). Légende complète.
\item \textbf{Analyse de document (Texte) :} « L'émergence n'est pas un long fleuve tranquille. Elle produit de nouvelles inégalités, détruit l'environnement et expose les pays aux crises financières venues du Nord. » À l'aide de deux exemples précis parmi les pays étudiés, discute cette affirmation en distinguant les facteurs internes et externes de ces fragilités. (15 lignes environ).
\item \textbf{Calcul :} Le PIB nominal du Vietnam était de $\SI{340}{milliards}\ \$US$ en 2020 et $\SI{410}{milliards}\ \$US$ en 2022. Calcule le taux de variation moyen annuel (TCAM) sur cette période. (Indice : $TCAM = \left(\frac{V_f}{V_i}\right)^{1/n} - 1$).
\end{enumerate}
\end{exercice}

\begin{corrige}[Exercices]
\begin{enumerate}
\item \textbf{Croquis attendu :} Cercles : Chine (rayon $\approx 9$ cm - \textit{attention, Chine trop grande pour cette échelle sur A4, adapter échelle ou utiliser surface proportionnelle / ou représenter seulement les 5 BRICS + Mexique/Indonésie/Turquie/Arabie/Nigéria/Maroc/Vietnam avec échelle 1 cm pour 3000 Md}). \textbf{Correction échelle suggérée :} Surface cercle $\propto$ PIB. $\text{Rayon} = k\sqrt{\text{PIB}}$. Chine $\SI{18\,000}{Md}$ $\to$ Rayon max 3 cm $\to k = 3/\sqrt{18000} \approx 0,022$. Inde ($\SI{3400}{Md}$) $\to$ 1,3 cm. Brésil ($\SI{1900}{Md}$) $\to$ 1,0 cm. Russie ($\SI{2200}{Md}$) $\to$ 1,0 cm. Af. Sud ($\SI{400}{Md}$) $\to$ 0,45 cm. Mexique ($\SI{1400}{Md}$) $\to$ 0,8 cm. Indonésie ($\SI{1300}{Md}$) $\to$ 0,8 cm. Turquie ($\SI{900}{Md}$) $\to$ 0,65 cm. Arabie ($\SI{1100}{Md}$) $\to$ 0,7 cm. Nigéria ($\SI{470}{Md}$) $\to$ 0,5 cm. Maroc ($\SI{130}{Md}$) $\to$ 0,25 cm. Vietnam ($\SI{410}{Md}$) $\to$ 0,45 cm. Flèches : Brésil $\to$ Chine (Fer/Soja), Golfe/Chine (Pétrole), Inde $\to$ Afrique (Médicaments), Chine $\to$ ASEAN (Composants électroniques). Légende : Titre, Auteur, Date, Échelle cercles, Signification couleurs, Flux.
\item \textbf{Analyse attendue :} Structure : Intro (définition émergence + thèse : affirmation juste mais partielle). §1 : Inégalités internes (Facteur interne : modèle de croissance excluant ; Exemple : Afrique du Sud Gini 0,63 / Brésil favelas vs Faria Lima). §2 : Environnement (Facteur interne : priorité croissance / Externe : demande mondiale matières 1ères ; Exemple : Déforestation Amazonie pour soja Chine / Pollution Delhi usines délocalisées). §3 : Vulnérabilité crises (Facteur externe : interdépendance financière/commerciale ; Exemple : Fuite capitaux Turquie/Argentine 2018 / Chute cours pétrole 2020 Nigeria/Angola). Conclusion : L'émergence est un processus contradictoire, pas un aboutissement. Nécessité politiques inclusives/vertes/résilientes.
\item \textbf{Calcul TCAM :} $n=2$ ans. $V_i=340$, $V_f=410$.
\[ TCAM = \left(\frac{410}{340}\right)^{1/2} - 1 = (1,2059)^{0,5} - 1 \approx 1,0981 - 1 = 0,0981 \approx \textbf{9,8 \% par an}. \]
Forte croissance confirmant le statut d'émergent dynamique (Vietnam).
\end{enumerate}
\end{corrige}

\coursec{TP guidé : cartographier et représenter les foyers émergents}

Après avoir analysé les facteurs, les manifestations et les limites de l'émergence, tu disposes désormais de toutes les connaissances nécessaires. Mais un géographe ne se contente pas de savoir : il \cle{représente}, il \cle{cartographie}, il \cle{met en images} l'espace qu'il étudie. Cette dernière partie de la leçon est entièrement consacrée aux travaux pratiques : tu vas apprendre à construire une carte mondiale des foyers émergents, un croquis de l'Afrique, des diagrammes statistiques, puis à analyser des documents pour nuancer la notion d'émergence. Ce sont exactement les compétences exigées au Baccalauréat dans les épreuves de croquis et d'étude de documents.

\courssub{TP 1 — Cartographier les foyers émergents du monde}

\textbf{Objectif.} À partir d'un fond de carte du monde vierge, localiser et légender les \cle{BRICS} (Brésil, Russie, Inde, Chine, Afrique du Sud), les quatre \cle{Dragons} (Corée du Sud, Taïwan, Singapour, Hong Kong), les pays du \cle{Golfe} (Arabie Saoudite, Émirats Arabes Unis) et les quatre émergents africains (Afrique du Sud, Nigéria, Égypte, Maroc).

\textbf{Matériel.} Fond de planisphère vierge, crayons de couleur, règle, feutre fin noir.

\textbf{Protocole.}
\begin{enumerate}
\item Repérer chaque pays sur l'atlas, puis le colorier en \emph{aplat léger} sur le fond vierge.
\item Attribuer un \cle{code couleur} par catégorie : une couleur pour les BRICS, une autre pour les Dragons, une autre pour les pays du Golfe, une autre pour les émergents africains.
\item Construire la légende \emph{en bas à droite}, organisée en catégories titrées, dans l'ordre logique : d'abord les ensembles (BRICS), puis les groupes régionaux (Dragons, Golfe), puis les cas africains.
\item Ajouter le titre en haut, en capitales, et une flèche du nord si le fond l'exige.
\end{enumerate}

\begin{methode}[Règles d'or du croquis et de la carte géographiques]
\begin{enumerate}
\item \textbf{Titre} : obligatoire, précis, placé en haut, en capitales (ex. : LES FOYERS ÉCONOMIQUES EN EXPANSION DANS LE MONDE).
\item \textbf{Échelle indicative} : mentionner « sans échelle » ou donner un ordre de grandeur ; ne jamais inventer une échelle fausse.
\item \textbf{Légende organisée en catégories} : regrouper les symboles par familles (foyers, villes, flux), chaque famille précédée d'un sous-titre ; jamais de liste désordonnée.
\item \textbf{Flèches orientées pour les flux} : un flux (pétrole, marchandises, capitaux) se représente par une flèche dont l'épaisseur est proportionnelle à l'importance du flux.
\item \textbf{Propreté du tracé} : traits nets, aplats légers qui laissent voir les noms, écriture horizontale et lisible.
\end{enumerate}
\end{methode}

\begin{attention}[Erreur fréquente : la carte saturée]
La tentation est grande de tout mettre sur la carte : capitales, reliefs, fleuves, climats... C'est une erreur grave. Retiens la règle : \cle{une information = un symbole}, et \emph{ne représenter que ce qui sert le sujet}. Si le sujet porte sur les foyers émergents, le fleuve Amazone n'a rien à y faire. Un correcteur pénalise une carte illisible autant qu'une carte incomplète.
\end{attention}

\begin{popfigure}
\begin{center}
\begin{tikzpicture}[font=\small]
\draw[fill=popcream,draw=popline] (0,0) rectangle (8.6,5.2);
\node[font=\bfseries\color{popdark}] at (4.3,4.8) {Modèle de légende hiérarchisée};
\draw[fill=popblue!70,draw=popdark] (0.5,3.9) rectangle (1.1,4.3);
\node[right] at (1.3,4.1) {\textbf{1. Les BRICS} : Brésil, Russie, Inde, Chine, Afrique du Sud};
\draw[fill=poporange!80,draw=popdark] (0.5,3.1) rectangle (1.1,3.5);
\node[right] at (1.3,3.3) {\textbf{2. Les Dragons} : Corée du Sud, Taïwan, Singapour, Hong Kong};
\draw[fill=popgold!85,draw=popdark] (0.5,2.3) rectangle (1.1,2.7);
\node[right] at (1.3,2.5) {\textbf{3. Le Golfe} : Arabie Saoudite, Émirats Arabes Unis};
\draw[fill=popgreen!70,draw=popdark] (0.5,1.5) rectangle (1.1,1.9);
\node[right] at (1.3,1.7) {\textbf{4. Émergents d'Afrique} : Nigéria, Égypte, Maroc (+ Afr. du Sud)};
\draw[-{Stealth[length=3mm]},line width=2pt,popink] (0.6,0.55) -- (1.6,0.55);
\node[right] at (1.8,0.55) {\textbf{5. Flux} : pétrole et marchandises (épaisseur = importance)};
\end{tikzpicture}
\end{center}
\legende{Modèle de légende hiérarchisée pour la carte mondiale des foyers émergents : les catégories sont numérotées et ordonnées, chaque couleur correspond à une seule catégorie de pays.}
\end{popfigure}

\courssub{TP 2 — Croquis : « Les foyers économiques en expansion en Afrique »}

\textbf{Consigne.} Sur un fond de carte de l'Afrique, construis un croquis faisant apparaître les quatre foyers émergents (Afrique du Sud, Nigéria, Égypte, Maroc), les grands ports (Le Cap, Durban, Lagos, Port-Saïd/Suez, Tanger-Med), et les flux pétroliers (golfe de Guinée) et commerciaux (détroit de Gibraltar, canal de Suez).

\begin{experience}[Démarche guidée pas à pas]
\begin{enumerate}
\item \textbf{Tracer le contour} du continent à main levée mais avec soin, en repérant trois points de calage : le cap Bon (Tunisie) au nord, le cap des Aiguilles au sud, le cap Vert à l'ouest.
\item \textbf{Placer les foyers} : colorier légèrement les quatre pays émergents ; marquer d'un point nommé les métropoles (Johannesburg, Lagos, Le Caire, Casablanca).
\item \textbf{Localiser les portes de sortie} : Tanger-Med face à Gibraltar, Port-Saïd à l'entrée du canal de Suez, Lagos sur le golfe de Guinée, Durban sur l'océan Indien.
\item \textbf{Dessiner les flux} : flèches épaisses pour le pétrole du golfe de Guinée vers l'Europe et l'Asie ; flèches moyennes pour les conteneurs passant par Suez et Gibraltar.
\item \textbf{Rédiger la légende} en trois catégories ordonnées : foyers, villes-ports, flux.
\end{enumerate}
\end{experience}

\begin{popfigure}
\includegraphics[width=\linewidth,height=11cm,keepaspectratio]{N14/afrique_pol.jpg}
\legende{Fond de carte pour le corrigé du croquis (carte politique de l'Afrique, en anglais). Le croquis attendu colore quatre foyers (Afrique du Sud, Nigéria, Égypte, Maroc), place les métropoles en points (Lagos, Le Caire, Johannesburg, Casablanca), les ports en carrés, les flux pétroliers du golfe de Guinée en flèches épaisses et les flux de conteneurs par Gibraltar et Suez en flèches fines. La carte fournie ne montre que les États, les capitales et les grandes villes ; échelle : voir la barre graduée en bas à droite. \\ Source : Wikimedia Commons, Central Intelligence Agency (États-Unis), domaine public.}
\end{popfigure}

\courssub{TP 3 — Construire des diagrammes statistiques}

\textbf{Données fournies.} PIB des BRICS en 2022 (ordres de grandeur, en milliers de milliards de dollars, source Banque mondiale) :

\begin{center}
\begin{tabularx}{\linewidth}{|l|X|X|X|X|X|}
\hline
\rowcolor{popblueL} \textbf{Pays} & Chine & Inde & Brésil & Russie & Afrique du Sud \\
\hline
\textbf{PIB (10\textsuperscript{3} Mds USD)} & 17,9 & 3,4 & 1,9 & 2,2 & 0,4 \\
\hline
\end{tabularx}
\end{center}

\begin{exoresolu}[Construire le diagramme en bâtons du PIB des BRICS]
Énoncé : représente graphiquement le PIB des cinq BRICS et calcule la part de la Chine dans le PIB total du groupe.
\tcblower
\textbf{Étape 1 — Calcul.} PIB total : $17,9 + 3,4 + 1,9 + 2,2 + 0,4 = 25,8$ milliers de milliards de dollars. Part de la Chine : $\dfrac{17,9}{25,8} \times 100 \approx 69,4\,\%$.

\textbf{Étape 2 — Choix de l'échelle.} La valeur maximale est 17,9 ; on choisit 1 cm pour 2 milliers de milliards, soit un bâton chinois de 9 cm environ.

\textbf{Étape 3 — Tracé.} Bâtons de largeur égale, régulièrement espacés, hauteur proportionnelle au PIB ; axe vertical gradué et titré (« PIB en milliers de milliards USD »), titre du graphique en haut.

\textbf{Lecture.} Le diagramme révèle l'écrasante domination chinoise : les BRICS sont un groupe très \cle{hétérogène}, ce qui limite leur capacité d'action commune.
\end{exoresolu}

\begin{popfigure}
\begin{center}
\begin{tikzpicture}
\begin{axis}[
  ybar, bar width=0.55cm,
  width=0.85\linewidth, height=6cm,
  ymin=0, ymax=20,
  ylabel={PIB (milliers de milliards USD)},
  symbolic x coords={Chine,Inde,Russie,Br\'esil,Afr. du Sud},
  xtick=data,
  nodes near coords, every node near coord/.append style={font=\scriptsize},
  ymajorgrids, grid style={popline!40},
  axis line style={popdark},
]
\addplot[fill=popblue!70,draw=popdark] coordinates {(Chine,17.9) (Inde,3.4) (Russie,2.2) (Br\'esil,1.9) (Afr. du Sud,0.4)};
\end{axis}
\end{tikzpicture}
\end{center}
\legende{Diagramme en bâtons du PIB des BRICS en 2022 (ordres de grandeur, source : Banque mondiale). La Chine pèse à elle seule près de 70\,\% du total du groupe.}
\end{popfigure}

\textbf{Diagramme circulaire de la population mondiale.} Pour convertir des effectifs en angles, applique la proportionnalité : $\text{angle} = \dfrac{\text{effectif}}{\text{total}} \times 360°$. Avec une population mondiale d'environ 8 milliards d'habitants : Asie $\approx 4,7$ Mds ($\approx 59\,\%$, soit $\approx 213°$), Afrique $\approx 1,4$ Mds ($\approx 18\,\%$, soit $\approx 63°$), Europe $\approx 0,75$ Md ($\approx 9\,\%$, soit $\approx 34°$), Amériques $\approx 1$ Md ($\approx 13\,\%$, soit $\approx 46°$), Océanie $\approx 0,04$ Md ($\approx 1\,\%$, soit $\approx 4°$). Vérifie toujours que la somme des angles fait 360°.

\courssub{TP 4 — Analyse croisée de documents : nuancer l'émergence}

\textbf{Document A (photographie décrite).} La \emph{skyline} de Dubaï : tours de verre, dont le Burj Khalifa (828 m), autoroutes à dix voies, marina luxueuse. \textbf{Document B (photographie décrite).} Une favela de Rio de Janeiro : habitat précaire accroché à la colline, à quelques centaines de mètres des quartiers riches. \textbf{Document C (tableau).} IDH et PIB par habitant de quelques pays émergents.

\begin{center}
\begin{tabularx}{\linewidth}{|l|X|X|}
\hline
\rowcolor{popblueL} \textbf{Pays} & \textbf{PIB/habitant (USD, ordre de grandeur)} & \textbf{IDH (classement approximatif)} \\
\hline
Émirats Arabes Unis & 40 000 -- 45 000 & Très élevé (top 40) \\
\hline
Brésil & 8 000 -- 9 000 & Élevé (vers la 90\textsuperscript{e} place) \\
\hline
Afrique du Sud & 6 000 -- 7 000 & Moyen (vers la 110\textsuperscript{e} place) \\
\hline
Inde & 2 300 -- 2 500 & Moyen (vers la 130\textsuperscript{e} place) \\
\hline
\end{tabularx}
\end{center}

\begin{exoresolu}[Commenter les documents pour nuancer la notion d'émergence]
Énoncé : à partir des trois documents, montre que l'émergence est un processus réel mais inachevé et inégal.
\tcblower
\textbf{Étape 1 — Décrire.} Le document A montre une modernité spectaculaire (gratte-ciel, infrastructures) ; le document B montre la persistance de la pauvreté au cœur d'un pays pourtant émergent ; le document C révèle des écarts considérables de richesse par habitant entre pays dits « émergents ».

\textbf{Étape 2 — Analyser.} L'émergence est d'abord une émergence de \emph{la richesse nationale globale} (PIB), non de \emph{tous les habitants}. La croissance profite inégalement : les Émirats affichent un niveau de vie de pays développé, tandis que l'Inde reste un pays à faible revenu par habitant malgré un PIB total immense.

\textbf{Étape 3 — Conclure.} L'émergence est un \cle{processus inachevé} : elle transforme les États et les littoraux avant les campagnes et les marges urbaines. D'où la nécessité d'indicateurs comme l'\cle{IDH}, qui intègrent santé et éducation, pour juger du développement réel.
\end{exoresolu}

\courssub{Étude de cas : le Cameroun, pays à revenus intermédiaires}

\begin{exemplebox}[Le Cameroun parmi les pays à revenus intermédiaires]
Le Cameroun est classé par la Banque mondiale parmi les \cle{pays à revenus intermédiaires} (tranche inférieure), avec un PIB par habitant d'environ \num{1500} à \num{1700} USD (ordre de grandeur). Il se situe ainsi au-dessus des PPTE mais loin derrière les émergents africains comme l'Afrique du Sud ou le Maroc.
\end{exemplebox}

\textbf{Les atouts du Cameroun.}
\begin{itemize}
\item Des \cle{ressources pétrolières et gazières} (golfe de Guinée, bassin du Rio del Rey), qui assurent une part importante des recettes d'exportation ;
\item Une \cle{agriculture diversifiée} et exportatrice : cacao (le Cameroun figure parmi les premiers producteurs mondiaux), café, coton, bananes, bois ;
\item Une \cle{position stratégique} sur le golfe de Guinée, charnière entre l'Afrique centrale et l'Afrique de l'Ouest, avec les ports de Douala et de Kribi ;
\item Une population jeune et un marché qui irrigue la CEMAC.
\end{itemize}

\textbf{Les obstacles à l'émergence.} Dépendance aux matières premières brutes (faible transformation locale), insuffisance des infrastructures (énergie, routes), gouvernance et corruption dénoncées par les rapports internationaux, crises sécuritaires (Extrême-Nord, régions anglophones), endettement croissant.

\begin{savaistu}[L'ambition camerounaise : émergence à l'horizon 2035]
Le Cameroun a inscrit son ambition d'émergence à l'horizon 2035 dans sa \cle{SND30} (Stratégie Nationale de Développement 2020-2030, prolongeant la vision 2035). La stratégie s'appuie sur les grandes infrastructures — le \cle{port en eau profonde de Kribi}, le barrage de Nachtigal — et sur l'industrialisation par la transformation locale des matières premières.
\end{savaistu}

\begin{definition}[PPTE]
Les \cle{PPTE} (Pays Pauvres Très Endettés) sont les pays les plus pauvres du monde, écrasés par une dette extérieure insoutenable ; ils bénéficient d'initiatives internationales d'\emph{allègement} ou d'annulation de la dette (initiative PPTE du FMI et de la Banque mondiale, 1996). Ils illustrent le \emph{décollage économique difficile}, qui fera l'objet de la leçon 15.
\end{definition}

\courssub{Restitution et évaluation des productions}

Chaque groupe présente sa production (carte, croquis ou diagramme) en trois minutes, en \textbf{justifiant} : le choix des couleurs et des symboles, l'organisation de la légende, l'échelle du diagramme. La classe évalue selon la grille suivante :

\begin{center}
\begin{tabularx}{\linewidth}{|l|X|}
\hline
\rowcolor{popblueL} \textbf{Critère} & \textbf{Attendu} \\
\hline
Exactitude & Pays et villes correctement localisés, calculs justes \\
\hline
Légende & Organisée en catégories numérotées, un symbole par information \\
\hline
Lisibilité & Tracé propre, carte non surchargée, titre en capitales \\
\hline
Argumentation & Choix justifiés, vocabulaire géographique précis \\
\hline
\end{tabularx}
\end{center}

\begin{aretenir}[L'essentiel du TP]
Cartographier les foyers émergents, c'est montrer une \cle{mondialisation à plusieurs vitesses} : BRICS hétérogènes dominés par la Chine, Dragons intégrés aux échanges high-tech, Golfe riche de sa rente pétrolière, émergents africains encore fragiles. Le Cameroun, pays à revenus intermédiaires, possède des atouts réels (pétrole, agriculture, façade atlantique) mais n'est pas encore un pays émergent : son horizon est 2035. Un bon croquis se reconnaît à trois qualités : un titre, une légende ordonnée, et rien de superflu.
\end{aretenir}

\newpage

\coursec{Activité d'intégration}

\courssub{Objectifs de l'activité}

Cette activité d'intégration te place dans la peau d'un expert-conseil. À l'issue du travail, tu devras être capable de :

\begin{itemize}
\item \cle{localiser} et \cle{cartographier} le Cameroun et les pays émergents étudiés sur un fond de carte de l'Afrique (croquis légendé, orienté, avec échelle indicative) ;
\item \cle{construire} un diagramme en bâtons à partir d'un tableau statistique, en respectant les règles de représentation graphique (titre, axes légendés, unités, source) ;
\item \cle{inventorier}, \cle{classer} et \cle{interpréter} des documents variés (carte, tableau, photographies décrites) pour dégager atouts et obstacles ;
\item \cle{formuler} des recommandations argumentées en mobilisant les modèles d'émergence étudiés dans la leçon (Maroc, pays du Golfe, Dragons asiatiques) ;
\item rédiger une note de synthèse brève et la présenter oralement, comme au Baccalauréat où la rigueur de la légende et la précision des chiffres sont évaluées.
\end{itemize}

\courssub{Situation}

\textbf{Yaoundé, ministère de l'Économie, de la Planification et de l'Aménagement du Territoire (MINEPAT).} Dans le cadre du suivi de la Stratégie Nationale de Développement 2020--2030 (SND30), qui ambitionne de faire du Cameroun un \cle{pays émergent}, le ministère lance un appel à notes d'expertise. Ta classe est constituée en quatre cabinets de conseil. Chaque cabinet doit rédiger une note de deux pages maximum répondant à la question : \emph{le Cameroun peut-il devenir un foyer économique en expansion, et à quelles conditions ?}

\textbf{Le dossier documentaire remis aux cabinets}

\textbf{Document 1 — Tableau comparatif (ordres de grandeur, début des années 2020, sources : Banque mondiale, PNUD ; valeurs arrondies).}

\begin{center}
\begin{tabularx}{\linewidth}{|l|X|X|X|X|}
\hline
\rowcolor{popblueL} \textbf{Indicateur} & \textbf{Cameroun} & \textbf{Nigéria} & \textbf{Maroc} & \textbf{Vietnam} \\
\hline
PIB (milliards USD) & environ \num{45} & environ \num{450} & environ \num{130} & environ \num{360} \\
\hline
PIB par habitant (USD) & environ \num{1600} & environ \num{2200} & environ \num{3500} & environ \num{3700} \\
\hline
Croissance annuelle moyenne & 3 à 4 \% & 2 à 3 \% & 2 à 3 \% & 6 à 7 \% \\
\hline
IDH (rang approx.) & autour de 150\textsuperscript{e} & hors top 150 & vers le 120\textsuperscript{e} & vers le 115\textsuperscript{e} \\
\hline
Exportations dominantes & pétrole brut, cacao, bois, coton & pétrole et gaz (plus de 85 \%) & phosphates, automobiles, textiles, agrumes & électronique, textiles, chaussures, riz \\
\hline
\end{tabularx}
\end{center}

\textbf{Document 2 — Photographies décrites.} Photo A : le \cle{port en eau profonde de Kribi} (Sud), inauguré en 2018, terminal à conteneurs moderne, zone industrielle et portuaire en extension ; le trafic y progresse fortement. Photo B : le port de \cle{Tanger Med} (Maroc), première capacité portuaire d'Afrique et de Méditerranée (capacité de 9 millions d'EVP, ordre de grandeur 2024), adossé à des zones industrielles où sont implantés Renault et Stellantis.

\textbf{Document 3 — Extraits de données camerounaises.} Le pétrole représente encore une part importante des recettes d'exportation (avec le cacao, premier poste agricole : le Cameroun est environ le 5\textsuperscript{e} producteur mondial). Le projet gazier de Kribi, le barrage de Nachtigal (\SI{420}{MW}, mis en service progressif à partir de 2024) et le corridor douanier vers le Tchad et la Centrafrique sont des chantiers structurants. Mais le pays reste classé parmi les \cle{pays à revenu intermédiaire inférieur} (PRII), avec une dette maîtrisée mais tendue et une industrialisation limitée (moins de 15 \% du PIB pour l'industrie manufacturière, ordre de grandeur).

\courssub{Consignes}

\begin{enumerate}
\item Sur un fond de carte de l'Afrique, localise et nomme le Cameroun, le Nigéria, le Maroc, l'Égypte et l'Afrique du Sud ; place les ports de Kribi et de Tanger Med. Légende, orientation et titre obligatoires.
\item À partir du document 1, construis un diagramme en bâtons comparant le PIB des quatre pays (Cameroun, Nigéria, Maroc, Vietnam). Choisis une échelle adaptée et justifie ton choix.
\item Calcule la part du PIB du Cameroun dans le total des quatre pays, puis compare le PIB par habitant du Cameroun à celui du Maroc. Que conclus-tu sur le niveau relatif du Cameroun ?
\item À partir des documents 2 et 3 et de ton cours, dégage \textbf{deux atouts} et \textbf{deux obstacles} à l'émergence du Cameroun, en les classant (atouts/obstacles naturels, économiques, humains).
\item Formule \textbf{deux recommandations} argumentées en t'inspirant des modèles étudiés : le modèle marocain (plateforme industrielle et portuaire) et le modèle du Golfe (diversification de la rente).
\item Rédige la note de synthèse (15 à 20 lignes) et prépare une présentation orale de 5 minutes. Barème : exactitude des chiffres (4 pts), qualité cartographique et graphique (6 pts), argumentation (6 pts), présentation (4 pts).
\end{enumerate}

\courssub{Ressources à mobiliser}

\begin{aretenir}[Ce qu'il faut réactiver avant de commencer]
\begin{itemize}
\item Un \cle{pays émergent} est un pays en voie de développement dont la croissance est forte et durable, qui s'industrialise, s'intègre aux échanges mondiaux et voit le niveau de vie de sa population progresser (BRICS, NPI, Dragons).
\item Les critères d'émergence : croissance du PIB, diversification de l'économie, industrialisation, insertion dans les flux mondiaux (ports, IDE), amélioration de l'IDH.
\item Règles du croquis : titre, légende numérotée ou fléchée, orientation, noms en capitales pour les États.
\item Règles du diagramme en bâtons : titre, axes légendés avec unité, échelle régulière, source.
\item Part en pourcentage : $\text{part} = \dfrac{\text{valeur de la partie}}{\text{valeur du total}} \times 100$.
\end{itemize}
\end{aretenir}

\courssub{Démarche et corrigé commenté}

\begin{exoresolu}[Note d'expertise : le Cameroun, futur foyer économique en expansion ?]
Énoncé : à l'aide du dossier documentaire, produire les documents demandés (croquis, diagramme, calculs) et rédiger la note d'expertise répondant à la question : le Cameroun peut-il devenir un foyer économique en expansion ?
\tcblower
\textbf{Étape 1 — Le croquis de localisation.} Sur le fond de carte de l'Afrique, on place : le MAROC et le port de Tanger Med au nord (façade méditerranéenne), l'ÉGYPTE au nord-est, le NIGÉRIA et le CAMEROUN en Afrique de l'Ouest et centrale (golfe de Guinée, port de Kribi), l'AFRIQUE DU SUD à l'extrémité australe. La légende distingue les pays émergents africains étudiés (une couleur), le Cameroun (une couleur contrastée) et les ports stratégiques (symbole d'ancre). La flèche du nord pointe vers le haut.

\begin{popfigure}
\includegraphics[width=\linewidth,height=11cm,keepaspectratio]{N14/afrique_pol.jpg}
\legende{Carte politique de l'Afrique (en anglais) : États, capitales (étoiles) et principales villes, du Maroc à l'Afrique du Sud, avec le golfe de Guinée à l'ouest et la mer Rouge à l'est. Elle sert de fond de carte pour le croquis de localisation : l'élève y repère le Maroc (au nord-ouest, face au détroit de Gibraltar où se trouve Tanger Med), l'Égypte (au nord-est), le Nigéria et le Cameroun autour du golfe de Guinée (Douala et Yaoundé sont figurées) et l'Afrique du Sud (à l'extrémité australe), puis place Kribi (au sud de Douala) et les ports stratégiques. La carte ne figure ni le port de Kribi ni celui de Tanger Med. \\ Source : Wikimedia Commons, Central Intelligence Agency (États-Unis), domaine public.}
\end{popfigure}

\textbf{Étape 2 — Le diagramme en bâtons.} Les valeurs vont de 45 à 450 milliards USD : on choisit une échelle de 1 cm pour 50 milliards, ce qui donne des bâtons de 0,9 cm (Cameroun) à 9 cm (Nigéria), lisibles sur une page.

\begin{popfigure}
\begin{tikzpicture}
\begin{axis}[
  width=0.85\linewidth, height=6cm,
  ybar, bar width=0.9cm,
  ymin=0, ymax=500,
  ylabel={PIB (milliards USD)},
  symbolic x coords={Cameroun, Nigeria, Maroc, Vietnam},
  xtick=data,
  ytick={0,100,200,300,400,500},
  nodes near coords, nodes near coords style={font=\scriptsize, anchor=south, yshift=2pt},
  axis line style={popdark}, tick label style={font=\scriptsize},
]
\addplot[fill=popblue, draw=popdark] coordinates {(Cameroun,45) (Nigeria,450) (Maroc,130) (Vietnam,360)};
\end{axis}
\end{tikzpicture}
\legende{PIB comparé du Cameroun, du Nigéria, du Maroc et du Vietnam (ordres de grandeur, début des années 2020). Source : Banque mondiale, valeurs arrondies.}
\end{popfigure}

\textbf{Étape 3 — Les calculs.} Total des quatre PIB :
\[ 45 + 450 + 130 + 360 = 985 \text{ milliards USD.} \]
Part du Cameroun :
\[ \text{part} = \frac{45}{985} \times 100 \approx 4,6\,\%. \]
Comparaison du PIB par habitant : $\dfrac{3500}{1600} \approx 2,2$. Le Marocain moyen dispose donc d'un revenu environ \textbf{2,2 fois supérieur} à celui du Camerounais moyen. Conclusion : le Cameroun pèse moins de 5 \% du total étudié et accuse un retard de revenu par habitant ; il n'est pas encore un foyer économique en expansion, mais un \cle{pays à revenu intermédiaire} en décollage difficile.

\textbf{Étape 4 — Atouts et obstacles.}
\begin{itemize}
\item \textbf{Atout 1 (naturel et économique)} : une agriculture et des ressources diversifiées (cacao, pétrole, gaz, bois, coton) — le Cameroun est « l'Afrique en miniature ».
\item \textbf{Atout 2 (économique)} : des infrastructures structurantes récentes (port de Kribi, barrage de Nachtigal) et une position de corridor vers le Tchad et la Centrafrique.
\item \textbf{Obstacle 1 (économique)} : une faible industrialisation (moins de 15 \% du PIB) et des exportations encore dominées par des produits bruts, comme le Nigéria avec son pétrole.
\item \textbf{Obstacle 2 (humain et institutionnel)} : un IDH encore faible, une gouvernance et un climat des affaires perfectibles qui freinent les IDE.
\end{itemize}

\textbf{Étape 5 — Recommandations.}
\begin{enumerate}
\item \emph{S'inspirer du Maroc} : transformer Kribi en véritable plateforme industrielle et portuaire (zones franches, transformation locale du cacao et du bois, attractivité fiscale) pour passer de l'exportation brute à l'exportation de produits transformés.
\item \emph{S'inspirer du Golfe} : investir la rente pétrolière et gazière dans la diversification (agro-industrie, énergie, numérique, formation) plutôt que dans la consommation, afin de préparer l'après-pétrole.
\end{enumerate}

\textbf{Étape 6 — Note de synthèse (modèle).} Le Cameroun possède des atouts réels : ressources variées, position stratégique en Afrique centrale, infrastructures nouvelles. Mais son PIB (environ 45 milliards USD, soit 4,6 \% du total étudié), son revenu par habitant deux fois inférieur à celui du Maroc et sa faible industrialisation montrent qu'il reste un pays à revenu intermédiaire au décollage difficile. Il peut devenir un foyer économique en expansion à condition de transformer localement ses matières premières, de faire de Kribi une plateforme industrielle à l'image de Tanger Med et de diversifier son économie comme les pays du Golfe. L'émergence n'est donc pas un destin, mais un chantier.
\end{exoresolu}

\begin{attention}[Pièges fréquents]
Ne confonds pas \cle{PIB} et \cle{PIB par habitant} : le Nigéria a le plus gros PIB du tableau, mais pas le revenu par habitant le plus élevé. N'oublie jamais le titre, la source et l'unité sur un diagramme. Enfin, évite les affirmations sans chiffre : au Baccalauréat, chaque argument doit être appuyé sur une donnée du dossier.
\end{attention}

\courssub{Bilan de l'activité}

Cette activité t'a permis de mobiliser ensemble les trois savoir-faire du programme : \cle{localiser et cartographier} (croquis de l'Afrique), \cle{construire un diagramme} (bâtons du PIB) et \cle{interpréter des documents} (tableau, photographies, données). Elle a surtout mis en évidence une idée centrale de la leçon : l'émergence se mesure à des critères précis (croissance, diversification, industrialisation, insertion mondiale, IDH) et non à une simple richesse en ressources. Le Cameroun, comme le Nigéria, illustre le paradoxe d'un pays riche en matières premières mais encore peu transformateur ; le Maroc et le Vietnam montrent au contraire que l'intégration aux chaînes de valeur mondiales (ports, usines, exportations manufacturées) est le véritable moteur des foyers économiques en expansion. Tu es désormais capable, à la manière d'un rapport d'experts, d'évaluer la trajectoire d'un pays et de proposer des stratégies argumentées : c'est exactement la compétence attendue au Baccalauréat.

\newpage

\coursec{Méthodes et exercices résolus}

\courssub{Fiche méthode 1 : Observer, localiser et décrire les foyers économiques en expansion}

\begin{definition}[Pays émergent]
Un \cle{pays émergent} est un pays en voie de développement rapide, qui connaît une forte croissance économique (souvent supérieure à 4 \% par an), une industrialisation en cours et une insertion croissante dans la mondialisation, tout en conservant un revenu par habitant intermédiaire et de fortes inégalités internes. On parle aussi de \cle{foyer économique en expansion}.
\end{definition}

\begin{methode}[Localiser et décrire un foyer émergent]
Pour observer, localiser et décrire un foyer économique en expansion (pays émergent), suis toujours la même démarche en quatre étapes :
\begin{enumerate}
\item \textbf{Situer} : nomme le continent, la sous-région, les pays voisins, les façades maritimes (ex. : le Brésil, en Amérique du Sud, façade atlantique, frontalier de dix pays).
\item \textbf{Identifier les indicateurs} : PIB (valeur absolue et rang mondial), PIB par habitant, IDH, taux de croissance, part dans le commerce mondial. Un pays émergent combine forte croissance (souvent supérieure à 4 \%) et revenu intermédiaire.
\item \textbf{Décrire les pôles internes} : localise les métropoles (São Paulo, Mexico, Johannesburg, Lagos, Le Caire, Casablanca, Dubaï), les zones industrialo-portuaires et les axes de développement.
\item \textbf{Qualifier} : classe le pays dans la bonne catégorie (BRICS, NPI, Dragon, pays du Golfe, émergent africain) et précise sa spécialisation (agro-industrie, hydrocarbures, manufacturière, services).
\end{enumerate}
\textbf{Réflexe} : une localisation sans échelle (monde $\rightarrow$ continent $\rightarrow$ pays $\rightarrow$ pôles internes) est une localisation incomplète.
\end{methode}

\begin{exoresolu}[Localiser et décrire un foyer émergent]
\textbf{Énoncé (type Bac).} À partir de vos connaissances, localisez le Brésil et décrivez en une dizaine de lignes pourquoi ce pays constitue un foyer économique en expansion.
\tcblower
\textbf{Solution détaillée.}
\begin{enumerate}
\item \emph{Localisation.} Le Brésil occupe la moitié orientale de l'Amérique du Sud ; il s'ouvre sur l'océan Atlantique par une façade d'environ \num{7400} km et partage des frontières avec dix pays sud-américains (tous sauf le Chili et l'Équateur). C'est le cinquième pays du monde par la superficie (environ 8,5 millions de km\textsuperscript{2}) et l'un des plus peuplés (environ 215 millions d'habitants, ordre de grandeur).
\item \emph{Indicateurs.} Avec un PIB de l'ordre de \num{2000} milliards de dollars (ordre de grandeur, parmi les douze premières économies mondiales) et un PIB par habitant voisin de \num{9000} dollars, le Brésil appartient aux pays à revenu intermédiaire et au groupe des BRICS (avec la Russie, l'Inde, la Chine et l'Afrique du Sud).
\item \emph{Pôles internes.} Le dynamisme se concentre dans le Sud-Est : São Paulo (première métropole industrielle et financière), Rio de Janeiro, Belo Horizonte (triangle minier du Minas Gerais), et les ports de Santos et de Paranaguá qui exportent le soja, le café, le sucre et la viande.
\item \emph{Spécialisations.} Le Brésil est une puissance agro-industrielle (premier exportateur mondial de soja, de café et de sucre, et l'un des premiers exportateurs de viande bovine), minière (minerai de fer de la société Vale) et industrielle (Embraer, l'un des premiers avionneurs mondiaux ; pétrole offshore de Petrobras).
\end{enumerate}
\textbf{Conclusion.} Foyer continental inséré dans la mondialisation par ses ports et ses firmes, le Brésil illustre un pays émergent complet, bien que marqué par de fortes inégalités internes entre le Sud-Est développé et le Nord amazonien en retard.
\end{exoresolu}

\courssub{Fiche méthode 2 : Construire et légender un croquis des foyers émergents}

\begin{methode}[Réaliser un croquis de localisation des foyers en expansion]
Au Baccalauréat, le croquis est noté sur trois critères : l'organisation de l'espace, la justesse des localisations, la propreté de la légende. Démarche :
\begin{enumerate}
\item \textbf{Tracer le fond} : esquisse un planisphère très simplifié (blocs continentaux) ou le contour du continent demandé. Ajoute la flèche du nord et le titre.
\item \textbf{Placer les foyers} : localise par des points ou des zones colorées les pays émergents demandés (BRICS, Dragons, NPI, pays du Golfe, émergents africains).
\item \textbf{Hiérarchiser} : utilise des symboles proportionnels (ronds de tailles différentes selon le poids économique) ou des couleurs distinctes par catégorie.
\item \textbf{Figurer les relations} : flèches de flux (exportations de pétrole du Golfe, flux manufacturiers des Dragons vers l'Europe et l'Amérique du Nord).
\item \textbf{Rédiger la légende} : numérotée ou fléchée, placée hors de la carte, chaque symbole doit y figurer. La légende vaut souvent autant de points que le dessin.
\end{enumerate}
\textbf{Réflexe} : reste schématique mais cohérent ; un croquis ne prétend jamais à la précision d'un atlas.
\end{methode}

\begin{exoresolu}[Croquis des BRICS et des principaux foyers émergents]
\textbf{Énoncé (type Bac).} Construisez un croquis intitulé « Les principaux foyers économiques en expansion dans le monde » en localisant : les cinq BRICS, les quatre Dragons asiatiques et deux foyers pétroliers du Moyen-Orient. Légendez votre croquis.
\tcblower
\textbf{Solution détaillée.} On procède en trois temps : fond de carte simplifié, placement des foyers, légende.

\begin{popfigure}
\includegraphics[width=\linewidth,height=8cm,keepaspectratio]{N14/brics.png}
\legende{\textbf{Les foyers économiques en expansion : le groupe des BRICS élargi.} Bleu foncé : membres du groupe (Brésil, Russie, Inde, Chine, Afrique du Sud, rejoints par l'Égypte, l'Éthiopie, l'Iran, l'Arabie saoudite, les Émirats arabes unis et l'Indonésie) ; bleu clair : pays partenaires ; gris : autres pays. Ce planisphère situe les BRICS sur quatre continents mais ne montre ni les quatre Dragons (Corée du Sud, Taïwan, Hong Kong, Singapour), ni les flux : le croquis attendu les ajoute sur un fond de carte comparable. \\ Source : Wikimedia Commons, Asikni (d'après BRICS.svg), CC0.}
\end{popfigure}

\textbf{Justification des choix.} Le planisphère fourni sert de fond de carte : le correcteur évalue la \emph{position relative} des foyers, pas la précision des côtes. Les BRICS sont répartis sur quatre continents, ce qui montre visuellement que l'émergence est un phénomène mondial ; sur son propre croquis, on ajoute en Asie orientale et du Sud-Est les quatre Dragons, et sur la péninsule Arabique les foyers pétroliers du Golfe (Arabie saoudite, Émirats arabes unis), reliés par des flèches d'exportation d'hydrocarbures vers l'Europe et l'Asie, qui rappellent leur insertion dans la mondialisation par le commerce.

\textbf{Conclusion.} Le croquis met en évidence trois grandes catégories de foyers en expansion : les BRICS (puissances continentales), les Dragons (ateliers et places financières) et les monarchies pétrolières du Golfe (rente énergétique réinvestie).
\end{exoresolu}

\courssub{Fiche méthode 3 : Inventorier et classer les pays émergents}

\begin{methode}[Classer les foyers économiques en expansion]
Face à une liste de pays ou à un tableau, le classement doit suivre des critères explicites :
\begin{enumerate}
\item \textbf{Critère continental} : Amérique (Brésil, Mexique, Argentine) ; Asie et Moyen-Orient (Inde, Malaisie, les quatre Dragons, Arabie Saoudite, Émirats arabes unis) ; Afrique (Afrique du Sud, Nigéria, Égypte, Maroc).
\item \textbf{Critère du statut} : BRICS (Brésil, Russie, Inde, Chine, Afrique du Sud) ; NPI, nouveaux pays industrialisés (Dragons : Corée du Sud, Taïwan, Hong Kong, Singapour ; puis Mexique, Malaisie...) ; pays du Golfe à rente pétrolière.
\item \textbf{Critère du moteur de croissance} : manufacturier (Dragons, Mexique avec les maquiladoras), agricole et minier (Brésil, Argentine), hydrocarbures (Golfe, Nigéria), services et tourisme (Émirats, Maroc, Égypte).
\item \textbf{Vérifier la cohérence} : un même pays peut appartenir à plusieurs catégories (l'Afrique du Sud est à la fois BRICS et premier foyer industriel africain) ; le préciser montre la maîtrise du sujet.
\end{enumerate}
\textbf{Réflexe} : annonce toujours le critère avant de classer ; un classement sans critère annoncé est sanctionné.
\end{methode}

\begin{exoresolu}[Classer douze foyers émergents]
\textbf{Énoncé (type Bac).} Voici douze pays : Afrique du Sud, Arabie Saoudite, Argentine, Brésil, Corée du Sud, Égypte, Émirats arabes unis, Inde, Malaisie, Maroc, Mexique, Singapour. Classez-les selon deux critères : le continent d'appartenance, puis le moteur principal de leur émergence.
\tcblower
\textbf{Solution détaillée.}

\emph{Premier classement : par continent.}
\begin{center}
\begin{tabularx}{\linewidth}{|l|X|}\hline
\rowcolor{popblueL} \textbf{Continent} & \textbf{Pays} \\ \hline
Amérique & Brésil, Mexique, Argentine \\ \hline
Asie et Moyen-Orient & Inde, Corée du Sud, Singapour, Malaisie, Arabie Saoudite, Émirats arabes unis \\ \hline
Afrique & Afrique du Sud, Égypte, Maroc \\ \hline
\end{tabularx}
\end{center}

\emph{Deuxième classement : par moteur de croissance.}
\begin{itemize}
\item \textbf{Moteur manufacturier} (industries d'exportation) : Corée du Sud et Singapour (Dragons : électronique, chantiers navals, place financière), Mexique (maquiladoras adossées au marché nord-américain), Malaisie (électronique).
\item \textbf{Moteur agricole et minier} : Brésil (soja, café, minerai de fer), Argentine (pampa : céréales, viande, soja).
\item \textbf{Moteur pétrolier et gazier} : Arabie Saoudite, Émirats arabes unis (qui diversifient vers la finance, le tourisme et l'aérien avec Dubaï et Abu Dhabi).
\item \textbf{Moteur mixte en construction} : Inde (services informatiques, industrie, démographie), Afrique du Sud (mines, industrie, finance), Égypte (canal de Suez, gaz, tourisme, textile), Maroc (automobile, aéronautique, phosphates, tourisme).
\end{itemize}

\textbf{Conclusion.} Ce double classement montre qu'il n'existe pas un modèle unique d'émergence : certains foyers s'insèrent dans la mondialisation par l'atelier (Dragons, Mexique), d'autres par le grenier et la mine (Brésil, Argentine), d'autres encore par la rente énergétique (Golfe) ; les émergents africains combinent plusieurs moteurs encore fragiles.
\end{exoresolu}

\courssub{Fiche méthode 4 : Interpréter un tableau statistique et calculer des indicateurs}

\begin{methode}[Exploiter des données chiffrées sur les pays émergents]
\begin{enumerate}
\item \textbf{Lire le tableau en entier} : titre, unités (milliards de dollars, \%, dollars par habitant), source, date. Ne commente jamais un chiffre sans son unité.
\item \textbf{Repérer les valeurs extrêmes} : maximum, minimum, et calcule l'écart (absolu et relatif).
\item \textbf{Calculer les indicateurs utiles} :
\begin{itemize}
\item Part en pourcentage : $\text{part} = \dfrac{\text{valeur de la partie}}{\text{valeur du total}} \times 100$.
\item Taux de variation : $\text{taux} = \dfrac{V_f - V_i}{V_i} \times 100$ (en \%).
\item Indice base 100 : $\text{indice} = \dfrac{V_f}{V_i} \times 100$.
\end{itemize}
\item \textbf{Interpréter géographiquement} : relie chaque chiffre à une réalité (un PIB par habitant de \num{8000} dollars signifie un revenu intermédiaire, donc un pays émergent, pas un pays développé).
\item \textbf{Conclure par une phrase de synthèse} qui répond à la question posée.
\end{enumerate}
\textbf{Réflexe} : au Bac, tout calcul doit être écrit (opération posée, résultat, unité) ; un résultat tombé du ciel n'est pas crédité.
\end{methode}

\begin{exoresolu}[Calculer et interpréter : le poids des BRICS]
\textbf{Énoncé (type Bac).} Le tableau suivant donne des ordres de grandeur du PIB en 2022 (en milliards de dollars courants) :

\begin{center}
\begin{tabularx}{\linewidth}{|l|X|X|X|X|X|X|}\hline
\rowcolor{popblueL} \textbf{Pays} & Chine & Inde & Brésil & Russie & Afrique du Sud & Monde \\ \hline
\textbf{PIB (mds \$)} & \num{18000} & \num{3400} & \num{1900} & \num{2200} & 400 & \num{101000} \\ \hline
\end{tabularx}
\end{center}

1. Calculez le PIB cumulé des cinq BRICS. 2. Calculez leur part dans le PIB mondial. 3. Calculez le taux de variation du PIB indien entre 2010 (\num{1700} milliards de dollars) et 2022. 4. Tirez-en deux conclusions géographiques.
\tcblower
\textbf{Solution détaillée.}

\emph{1. PIB cumulé des BRICS.}
\[ \num{18000} + \num{3400} + \num{1900} + \num{2200} + 400 = \num{25900} \text{ milliards de dollars.} \]

\emph{2. Part dans le PIB mondial.}
\[ \text{part} = \frac{\num{25900}}{\num{101000}} \times 100 \approx 25,6\,\%. \]
Les cinq BRICS pèsent donc environ un quart de la richesse mondiale produite en un an.

\emph{3. Taux de variation du PIB indien entre 2010 et 2022.}
\begin{align*}
\text{taux} &= \frac{V_f - V_i}{V_i} \times 100 \\
&= \frac{\num{3400} - \num{1700}}{\num{1700}} \times 100 \\
&= \frac{\num{1700}}{\num{1700}} \times 100 = 100\,\%.
\end{align*}
Le PIB indien a doublé en douze ans (indice 200, base 100 en 2010), ce qui traduit une croissance moyenne très rapide, caractéristique d'un pays émergent.

\emph{4. Conclusions géographiques.}
\begin{itemize}
\item Les BRICS sont devenus un poids lourd de l'économie mondiale : avec environ un quart du PIB mondial, ils contestent la domination historique de la Triade (Amérique du Nord, Europe de l'Ouest, Japon).
\item Ce poids est très inégalement réparti : la Chine concentre à elle seule $\frac{\num{18000}}{\num{25900}} \times 100 \approx 69\,\%$ du PIB du groupe, tandis que l'Afrique du Sud n'en représente que $\frac{400}{\num{25900}} \times 100 \approx 1,5\,\%$. Le « groupe » des BRICS est donc hétérogène : une superpuissance (Chine), des puissances moyennes en forte croissance (Inde, Brésil, Russie) et une puissance régionale africaine (Afrique du Sud).
\end{itemize}
\textbf{Conclusion.} Les chiffres confirment l'émergence comme fait majeur de la mondialisation, mais invitent à nuancer : l'émergence est d'abord asiatique et chinoise.
\end{exoresolu}

\courssub{Fiche méthode 5 : Construire un diagramme statistique (barres ou circulaire)}

\begin{methode}[Représenter graphiquement des données sur les foyers émergents]
\begin{enumerate}
\item \textbf{Choisir le bon graphique} : diagramme en barres pour comparer des valeurs (PIB de plusieurs pays) ; diagramme circulaire pour des parts d'un total (répartition en \%).
\item \textbf{Pour un diagramme circulaire}, convertir chaque part en angle : $\text{angle} = \dfrac{\text{part} \times 360}{100}$ (en degrés). Vérifier que la somme des angles vaut 360°.
\item \textbf{Tracer proprement} : axes gradués et légendés (grandeur + unité) pour les barres ; secteurs colorés et légendés pour le circulaire.
\item \textbf{Titrer et sourcer} : tout graphique non titré est pénalisé.
\item \textbf{Exploiter} : dégager un ou deux enseignements (hiérarchie, domination d'un foyer).
\end{enumerate}
\end{methode}

\begin{exoresolu}[Diagramme en barres : le PIB des foyers émergents]
\textbf{Énoncé (type Bac).} Représentez par un diagramme en barres le PIB (ordres de grandeur 2022, en milliards de dollars) des pays suivants : Chine \num{18000} ; Inde \num{3400} ; Brésil \num{1900} ; Mexique \num{1400} ; Afrique du Sud 400 ; Nigéria 500. Quel enseignement en tirez-vous ?
\tcblower
\textbf{Solution détaillée.} On choisit un diagramme en barres car on compare des valeurs absolues entre pays. L'axe vertical est gradué en milliards de dollars, l'axe horizontal porte les pays classés par PIB décroissant.

\begin{popfigure}
\begin{tikzpicture}
\begin{axis}[
  ybar, bar width=0.55cm,
  width=0.95\linewidth, height=6.2cm,
  ymin=0, ymax=20000,
  ylabel={PIB (milliards de dollars)},
  symbolic x coords={Chine, Inde, Bresil, Mexique, Nigeria, AfrSud},
  xtick=data, x tick label style={font=\scriptsize},
  ytick={0,5000,10000,15000,20000},
  y tick label style={font=\scriptsize},
  nodes near coords, nodes near coords style={font=\scriptsize, popdark},
  axis lines=left,
  every axis plot/.append style={fill=poporange, draw=popdark}
]
\addplot coordinates {(Chine,18000) (Inde,3400) (Bresil,1900) (Mexique,1400) (Nigeria,500) (AfrSud,400)};
\end{axis}
\end{tikzpicture}
\legende{\textbf{PIB de six foyers émergents en 2022} (ordres de grandeur, milliards de dollars courants ; source : Banque mondiale). Bresil = Brésil ; Nigeria = Nigéria ; AfrSud = Afrique du Sud.}
\end{popfigure}

\textbf{Lecture.} La barre chinoise (\num{18000} milliards) écrase toutes les autres : elle vaut plus de cinq fois celle de l'Inde et environ treize fois celle du Mexique. Les deux foyers africains (Nigéria, Afrique du Sud) ferment la hiérarchie avec des PIB inférieurs à 500 milliards, comparables à ceux de pays européens moyens.

\textbf{Conclusion.} Le diagramme hiérarchise les foyers en expansion : une superpuissance émergente (Chine), des puissances moyennes (Inde, Brésil, Mexique) et des foyers africains encore modestes, ce qui illustre l'inégale intensité de l'émergence à l'échelle mondiale.
\end{exoresolu}

\courssub{Fiche méthode 6 : Rédiger un commentaire de document ou un paragraphe argumenté}

\begin{methode}[Commenter un document sur les pays émergents]
\begin{enumerate}
\item \textbf{Présenter le document} : nature (carte, tableau, texte), auteur ou source, date, thème.
\item \textbf{Dégager la problématique} : quelle question le document permet-il de traiter ? (ex. : quels facteurs expliquent l'émergence du Mexique ?)
\item \textbf{Analyser} : décris les faits en citant le document (chiffres, localisations), puis explique-les avec tes connaissances (facteurs : main-d'œuvre abondante et bon marché, IDE, matières premières, stabilité politique, ouverture commerciale).
\item \textbf{Nuance} : mentionne toujours les limites (inégalités sociales et spatiales, dépendance aux matières premières, pollution, endettement).
\item \textbf{Conclure} : réponds à la problématique en une phrase et ouvre éventuellement (ex. : du statut d'émergent vers celui de puissance développée ?).
\end{enumerate}
\textbf{Plan type d'un paragraphe argumenté} : idée directrice $\rightarrow$ faits et chiffres $\rightarrow$ explication $\rightarrow$ nuance $\rightarrow$ phrase de transition.
\end{methode}

\begin{exoresolu}[Paragraphe argumenté : les facteurs et les limites de l'émergence mexicaine]
\textbf{Énoncé (type Bac).} « Le Mexique, atelier manufacturier adossé aux États-Unis, reste un émergent inachevé. » Rédigez un paragraphe argumenté d'une quinzaine de lignes discutant cette affirmation.
\tcblower
\textbf{Solution détaillée.}

\emph{Idée directrice.} Le Mexique illustre une émergence tirée par l'industrie manufacturière d'exportation, mais freinée par de lourdes limites internes.

\emph{Faits et explication.} Membre de l'ALENA (accord de libre-échange nord-américain avec les États-Unis et le Canada, entré en vigueur en 1994 et remplacé en 2020 par l'ACEUM, l'Accord Canada--États-Unis--Mexique), le Mexique a attiré des milliers de \cle{maquiladoras}, usines de montage (automobile, électronique, textile) installées près de la frontière nord (Tijuana, Ciudad Juárez, Monterrey) pour profiter d'une main-d'œuvre abondante et bon marché et de la proximité du marché américain. Les IDE des firmes transnationales (General Motors, Volkswagen à Puebla) ont fait du Mexique le premier exportateur manufacturier d'Amérique latine, avec un PIB de l'ordre de \num{1400} milliards de dollars (ordre de grandeur 2022) et une population d'environ 128 millions d'habitants.

\emph{Nuance.} Mais cette émergence demeure inachevée : dépendance presque totale au marché américain (environ 80 \% des exportations mexicaines y sont destinées), faible valeur ajoutée locale (les maquiladoras importent les composants et réexportent), fractures spatiales entre le Nord industrialisé et le Sud rural indien (Chiapas, Oaxaca), insécurité liée au narcotrafic et forte émigration vers les États-Unis.

\emph{Conclusion.} Le Mexique est donc bien un foyer économique en expansion, inséré dans la mondialisation par l'atelier frontalier, mais son émergence, extravertie et inégale, reste à consolider par la montée en gamme industrielle et la réduction des inégalités.
\end{exoresolu}

\begin{attention}[Les 5 erreurs qui coûtent des points au Bac]
\begin{enumerate}
\item \textbf{Confondre les catégories.} Ne mélange pas BRICS, Dragons et NPI : les quatre Dragons sont la Corée du Sud, Taïwan, Hong Kong et Singapour ; les BRICS sont le Brésil, la Russie, l'Inde, la Chine et l'Afrique du Sud. Placer la Chine parmi les Dragons ou Singapour parmi les BRICS est une faute d'élimination.
\item \textbf{Oublier la légende ou le titre du croquis.} Un croquis sans titre, sans flèche du nord et sans légende numérotée perd la moitié de ses points, même si les localisations sont justes. Chaque symbole utilisé doit figurer dans la légende.
\item \textbf{Commenter un chiffre sans unité ni calcul posé.} Écrire « le PIB augmente beaucoup » ne vaut rien ; il faut écrire l'opération : $\frac{\num{3400}-\num{1700}}{\num{1700}}\times 100 = 100\,\%$, puis conclure : « le PIB indien a doublé ». Un taux de variation sans formule explicite n'est pas crédité.
\item \textbf{Confondre PIB et PIB par habitant.} Un PIB total élevé (Inde, Brésil) peut coexister avec un PIB par habitant moyen : c'est précisément la définition d'un pays émergent à revenu intermédiaire. Ne conclus jamais « pays riche » à partir du seul PIB total.
\item \textbf{Rédiger un commentaire sans nuance.} Décrire l'émergence comme un succès total est sanctionné : le correcteur attend les limites (inégalités sociales et spatiales, dépendance aux matières premières ou au marché américain, pollution, endettement). Un paragraphe sans nuance plafonne la note, même s'il est factuellement exact.
\end{enumerate}
\end{attention}

\newpage

\coursec{Exercices}

\courssub{Je vérifie mes connaissances}

\begin{exercice}[QCM : bien distinguer les catégories de pays]
Pour chaque question, une seule réponse est exacte. Entoure la lettre correspondante.

\textbf{1.} Les BRICS regroupent :
\begin{itemize}
\item[a)] le Brésil, la Russie, l'Inde, la Chine et le Sénégal ;
\item[b)] le Brésil, la Russie, l'Inde, la Chine et l'Afrique du Sud ;
\item[c)] le Brésil, le Royaume-Uni, l'Inde, la Chine et le Soudan.
\end{itemize}

\textbf{2.} Les « quatre Dragons » désignent :
\begin{itemize}
\item[a)] la Corée du Sud, Taïwan, Singapour et Hong Kong ;
\item[b)] le Japon, la Chine, la Thaïlande et le Vietnam ;
\item[c)] l'Inde, le Pakistan, l'Indonésie et les Philippines.
\end{itemize}

\textbf{3.} Un PPTE est :
\begin{itemize}
\item[a)] un pays pauvre très endetté, bénéficiant d'initiatives d'allègement de la dette ;
\item[b)] un pays pétrolier à très forte exportation ;
\item[c)] un pays à très forte croissance industrielle.
\end{itemize}

\textbf{4.} Un NPI (nouveau pays industriel) se caractérise surtout par :
\begin{itemize}
\item[a)] une économie fondée exclusivement sur l'agriculture vivrière ;
\item[b)] une industrialisation rapide tournée vers l'exportation de produits manufacturés ;
\item[c)] l'absence totale d'insertion dans les échanges mondiaux.
\end{itemize}
\end{exercice}

\begin{exercice}[Vrai ou faux, en le justifiant]
Réponds par vrai ou faux, puis justifie ta réponse en une ou deux phrases.
\begin{enumerate}
\item Le Nigéria est un pays émergent d'Afrique dont l'atout principal est le pétrole.
\item Les pays émergents ont tous un IDH très élevé, supérieur à celui de la France.
\item Dubaï, aux Émirats arabes unis, illustre la reconversion d'une rente pétrolière vers les services (finance, tourisme, transport aérien).
\item Le Mexique bénéficie de sa proximité avec les États-Unis, notamment à travers l'accord de libre-échange nord-américain (ALENA, devenu ACEUM).
\end{enumerate}
\end{exercice}

\begin{exercice}[Définitions éclair]
Définis en deux ou trois lignes chacun des termes suivants : \cle{pays émergent}, \cle{BRICS}, \cle{pays à revenus intermédiaires}, \cle{foyer économique}.
\end{exercice}

\begin{exercice}[Calculs directs sur des données]
Le tableau ci-dessous donne le PIB de deux pays émergents (ordres de grandeur, sources : Banque mondiale, données arrondies).

\begin{center}
\begin{tabularx}{\linewidth}{|l|X|X|}\hline
\rowcolor{popblueL} \textbf{Pays} & \textbf{PIB 2000 (milliards USD)} & \textbf{PIB 2022 (milliards USD)} \\ \hline
Brésil & environ 650 & environ 1 920 \\ \hline
Inde & environ 470 & environ 3 390 \\ \hline
\end{tabularx}
\end{center}

\begin{enumerate}
\item Calcule le taux de variation du PIB entre 2000 et 2022 pour chacun des deux pays, selon la formule : $\text{taux} = \dfrac{\text{valeur finale} - \text{valeur initiale}}{\text{valeur initiale}} \times 100$.
\item Quel pays a connu la croissance relative la plus forte ? Justifie.
\item Pourquoi ces chiffres traduisent-ils l'idée de « foyers économiques en expansion » ?
\end{enumerate}
\end{exercice}

\courssub{Je m'entraîne}

\begin{exercice}[Localiser et classer]
À partir de la liste suivante : Brésil, Mexique, Inde, Malaisie, Corée du Sud, Arabie saoudite, Émirats arabes unis, Afrique du Sud, Nigéria, Égypte, Maroc, Argentine, Singapour, Taïwan, Hong Kong.
\begin{enumerate}
\item Classe ces pays en trois colonnes : Amérique, Asie et Moyen-Orient, Afrique.
\item Souligne dans ta liste les membres des BRICS.
\item Entoure les « quatre Dragons ».
\item Pour chaque continent, choisis un pays et rédige trois lignes présentant son principal atout économique.
\end{enumerate}
\end{exercice}

\begin{exercice}[Construire un diagramme en bâtons]
Le tableau ci-dessous présente le PIB par habitant et l'IDH des cinq BRICS (ordres de grandeur arrondis, sources : Banque mondiale et PNUD, début des années 2020).

\begin{center}
\begin{tabularx}{\linewidth}{|l|X|X|}\hline
\rowcolor{popblueL} \textbf{Pays} & \textbf{PIB/habitant (USD)} & \textbf{IDH (sur 1)} \\ \hline
Afrique du Sud & environ 6 800 & environ 0,71 \\ \hline
Brésil & environ 8 900 & environ 0,75 \\ \hline
Chine & environ 12 700 & environ 0,77 \\ \hline
Inde & environ 2 400 & environ 0,63 \\ \hline
Russie & environ 15 300 & environ 0,82 \\ \hline
\end{tabularx}
\end{center}

\begin{enumerate}
\item Construis un diagramme en bâtons représentant le PIB par habitant des cinq pays (échelle : 1 cm pour 2 000 USD).
\item Classe les pays du plus riche au moins riche selon le PIB par habitant, puis selon l'IDH.
\item En quatre lignes, montre que la richesse nationale ne signifie pas automatiquement le développement humain.
\end{enumerate}
\end{exercice}

\begin{exercice}[Légender un fond de carte du monde]
Sur un fond de planisphère (fourni par ton professeur ou reproduit à main levée) :
\begin{enumerate}
\item Localise et nomme les cinq BRICS, les quatre Dragons et deux pays du Golfe (Arabie saoudite, Émirats arabes unis).
\item Construis une légende hiérarchisée utilisant trois couleurs : une pour les BRICS, une pour les Dragons, une pour les pays du Golfe.
\item Rédige un titre précis et daté pour ta carte.
\item En trois lignes, dégage la répartition spatiale des foyers économiques en expansion (continents concernés, hémisphères).
\end{enumerate}
\end{exercice}

\begin{exercice}[Étude de cas : le Maroc, un émergent africain]
Document : le Maroc a développé depuis les années 2000 des pôles industriels tournés vers l'exportation : l'automobile autour de Tanger (port de Tanger Med, devenu le premier port de la Méditerranée en capacité de conteneurs) et l'aéronautique près de Casablanca. Le pays mise aussi sur les énergies renouvelables (centrale solaire de Noor à Ouarzazate) et sur le tourisme.
\begin{enumerate}
\item Releve dans le document trois atouts qui expliquent l'émergence du Maroc.
\item Montre en quatre lignes comment le port de Tanger Med insère le Maroc dans la mondialisation.
\item Cite une limite possible à cette émergence (inégalités, dépendance, chômage des jeunes) et justifie ton choix.
\end{enumerate}
\end{exercice}

\courssub{Je me prépare au Bac}

\begin{exercice}[Étude de documents : les BRICS, géants inégalement développés]
\textbf{Document 1.} Quelques indicateurs des BRICS (ordres de grandeur arrondis, sources : Banque mondiale, PNUD, début des années 2020).

\begin{center}
\begin{tabularx}{\linewidth}{|l|X|X|X|}\hline
\rowcolor{popblueL} \textbf{Pays} & \textbf{PIB total (milliards USD)} & \textbf{PIB/habitant (USD)} & \textbf{IDH} \\ \hline
Chine & environ 17 900 & environ 12 700 & environ 0,77 \\ \hline
Inde & environ 3 390 & environ 2 400 & environ 0,63 \\ \hline
Russie & environ 2 240 & environ 15 300 & environ 0,82 \\ \hline
Brésil & environ 1 920 & environ 8 900 & environ 0,75 \\ \hline
Afrique du Sud & environ 400 & environ 6 800 & environ 0,71 \\ \hline
\end{tabularx}
\end{center}

\textbf{Document 2.} Extrait de texte : « Réunis en sommets depuis 2009, les BRICS pèsent environ 40 \% de la population mondiale et plus du quart du PIB mondial. Ils revendiquent une plus grande place dans la gouvernance économique internationale et ont créé en 2015 la Nouvelle Banque de développement, dont le siège est à Shanghai. »

\begin{enumerate}
\item Définis les sigles PIB et IDH. (2 pts)
\item À partir du document 1, construis un diagramme en bâtons du PIB total des cinq pays (échelle : 1 cm pour 2 000 milliards USD). (4 pts)
\item Compare le classement des pays selon le PIB total et selon le PIB par habitant. Que constates-tu pour l'Inde et la Russie ? (4 pts)
\item En une dizaine de lignes, montre à partir des deux documents que les BRICS sont de nouvelles puissances économiques, mais que leur développement reste inégal et inachevé. (10 pts)
\end{enumerate}
\end{exercice}

\begin{exercice}[Croquis : les foyers économiques en expansion en Afrique]
Sur la base de tes connaissances et d'un fond de carte de l'Afrique :
\begin{enumerate}
\item Construis un croquis intitulé « Les foyers économiques en expansion en Afrique ». (8 pts)
\item Tu y feras figurer obligatoirement : quatre foyers économiques (Afrique du Sud, Nigéria, Égypte, Maroc), deux ports majeurs (Tanger Med, Durban ou Lagos), un corridor d'échanges (par exemple l'axe Gauteng--Johannesburg ou le corridor de Suez) et deux flux (exportations pétrolières nigérianes, flux commerciaux vers l'Europe ou l'Asie). (8 pts)
\item Rédige une légende organisée en trois rubriques maximum, avec des figurés simples et propres. (2 pts)
\item N'oublie ni l'orientation (flèche du nord) ni l'échelle indicative. (2 pts)
\end{enumerate}
\end{exercice}

\begin{exercice}[Dissertation guidée : émergence et mondialisation]
Sujet : \emph{« Les pays émergents : de nouvelles puissances qui bouleversent la mondialisation ? »}

Consignes : rédige une dissertation organisée en suivant le plan imposé ci-dessous.

\begin{enumerate}
\item \textbf{Introduction} (environ 10 lignes) : amène le sujet à partir d'un fait concret (par exemple la place de la Chine ou du Brésil dans le commerce mondial), définis l'expression « pays émergent », pose une problématique et annonce le plan. (4 pts)
\item \textbf{Première partie : les manifestations de l'émergence} (environ 15 lignes) : croissance rapide, industrialisation et exportations (Dragons, Mexique), poids démographique et commercial des BRICS, rôle des métropoles et des ports (Shanghai, Dubaï, São Paulo, Johannesburg), création d'institutions comme la Nouvelle Banque de développement. Appuie chaque idée sur au moins un exemple précis. (8 pts)
\item \textbf{Deuxième partie : les limites et les inégalités de l'émergence} (environ 15 lignes) : IDH parfois moyens, fortes inégalités sociales et spatiales (favelas au Brésil, bidonvilles en Inde), dépendances (matières premières au Nigéria, capitaux étrangers), vulnérabilités environnementales. (8 pts)
\item \textbf{Conclusion} (environ 8 lignes) : réponds clairement à la problématique et ouvre sur la place possible de l'Afrique, dont le Cameroun, dans cette mondialisation recomposée. (2 pts)
\end{enumerate}
\end{exercice}

\newpage

\coursec{Corrigés des exercices}

\begin{corrige}[Exercice 1 — QCM : bien distinguer les catégories de pays]
\textbf{1. Réponse b).} Les BRICS regroupent le \cle{Brésil}, la \cle{Russie}, l'\cle{Inde}, la \cle{Chine} et l'\cle{Afrique du Sud} (le « S » final correspond à \emph{South Africa}, intégrée en 2010). Le Sénégal et le Soudan n'en font pas partie.

\textbf{2. Réponse a).} Les « quatre Dragons » sont la \cle{Corée du Sud}, \cle{Taïwan}, \cle{Singapour} et \cle{Hong Kong}, territoires d'Asie de l'Est industrialisés dès les années 1960--1980 et tournés vers l'exportation. Le Japon est une puissance ancienne (membre du G7), il n'est pas un Dragon.

\textbf{3. Réponse a).} PPTE signifie \cle{pays pauvre très endetté} : il s'agit de pays à faible revenu dont la dette extérieure est insoutenable et qui bénéficient d'initiatives d'allègement pilotées par le FMI et la Banque mondiale (initiative PPTE lancée en 1996). Le Cameroun en a bénéficié (point d'achèvement atteint en 2006).

\textbf{4. Réponse b).} Un \cle{NPI} (nouveau pays industriel) s'est industrialisé rapidement en s'insérant dans la mondialisation par l'\textbf{exportation de produits manufacturés} (textile, électronique, automobile). Les réponses a) et c) décrivent au contraire des économies marginalisées.

\textbf{Point de vigilance :} ne pas confondre NPI (catégorie ancienne, années 1970--1980), pays émergents (catégorie plus récente et plus large) et BRICS (club diplomatique de cinq grands émergents). Un pays peut cumuler plusieurs étiquettes (la Corée du Sud fut un Dragon, donc un NPI, et est aujourd'hui souvent classée parmi les pays développés).
\end{corrige}

\begin{corrige}[Exercice 2 — Vrai ou faux, en le justifiant]
\textbf{1. Vrai.} Le Nigéria est le premier producteur de pétrole d'Afrique subsaharienne (en concurrence avec l'Angola) ; les hydrocarbures fournissent l'essentiel de ses recettes d'exportation et une large part des recettes de l'État. C'est aussi le pays africain le plus peuplé (plus de 200 millions d'habitants), ce qui renforce son poids économique.

\textbf{2. Faux.} Les pays émergents ont des IDH \emph{intermédiaires}, souvent compris entre 0,6 et 0,8 environ, inférieurs à celui de la France (environ 0,90). L'Inde, par exemple, a un IDH voisin de 0,63. La croissance économique ne s'accompagne pas automatiquement d'un développement humain achevé.

\textbf{3. Vrai.} Dubaï a investi sa rente pétrolière dans la finance, l'immobilier, le tourisme de luxe et le transport aérien (compagnie Emirates, aéroport international de Dubaï parmi les plus fréquentés du monde), anticipant l'épuisement des hydrocarbures.

\textbf{4. Vrai.} Le Mexique est lié aux États-Unis et au Canada par l'ALENA (1994), remplacé en 2020 par l'ACEUM. Ses usines d'assemblage (\emph{maquiladoras}) exportent massivement vers le marché nord-américain : la proximité avec la première puissance mondiale est son principal atout.

\textbf{Point de vigilance :} une justification d'une phrase doit contenir un \textbf{fait précis} (chiffre, nom, date) et non une simple reformulation de l'énoncé.
\end{corrige}

\begin{corrige}[Exercice 3 — Définitions éclair]
\cle{Pays émergent} : pays en voie de développement dont la croissance économique rapide et durable, l'industrialisation et l'insertion croissante dans les échanges mondiaux le rapprochent progressivement du niveau des pays développés, sans que le développement humain soit achevé.

\cle{BRICS} : acronyme désignant le groupe informel formé par le Brésil, la Russie, l'Inde, la Chine et l'Afrique du Sud, réunis en sommets depuis 2009 ; ils pèsent environ 40 \% de la population mondiale et plus du quart du PIB mondial, et revendiquent une gouvernance mondiale rééquilibrée (Nouvelle Banque de développement créée en 2015, siège à Shanghai).

\cle{Pays à revenus intermédiaires} : catégorie de la Banque mondiale désignant les pays dont le revenu national brut par habitant se situe entre les pays à faible revenu et les pays à haut revenu (seuils révisés chaque année) ; la plupart des pays émergents appartiennent à cette catégorie.

\cle{Foyer économique} : espace de dimension variable (ville, région, pays) qui concentre des activités productives, des richesses, des flux et des fonctions de commandement, et qui rayonne sur son arrière-pays. Dans cette leçon, les foyers en expansion sont les États et métropoles émergents dont le poids dans l'économie mondiale croît rapidement.
\end{corrige}

\begin{corrige}[Exercice 4 — Calculs directs sur des données]
\textbf{1. Taux de variation du PIB entre 2000 et 2022.}

Pour le Brésil :
\[
\text{taux}_{\text{Brésil}} = \dfrac{1\,920 - 650}{650} \times 100 = \dfrac{1\,270}{650} \times 100 \approx 195,4\,\%
\]
Le PIB brésilien a été multiplié par environ 2,95 (il a presque triplé).

Pour l'Inde :
\[
\text{taux}_{\text{Inde}} = \dfrac{3\,390 - 470}{470} \times 100 = \dfrac{2\,920}{470} \times 100 \approx 621,3\,\%
\]
Le PIB indien a été multiplié par environ 7,2.

\textbf{2. Croissance relative la plus forte.} C'est l'\textbf{Inde} : $621,3\,\% > 195,4\,\%$. Partie d'un niveau plus faible (470 milliards USD contre 650), elle a rattrapé puis dépassé le Brésil : en 2022, son PIB (environ 3\,390 milliards USD) est nettement supérieur à celui du Brésil (environ 1\,920 milliards USD).

\textbf{3. Lien avec la notion de « foyers en expansion ».} Une croissance de cette ampleur sur vingt-deux ans montre que ces pays accumulent rapidement des richesses, s'industrialisent et pèsent de plus en plus lourd dans le PIB et le commerce mondiaux : ce ne sont plus des périphéries passives, mais des pôles dynamiques qui attirent capitaux, usines et flux, et qui concurrencent les puissances anciennes. C'est précisément la définition d'un foyer économique \emph{en expansion}.

\textbf{Point de vigilance :} bien distinguer le \textbf{taux de variation} (en \%, mesure la croissance \emph{relative}) de la \textbf{variation absolue} (en milliards USD). Erreur fréquente : oublier de multiplier par 100, ou diviser par la valeur finale au lieu de la valeur initiale.
\end{corrige}

\begin{corrige}[Exercice 5 — Localiser et classer]
\textbf{1. Classement par continent.}

\begin{center}
\begin{tabularx}{\linewidth}{|X|X|X|}\hline
\rowcolor{popblueL} \textbf{Amérique} & \textbf{Asie et Moyen-Orient} & \textbf{Afrique} \\ \hline
Brésil, Mexique, Argentine & Inde, Malaisie, Corée du Sud, Arabie saoudite, Émirats arabes unis, Singapour, Taïwan, Hong Kong & Afrique du Sud, Nigéria, Égypte, Maroc \\ \hline
\end{tabularx}
\end{center}

\textbf{2. Membres des BRICS (à souligner) :} Brésil, Inde, Afrique du Sud. (La Chine et la Russie ne figurent pas dans la liste proposée.)

\textbf{3. Les quatre Dragons (à entourer) :} Corée du Sud, Singapour, Taïwan, Hong Kong.

\textbf{4. Exemples d'atouts (un par continent).}
\begin{itemize}
\item \textbf{Amérique — Brésil :} puissance agricole et minière de premier plan (soja, café, minerai de fer), dotée d'une industrie diversifiée (aéronautique avec Embraer) et d'un immense marché intérieur de plus de 200 millions d'habitants.
\item \textbf{Asie — Corée du Sud :} championne des industries de haute technologie (électronique avec Samsung, automobile avec Hyundai), portée par une population très qualifiée et des exportations mondialisées.
\item \textbf{Afrique — Nigéria :} premier producteur de pétrole d'Afrique subsaharienne et pays le plus peuplé du continent, il dispose d'une rente pétrolière considérable et d'un marché intérieur en pleine croissance.
\end{itemize}

\textbf{Point de vigilance :} Hong Kong est une région administrative spéciale de la Chine et non un État souverain ; au lycée, on l'admet dans la liste des Dragons pour des raisons historiques et économiques.
\end{corrige}

\begin{corrige}[Exercice 6 — Construire un diagramme en bâtons]
\textbf{1. Hauteurs des bâtons} (échelle : 1 cm pour 2\,000 USD, donc hauteur $=$ PIB/hab $\div$ 2\,000) :

\begin{center}
\begin{tabularx}{\linewidth}{|l|X|X|}\hline
\rowcolor{popblueL} \textbf{Pays} & \textbf{PIB/habitant (USD)} & \textbf{Hauteur du bâton (cm)} \\ \hline
Afrique du Sud & 6\,800 & $6\,800/2\,000 = 3,4$ \\ \hline
Brésil & 8\,900 & $8\,900/2\,000 = 4,45$ \\ \hline
Chine & 12\,700 & $12\,700/2\,000 = 6,35$ \\ \hline
Inde & 2\,400 & $2\,400/2\,000 = 1,2$ \\ \hline
Russie & 15\,300 & $15\,300/2\,000 = 7,65$ \\ \hline
\end{tabularx}
\end{center}

Modèle de diagramme attendu (à reproduire sur ta copie avec les hauteurs calculées) :

\begin{popfigure}
\begin{center}
\begin{tikzpicture}[scale=0.9]
\draw[->,thick] (0,0) -- (10.5,0) node[below left]{\footnotesize Pays};
\draw[->,thick] (0,0) -- (0,8.6) node[above left]{\footnotesize PIB/hab (USD)};
\foreach \y/\lab in {2/4000,4/8000,6/12000,7.65/15300}{\draw (0,\y) node[left]{\footnotesize \lab} -- (-0.08,\y);}
\fill[popgold] (0.7,0) rectangle (1.9,3.4);
\fill[popgreen] (2.7,0) rectangle (3.9,4.45);
\fill[poporange] (4.7,0) rectangle (5.9,6.35);
\fill[poppink] (6.7,0) rectangle (7.9,1.2);
\fill[popblue] (8.7,0) rectangle (9.9,7.65);
\node[below,font=\footnotesize] at (1.3,0) {Afr. Sud};
\node[below,font=\footnotesize] at (3.3,0) {Brésil};
\node[below,font=\footnotesize] at (5.3,0) {Chine};
\node[below,font=\footnotesize] at (7.3,0) {Inde};
\node[below,font=\footnotesize] at (9.3,0) {Russie};
\end{tikzpicture}
\end{center}
\legende{Diagramme en bâtons du PIB par habitant des BRICS (ordres de grandeur, début des années 2020). Échelle : 1 cm pour 2\,000 USD. Sources : Banque mondiale, données arrondies.}
\end{popfigure}

\textbf{2. Classements.}
\begin{itemize}
\item Selon le PIB par habitant (du plus riche au moins riche) : Russie (15\,300) $>$ Chine (12\,700) $>$ Brésil (8\,900) $>$ Afrique du Sud (6\,800) $>$ Inde (2\,400).
\item Selon l'IDH : Russie (0,82) $>$ Chine (0,77) $>$ Brésil (0,75) $>$ Afrique du Sud (0,71) $>$ Inde (0,63).
\end{itemize}

\textbf{3. Richesse et développement humain.} Les deux classements se ressemblent ici, mais les écarts diffèrent : l'Inde, deuxième PIB \emph{total} du groupe, reste dernière en richesse par habitant et en IDH, car sa richesse est diluée dans une population immense et très inégalement partagée. Inversement, l'Afrique du Sud affiche un PIB par habitant honorable mais un IDH limité par de fortes inégalités et un chômage massif. Le PIB mesure la production de richesse, l'IDH intègre aussi la santé et l'éducation : une nation riche peut donc laisser une partie de sa population à l'écart du développement humain.

\textbf{Point de vigilance :} sur la copie, chaque bâton doit être de même largeur, séparé du suivant, avec les axes gradués, l'unité indiquée et un titre daté. Un diagramme sans échelle ni unité perd des points.
\end{corrige}

\begin{corrige}[Exercice 7 — Légender un fond de carte du monde]
\textbf{1. Localisations attendues.}
\begin{itemize}
\item \textbf{BRICS :} Brésil (est de l'Amérique du Sud), Russie (nord de l'Eurasie), Inde (péninsule indienne), Chine (Asie de l'Est), Afrique du Sud (extrême sud de l'Afrique).
\item \textbf{Dragons :} Corée du Sud (péninsule coréenne), Taïwan (île au large de la Chine), Hong Kong (côte sud de la Chine), Singapour (pointe de la péninsule malaise).
\item \textbf{Golfe :} Arabie saoudite (péninsule arabique), Émirats arabes unis (côte sud du golfe Persique).
\end{itemize}

\textbf{2. Légende hiérarchisée (trois couleurs).}
\begin{itemize}
\item Figuré 1 (surface colorée, par exemple orange) : les BRICS, grands pays émergents continentaux.
\item Figuré 2 (points ou petites surfaces, par exemple bleu) : les quatre Dragons, petits territoires industriels d'Asie de l'Est.
\item Figuré 3 (surface, par exemple vert) : les monarchies pétrolières du Golfe.
\end{itemize}

\textbf{3. Titre précis et daté :} « Les foyers économiques en expansion dans le monde : BRICS, Dragons et monarchies du Golfe (début des années 2020) ».

\textbf{4. Répartition spatiale.} Les foyers en expansion se répartissent sur trois continents : l'Amérique latine (Brésil, Mexique, Argentine), l'Asie--Moyen-Orient (Chine, Inde, Dragons, Golfe) et l'Afrique (Afrique du Sud, Nigéria, Égypte, Maroc). Ils appartiennent majoritairement à l'\textbf{hémisphère Sud} ou aux marges du Nord, ce qui traduit un basculement partiel du centre de gravité de l'économie mondiale vers les « Suds ». Aucun n'est situé en Europe occidentale, cœur historique de la puissance mondiale.

\textbf{Point de vigilance :} chaque figuré de la légende doit être \emph{exactement reproduit} sur la carte (même couleur, même forme) ; un nom placé sans figuré correspondant dans la légende est une erreur classique.
\end{corrige}

\begin{corrige}[Exercice 8 — Étude de cas : le Maroc, un émergent africain]
\textbf{1. Trois atouts relevés dans le document :}
\begin{itemize}
\item des pôles industriels d'exportation : l'automobile autour de Tanger et l'aéronautique près de Casablanca ;
\item un équipement portuaire de rang mondial : Tanger Med, premier port de Méditerranée en capacité de conteneurs ;
\item une stratégie de diversification : énergies renouvelables (centrale solaire de Noor à Ouarzazate) et tourisme.
\end{itemize}

\textbf{2. Tanger Med et la mondialisation.} Situé sur le détroit de Gibraltar, à la croisée de l'Europe et de l'Afrique, Tanger Med accueille les grandes lignes maritimes de conteneurs qui relient l'Asie, l'Europe et l'Amérique. Il permet aux usines marocaines (automobile, textile) d'exporter rapidement vers l'Union européenne, toute proche, et attire les investissements étrangers (constructeurs automobiles français notamment). Le port fait ainsi du Maroc une plateforme logistique et industrielle insérée dans les chaînes de valeur mondiales.

\textbf{3. Une limite possible.} Les inégalités sociales et spatiales : la croissance profite surtout au littoral industrialisé (axe Tanger--Casablanca), tandis que l'intérieur rural reste pauvre et que le chômage touche fortement les jeunes diplômés. On peut aussi citer la dépendance aux capitaux et aux marchés européens, qui rend l'économie vulnérable aux crises extérieures.

\textbf{Point de vigilance :} à la question 1, il faut \emph{relever} des éléments du document (citation ou reformulation fidèle), et non mobiliser ses connaissances générales sur le Maroc.
\end{corrige}

\begin{corrige}[Exercice 9 — Étude de documents : les BRICS, géants inégalement développés]
\textbf{Barème indicatif : 20 points (2 + 4 + 4 + 10).}

\textbf{1. Définitions (2 pts, 1 pt par sigle).}
\begin{itemize}
\item \cle{PIB} (produit intérieur brut) : valeur totale des biens et services produits sur le territoire d'un pays en une année ; il mesure la richesse produite.
\item \cle{IDH} (indice de développement humain) : indicateur synthétique du PNUD, compris entre 0 et 1, combinant le revenu par habitant, l'espérance de vie et le niveau d'éducation.
\end{itemize}

\textbf{2. Diagramme en bâtons du PIB total (4 pts).} Échelle : 1 cm pour 2\,000 milliards USD, d'où les hauteurs : Chine $17\,900/2\,000 = 8,95$ cm ; Inde $3\,390/2\,000 \approx 1,7$ cm ; Russie $2\,240/2\,000 = 1,12$ cm ; Brésil $1\,920/2\,000 = 0,96$ cm ; Afrique du Sud $400/2\,000 = 0,2$ cm. Attribution des points : axes gradués et légendés (1 pt), hauteurs exactes (2 pts), titre daté (1 pt).

\begin{popfigure}
\begin{center}
\begin{tikzpicture}[scale=0.9]
\draw[->,thick] (0,0) -- (10.5,0) node[below left]{\footnotesize Pays};
\draw[->,thick] (0,0) -- (0,9.8) node[above left]{\footnotesize PIB (mds USD)};
\foreach \y/\lab in {2/4000,4/8000,6/12000,8/16000}{\draw (0,\y) node[left]{\footnotesize \lab} -- (-0.08,\y);}
\fill[poporange] (0.7,0) rectangle (1.9,8.95);
\fill[poppink] (2.7,0) rectangle (3.9,1.7);
\fill[popblue] (4.7,0) rectangle (5.9,1.12);
\fill[popgreen] (6.7,0) rectangle (7.9,0.96);
\fill[popgold] (8.7,0) rectangle (9.9,0.2);
\node[below,font=\footnotesize] at (1.3,0) {Chine};
\node[below,font=\footnotesize] at (3.3,0) {Inde};
\node[below,font=\footnotesize] at (5.3,0) {Russie};
\node[below,font=\footnotesize] at (7.3,0) {Brésil};
\node[below,font=\footnotesize] at (9.3,0) {Afr. Sud};
\end{tikzpicture}
\end{center}
\legende{PIB total des BRICS (ordres de grandeur, début des années 2020). Échelle : 1 cm pour 2\,000 milliards USD. Source : Banque mondiale, données arrondies.}
\end{popfigure}

\textbf{3. Comparaison des classements (4 pts).} Selon le PIB total : Chine $\gg$ Inde $>$ Russie $>$ Brésil $>$ Afrique du Sud. Selon le PIB par habitant : Russie (15\,300) $>$ Chine (12\,700) $>$ Brésil (8\,900) $>$ Afrique du Sud (6\,800) $>$ Inde (2\,400). Deux constats spectaculaires : l'\textbf{Inde}, deuxième du groupe en PIB total, devient \emph{dernière} en richesse par habitant, car sa production se répartit sur plus de 1,4 milliard d'habitants ; la \textbf{Russie}, seulement troisième en PIB total, devient \emph{première} en PIB par habitant grâce à une population bien moindre (environ 145 millions). Le PIB total mesure la puissance, le PIB par habitant mesure le niveau de vie moyen : ce sont deux lectures différentes de la richesse.

\textbf{4. Paragraphe argumenté (10 pts).} Éléments attendus :
\begin{itemize}
\item \emph{Nouvelles puissances} : poids démographique (environ 40 \% de l'humanité) et économique (plus du quart du PIB mondial) ; la Chine est la deuxième économie mondiale (environ 17\,900 milliards USD) ; institution propre (Nouvelle Banque de développement, 2015, siège à Shanghai) ; revendication d'une gouvernance mondiale rééquilibrée.
\item \emph{Développement inégal} : écart considérable entre la Chine (17\,900 milliards USD) et l'Afrique du Sud (400 milliards USD) ; PIB par habitant de 2\,400 USD en Inde contre 15\,300 USD en Russie.
\item \emph{Développement inachevé} : IDH compris entre 0,63 et 0,82, inférieurs à ceux des pays développés (environ 0,90) ; inégalités sociales internes fortes.
\end{itemize}
Barème : définitions implicites et mobilisation des deux documents (4 pts), exactitude des chiffres cités (3 pts), organisation et qualité de l'expression (3 pts).

\textbf{Points de vigilance :} citer les deux documents explicitement (« d'après le document 2\ldots ») ; ne pas écrire que les BRICS sont « développés » : ils sont \emph{émergents}, c'est-à-dire en cours de rattrapage ; vérifier que chaque bâton respecte l'échelle annoncée.
\end{corrige}

\begin{corrige}[Exercice 10 — Croquis : les foyers économiques en expansion en Afrique]
\textbf{Barème indicatif : 20 points} (construction et exactitude du fond : 8 pts ; éléments obligatoires : 8 pts ; légende : 2 pts ; orientation et échelle : 2 pts).

\textbf{1--2. Organisation attendue du croquis.}
\begin{itemize}
\item \textbf{Fond :} contour simplifié de l'Afrique, propre, au crayon ; les quatre foyers placés correctement : Maroc (nord-ouest), Égypte (nord-est, sur le Nil et Suez), Nigéria (golfe de Guinée, ouest), Afrique du Sud (extrême sud).
\item \textbf{Ports :} Tanger Med (détroit de Gibraltar), Durban (côte est sud-africaine) ou Lagos (Nigéria), figurés par un symbole d'ancre ou un point noir nommé.
\item \textbf{Corridor :} l'axe Gauteng--Johannesburg (trait plein) ou le corridor maritime de Suez (pointillés le long de la mer Rouge).
\item \textbf{Flux :} flèches d'épaisseur proportionnelle : exportations pétrolières du golfe de Guinée vers l'Europe et l'Amérique ; flux commerciaux de Tanger Med vers l'Europe et de Durban vers l'Asie.
\end{itemize}

Modèle de croquis (à adapter sur fond de carte réel) :

\begin{popfigure}
\includegraphics[width=\linewidth,height=11cm,keepaspectratio]{N14/afrique_pol.jpg}
\legende{Fond de carte pour le croquis modèle (carte politique de l'Afrique, en anglais). Le croquis attendu y situe : 1. les foyers émergents (Nigéria, Égypte, Afrique du Sud, Maroc) ; 2. les ports majeurs (Tanger Med, Alexandrie et Port-Saïd, Lagos, Durban, Douala) ; 3. les corridors d'échanges ; 4. les flux commerciaux et pétroliers (flèches). Seuls les États, les capitales et les grandes villes sont déjà représentés ; les figurés 1 à 4 sont à ajouter, avec une flèche du nord et une échelle. \\ Source : Wikimedia Commons, Central Intelligence Agency (États-Unis), domaine public.}
\end{popfigure}

\textbf{3. Légende en trois rubriques maximum :} par exemple « Foyers et métropoles » (disques, points), « Axes et corridors » (traits pleins et pointillés), « Flux » (flèches). Chaque figuré doit être dessiné dans la légende \emph{à l'identique} de la carte.

\textbf{4. Orientation et échelle :} flèche du nord en haut à droite ; échelle indicative (un segment coté, par exemple 2\,000 km).

\textbf{Points de vigilance :} le croquis se fait \textbf{au crayon}, sans rature, avec un titre souligné au-dessus ; les noms s'écrivent à l'horizontale, sans flèches de rappel croisées ; ne pas surcharger : quatre foyers, deux ports, un corridor, deux flux suffisent ; l'épaisseur des flèches doit être visiblement différente selon l'importance du flux.
\end{corrige}

\begin{corrige}[Exercice 11 — Dissertation guidée : émergence et mondialisation]
\textbf{Barème indicatif : 22 points} (introduction 4 pts ; première partie 8 pts ; deuxième partie 8 pts ; conclusion 2 pts).

\textbf{1. Introduction (modèle rédigé).} En 2022, la Chine est la deuxième économie mondiale et le premier exportateur de marchandises de la planète : un statut impensable il y a quarante ans. Comme elle, le Brésil, l'Inde, le Mexique ou l'Afrique du Sud connaissent une croissance rapide qui les rapproche des puissances anciennes. On appelle \cle{pays émergents} les pays en voie de développement dont l'industrialisation et l'insertion dans les échanges mondiaux progressent vite, sans que le développement humain soit achevé. Dès lors, dans quelle mesure ces pays émergents constituent-ils de nouvelles puissances qui bouleversent la mondialisation ? Nous verrons d'abord les manifestations de cette émergence, puis ses limites et ses inégalités.

\textbf{2. Première partie : les manifestations de l'émergence (éléments attendus).}
\begin{itemize}
\item \emph{Croissance et industrialisation tournées vers l'exportation} : les quatre Dragons (Corée du Sud, Taïwan, Singapour, Hong Kong) ont bâti leur richesse sur les produits manufacturés exportés ; le Mexique assemble dans ses \emph{maquiladoras} pour le marché nord-américain (ALENA puis ACEUM).
\item \emph{Poids démographique et commercial des BRICS} : environ 40 \% de la population mondiale, plus du quart du PIB mondial ; la Chine atteint environ 17\,900 milliards USD de PIB.
\item \emph{Rôle des métropoles et des ports} : Shanghai, premier port de conteneurs du monde ; Dubaï, hub aérien et financier reconverti de la rente pétrolière ; São Paulo, capitale économique du Brésil ; Johannesburg, cœur industriel et financier de l'Afrique australe ; Tanger Med, premier port de Méditerranée en capacité.
\item \emph{Institutions nouvelles} : sommets des BRICS depuis 2009, Nouvelle Banque de développement (2015, siège à Shanghai), qui concurrencent les institutions dirigées par les puissances anciennes.
\end{itemize}
Barème : quatre idées au moins, chacune appuyée sur un exemple précis et chiffré (2 pts par idée illustrée).

\textbf{3. Deuxième partie : les limites et les inégalités (éléments attendus).}
\begin{itemize}
\item \emph{Développement humain inachevé} : IDH de l'Inde voisin de 0,63, loin des 0,90 des pays développés ; PIB par habitant indien d'environ 2\,400 USD.
\item \emph{Inégalités sociales et spatiales} : favelas de Rio et de São Paulo face aux quartiers d'affaires ; bidonvilles de Mumbai ; chômage massif des jeunes en Afrique du Sud et au Maroc ; littoraux dynamiques contre intérieurs ruraux délaissés (Brésil, Maroc).
\item \emph{Dépendances} : le Nigéria dépend de ses exportations de pétrole (volatilité des cours) ; beaucoup d'émergents dépendent des capitaux et des technologies étrangères et des marchés du Nord (Maroc et Union européenne).
\item \emph{Vulnérabilités environnementales} : déforestation amazonienne liée à l'expansion agricole brésilienne, pollution industrielle en Chine, pression sur l'eau et les sols.
\end{itemize}
Barème : quatre idées au moins, chacune illustrée (2 pts par idée illustrée).

\textbf{4. Conclusion (modèle rédigé).} Les pays émergents sont bien devenus des puissances économiques qui recomposent la mondialisation : le centre de gravité du commerce mondial bascule vers l'Asie et les Suds, et les BRICS contestent l'ordre économique hérité de 1945. Mais cette puissance demeure incomplète, inégale entre pays et à l'intérieur de chacun. L'enjeu, pour l'Afrique et pour le Cameroun, est de transformer cette recomposition en opportunité : en industrialisant ses matières premières, en développant ses corridors (port de Kribi) et en s'appuyant sur la Zone de libre-échange continentale africaine, le continent peut à son tour devenir un foyer en expansion.

\textbf{Points de vigilance :} la problématique doit être une \emph{question} ; chaque partie doit s'appuyer sur au moins trois exemples nommés et chiffrés ; éviter le catalogue sans articulation (utiliser « d'abord\ldots ensuite\ldots enfin », « cependant », « en revanche ») ; la conclusion doit répondre explicitement à la question posée, puis ouvrir sans introduire d'argument nouveau.
\end{corrige}

\newpage

\coursec{Fiche bilan}

\courssub{Carte mentale de synthèse}
\begin{popfigure}
\begin{tikzpicture}[
    mindmap,
    grow cyclic,
    every node/.style={concept, execute at begin node=\hskip0pt},
    concept color=popblue,
    level 1/.append style={level distance=4.5cm, sibling angle=360/7},
    level 2/.append style={level distance=3.2cm, sibling angle=360/4},
    text=white,
    font=\small\bfseries
]
\node[concept, scale=1.2, align=center]{Foyers économiques\\ en expansion}
    child[concept color=poporange] {node[concept, align=center]{1. Définition\\ \& Concepts}
        child {node[concept, scale=0.7] {Pays émergents}}
        child {node[concept, scale=0.7] {NPI (Nouv. Pays Industriels)}}
        child {node[concept, scale=0.7] {BRICS / BRICS+}}
        child {node[concept, scale=0.7] {PIB/hab., IDH, IDE}}
    }
    child[concept color=popgreen] {node[concept] {2. Amérique}
        child {node[concept, scale=0.7] {Brésil (agro, aéro)}}
        child {node[concept, scale=0.7] {Mexique (maquila, auto)}}
        child {node[concept, scale=0.7] {Argentine (agro, énergie)}}
    }
    child[concept color=poppurple] {node[concept, align=center]{3. Asie\\ \& Moyen-Orient}
        child {node[concept, scale=0.7] {Chine (atelier monde)}}
        child {node[concept, scale=0.7] {Inde (services, tech)}}
        child {node[concept, scale=0.7] {4 Dragons (Corée, Taïwan, HK, Singapour)}}
        child {node[concept, scale=0.7] {EAU / Arabie (pétrole, finance, logistique)}}
    }
    child[concept color=poppink] {node[concept] {4. Afrique}
        child {node[concept, scale=0.7] {Afrique du Sud (mines, finance, auto)}}
        child {node[concept, scale=0.7] {Nigéria (pétrole, démographie, tech)}}
        child {node[concept, scale=0.7] {Égypte (Suez, gaz, tourisme)}}
        child {node[concept, scale=0.7] {Maroc (auto, aéro, phosphate, ZES)}}
    }
    child[concept color=popgold] {node[concept, align=center]{5. Facteurs\\ d'émergence}
        child {node[concept, scale=0.7] {Ressources naturelles}}
        child {node[concept, scale=0.7] {Dividende démographique}}
        child {node[concept, scale=0.7] {Réformes libérales \& ZES}}
        child {node[concept, scale=0.7] {Attractivité IDE}}
    }
    child[concept color=popdark] {node[concept] {6. Manifestations}
        child {node[concept, scale=0.7] {Croissance PIB > 4-5\%}}
        child {node[concept, scale=0.7] {Urbanisation / Métropolisation}}
        child {node[concept, scale=0.7] {Montée en gamme exportations}}
        child {node[concept, scale=0.7] {Multinationales locales}}
    }
    child[concept color=popmuted] {node[concept, align=center]{7. Limites\\ \& Défis}
        child {node[concept, scale=0.7] {Inégalités sociales fortes}}
        child {node[concept, scale=0.7] {Dépendance extraversion}}
        child {node[concept, scale=0.7] {Pression environnementale}}
        child {node[concept, scale=0.7] {Instabilité politique / gouvernance}}
    };
\end{tikzpicture}
\legende{Carte mentale récapitulative de la leçon 14 : les 7 piliers pour comprendre les foyers économiques en expansion.}
\end{popfigure}

\courssub{Bilan structuré}
\begin{bilan}

\textbf{1. Définitions clés à connaître par cœur}
\begin{itemize}
    \item \cle{Pays émergent} : Pays en développement dont la croissance économique soutenue (généralement $> 4\%$/an), l'industrialisation rapide et l'intégration aux marchés financiers mondiaux le rapprochent des pays développés.
    \item \cle{NPI (Nouveaux Pays Industriels)} : 1\up{re} vague (4 Dragons asiatiques : Corée du Sud, Taïwan, Hong Kong, Singapour) ; 2\up{e} vague (Malaisie, Thaïlande, Indonésie, Vietnam) ; 3\up{e} vague (Chine, Inde).
    \item \cle{BRICS} : Groupe géopolitique et économique (Brésil, Russie, Inde, Chine, Afrique du Sud) élargi en 2024 (BRICS+ : Égypte, Émirats, Éthiopie, Iran, Arabie Saoudite). Représente $\approx 45\%$ population mondiale, $\approx 35\%$ PIB mondial (PPA).
    \item \cle{IDE (Investissements Directs à l'Étranger)} : Investissement impliquant une relation durable et un contrôle significatif (\textgreater 10\% capital) d'une entreprise dans un pays par une entité d'un autre pays. Moteur principal de l'émergence.
    \item \cle{ZES (Zones Économiques Spéciales)} : Zones délimitées offrant avantages fiscaux, douaniers, administratifs pour attirer les IDE (ex: Shenzhen Chine, Tanger Med Maroc, Lagos Nigéria).
\end{itemize}

\textbf{2. Repères chiffrés (ordres de grandeur 2023-2024)}
\begin{center}
\begin{tabularx}{\linewidth}{|l|X|X|X|}\hline
\rowcolor{popblueL}
\textbf{Pays / Groupe} & \textbf{PIB nominal (Mds \$)} & \textbf{PIB/hab (PPA, \$)} & \textbf{Part export manuf.} \\ \hline
Chine & $\sim 18\,000$ & $\sim 23\,000$ & $\sim 95\%$ \\ \hline
Inde & $\sim 3\,500$ & $\sim 9\,000$ & $\sim 70\%$ \\ \hline
Brésil & $\sim 2\,100$ & $\sim 17\,000$ & $\sim 60\%$ \\ \hline
Afrique du Sud & $\sim 400$ & $\sim 14\,000$ & $\sim 60\%$ \\ \hline
Nigéria & $\sim 390$ & $\sim 5\,500$ & $< 10\%$ (pétrole dominant) \\ \hline
Mexique & $\sim 1\,800$ & $\sim 22\,000$ & $\sim 85\%$ \\ \hline
\textbf{BRICS (5 initiaux)} & $\sim 26\,000$ & --- & --- \\ \hline
\end{tabularx}
\end{center}
\textit{Sources : FMI, BM (estimations). PIB nominal en Mds \$ courants ; PPA = Parité de Pouvoir d'Achat.}

\textbf{3. Méthodes express}
\begin{methode}[Localiser et caractériser un foyer émergent sur une carte muette]
\begin{enumerate}
    \item Repérer le continent et la façade maritime (accès mer = commerce).
    \item Placer la capitale politique et la métropole économique (souvent distinctes : ex. Abuja/Lagos, Pretoria/Johannesburg, New Delhi/Mumbai).
    \item Identifier la/les ZES ou ports majeurs (flux entrants/sortants).
    \item Légender : ressources (mines, pétrole, agro), pôles industriels (auto, aéro, textile), flux d'IDE (flèches proportionnelles).
\end{enumerate}
\end{methode}

\begin{methode}[Comparer deux foyers émergents (type Bac : étude de cas ou commentaire)]
\begin{enumerate}
    \item \textbf{Similitudes} : Croissance forte, urbanisation rapide, ouverture commerciale, classe moyenne naissante.
    \item \textbf{Différences structurelles} : Modèle de croissance (exportation bas de gamme vs haute tech vs ressources), rôle de l'État (directeur vs facilitateur), insertion dans les chaînes de valeur mondiales (position amont/aval), vulnérabilités (climatique, dette, inégalités).
\end{enumerate}
\end{methode}

\end{bilan}

\courssub{Checklist avant le Bac}
\begin{aretenir}[Checklist avant le Bac : 10 points de vigilance]
\begin{itemize}
    \item $\square$ Je sais \textbf{localiser} précisément sur une carte mondiale : Brésil, Mexique, Argentine, Chine, Inde, 4 Dragons, EAU, Arabie Saoudite, Afrique du Sud, Nigéria, Égypte, Maroc.
    \item $\square$ Je distingue \textbf{NPI (vagues 1, 2, 3)} et \textbf{Pays émergents actuels} (dont BRICS) : chronologie et critères.
    \item $\square$ Je connais les \textbf{3 moteurs principaux} de l'émergence par cœur : Ressources / Démographie / Réformes structurelles + IDE.
    \item $\square$ Je sais \textbf{caractériser le modèle brésilien} (agrobusiness + aéronautique Embraer + inégalités) vs \textbf{mexicain} (maquiladoras frontalières + industrie auto intégrée USA) vs \textbf{argentin} (agro-export + instabilité macroéco).
    \item $\square$ Je sais expliquer la \textbf{spécificité asiatique} : rôle de l'État développeur, montée en gamme (bas salaires $\to$ haute tech), chaînes de valeur électroniques/automobiles.
    \item $\square$ Je maîtrise les \textbf{4 moteurs africains} : ZA (diversifié, financier), Nigéria (pétrolier, démographique, fintech), Égypte (rente Suez + gaz + démographie), Maroc (hub auto/aéro/phosphate + ZES Tanger Med).
    \item $\square$ Je sais \textbf{lister 4 limites communes} : inégalités (coeff. Gini élevé), extraversion/dépendance (cours matières premières, tech), environnement (pollution, déforestation), gouvernance (corruption, instabilité).
    \item $\square$ Je suis capable de \textbf{réaliser un croquis légendé} d'un foyer émergent (organisation de la légende : Ressources / Infrastructures / Activités / Flux / Villes) en 15 min.
    \item $\square$ Je sais \textbf{analyser un tableau de stats} (PIB, PIB/hab, IDE, Part export manuf, IDH) pour classer des pays et justifier le statut "émergent".
    \item $\square$ Je connais le \textbf{rôle géopolitique des BRICS+} : contester l'ordre occidental, nouvelle banque de développement, commerce en monnaies locales, élargissement 2024.
    \item $\square$ J'évite les \textbf{pièges classiques} : confondre "émergent" et "pays en développement" (PED), croire que l'émergence = développement durable (IDH souvent moyen), oublier l'hétérogénéité interne (côtes vs arrière-pays).
\end{itemize}
\end{aretenir}

\findecours

\end{document}
